\section{基本定义}

    \subsection{Banach 空间}

        \begin{dfn}
            [norm]
            {范数}
            [Norm]
            数域 \(\F\) 上的线性空间 $X$ 上的一个函数 $\norm{\cdot}:X\to\R$ 是 $X$ 上的一个\textbf{范数}, 当且仅当:
            \begin{enumerate}[label=(\roman*)]
                \item \textbf{非退化性 / 正定性}: $\norm{x}\geq 0, \norm{x}=0\iff x=0$;
                \item \textbf{绝对齐次性}: $\norm{\alpha x}=\abs{\alpha}\norm{x}$, $\forall \alpha\in\F, x\in X$;
                \item \textbf{次可加性 / 三角不等式}: $\norm{x+y}\leq\norm{x}+\norm{y}$, $\forall x,y\in X$.
            \end{enumerate}
        \end{dfn}

        \begin{dfn}
            [nvs]
            {赋范线性空间}
            [Normed Vector Space]
            线性空间 \(X\) 是赋范线性空间当且仅当 \(X\) 上存在范数函数.
        \end{dfn}

        \begin{rmk}
            [Dedicatia]
            一般称 \((X, \norm{\cdot})\) 是一个赋范线性空间, 其中 \(X\) 是线性空间, \(\norm{\cdot}\) 是 \(X\) 上的范数函数. 也称 \(X\) 是配备了范数 \(\norm{\cdot}\) 的赋范线性空间.
        \end{rmk}

        \begin{dfn}
            {Cauchy 列}
            [Cauchy Sequence]
            设 \((X, \norm{\cdot})\) 是一个赋范线性空间. 定义 \(X\) 中的一个序列 \(\{x_n\}\) 是一个\textbf{Cauchy 列}, 当且仅当对于任意 \(\varepsilon>0\), 存在 \(N\in\N\) 使得当 \(m,n>N\) 时, \(\norm{x_n-x_m}<\varepsilon\).
        \end{dfn}
        
        \begin{dfn}
            [complete]
            {完备}
            [Completeness]
            设 \((X, \norm{\cdot})\) 是一个赋范线性空间. 定义 \(X\) 是\textbf{完备的}, 当且仅当 \(X\) 中的每个 Cauchy 列都收敛于 \(X\) 中的某个元素.
        \end{dfn}

        \begin{dfn}
            [bs]
            {Banach 空间}
            [Banach Space]
            赋范线性空间 \((X, \norm{\cdot})\) 是 Banach 空间当且仅当 \(X\) 关于由范数诱导的度量 \(d(x,y)=\norm{x-y}\) 是\hyperref[dfn:complete]{完备的}.
        \end{dfn}

        \begin{rmk}
            [Dedicatia]
            所谓\textbf{完备}, 就是 \(X\) 中的每个 Cauchy 列都收敛于 \(X\) 中的某个元素.
        \end{rmk}

        \begin{ins}
            [euclid]
            {配备 \(p\) 范数的有限维 Euclid 空间是 Banach 空间}
            [Euclidean Space with \(p\) Norm is Banach Space]
            设 \(n\in\N\), \(1<p<\infty\), 定义 \(\R^n\) 上的范数 \(\norm{\cdot}_p\) 为
            \[\norm{x}_p=\left(\sum_{i=1}^n\abs{x_i}^p\right)^{\frac{1}{p}}, \forall x=(x_1,\cdots,x_n)\in\R^n\]
            则 \((\R^n, \norm{\cdot}_p)\) 是 Banach 空间.\\
            当 \(p=\infty\) 时, 定义 \(\norm{\cdot}_\infty\) 为
            \[\norm{x}_\infty=\max_{1\leq i\leq n}\abs{x_i}, \forall x=(x_1,\cdots,x_n)\in\R^n\]
            则 \((\R^n, \norm{\cdot}_\infty)\) 也是 Banach 空间.
            \tcbline
            可以验证范数的三个条件:
            \begin{enumerate}[label=(\roman*)]
                \item 正定性: \(\norm{x}_p=0\) 当且仅当 \(x=0\). 
                \item 绝对齐次性: \(\norm{\alpha x}_p=\abs{\alpha}\norm{x}_p\), \(\forall \alpha\in\R, x\in\R^n\). 
                \item 次可加性: 使用 Holder 不等式证明.
            \end{enumerate}
            然后证明完备性即可.
        \end{ins}

        \begin{ins}
            {配备 \(p\) 范数的可和序列空间是 Banach 空间}
            设 \(1\leq p<\infty\), 定义 \(\ell^p\) 为所有满足 \(\sum_{i=1}^\infty\abs{x_i}^p<\infty\) 的序列 \(x=(x_1,x_2,\cdots)\) 的集合. 定义 \(\norm{\cdot}_p\) 为
            \[\norm{x}_p=\left(\sum_{i=1}^\infty\abs{x_i}^p\right)^{\frac{1}{p}}, \forall x=(x_1,x_2,\cdots)\in\ell^p\]
            则 \((\ell^p, \norm{\cdot}_p)\) 是 Banach 空间.\\
            当 \(p=\infty\) 时, 定义 \(\ell^\infty\) 为所有有界序列, 即满足 \(\sup_{i\in\N}\abs{x_i}<\infty\) 的序列 \(x=(x_1,x_2,\cdots)\) 的集合. 定义 \(\norm{\cdot}_\infty\) 为
            \[\norm{x}_\infty=\sup_{i\in\N}\abs{x_i}, \forall x=(x_1,x_2,\cdots)\in\ell^\infty\]
        \end{ins}

        \begin{ins}
            [cont-with-sup-norm]
            {配备上确界范数的连续函数空间是 Banach 空间}
            设 \(X\) 是紧度量空间, 定义 \(\Cont(X)\) 为 \(X\) 上的连续函数空间. 定义 \(\norm{\cdot}_\infty\) 为
            \[\norm{f}_\infty=\sup_{x\in X}\abs{f(x)}, \forall f\in\Cont(X)\]
            则 \((\Cont(X), \norm{\cdot}_\infty)\) 是 Banach 空间.
            \tcbline
            证明的思路如下: 首先 \(\norm{\cdot}_\infty\) 确实是一个范数. 接下来我们只需考虑它是完备的, 也就是证明一致收敛保持连续性.\\
            对于 \(\Cont(X)\) 中的函数列 \(\{f_n\}_{n=1}^\infty\), 考虑极限函数 \(f(x)=\lim_{n\to\infty}f_n(x)\). 因为 \(\{f_n\}\) 是 Cauchy 列, 所以对于任意 \(\varepsilon>0\), 存在 \(N\in\N\) 使得当 \(m,n>N\) 时, \(\norm{f_n-f_m}_\infty<\varepsilon\). 从而对于任意 \(x\in X\), \(\abs{f_n(x)-f_m(x)}<\varepsilon\). 这说明 \(\{f_n(x)\}\) 是一个 Cauchy 列. \\
            因为 \(\R\) 是完备的, 所以 \(\{f_n(x)\}\) 收敛于某个实数 \(f(x)\). 接下来证明 \(f\) 是连续的. 设 \(x_0\in X\), 设 \(\varepsilon>0\). 因为 \(\{f_n\}\) 是 Cauchy 列, 所以存在 \(N\in\N\) 使得当 \(n>N\) 时, \(\norm{f_n-f}_\infty<\varepsilon/3\). 从而对于任意 \(x\in X\), \(\abs{f_n(x)-f(x)}<\varepsilon/3\). 因为 \(f_{N+1}\) 是连续的, 所以存在 \(\delta>0\) 使得当 \(d(x,x_0)<\delta\) 时, \(\abs{f_{N+1}(x)-f_{N+1}(x_0)}<\varepsilon/3\). 从而当 \(d(x,x_0)<\delta\) 时,
            \begin{align*}
            \abs{f(x)-f(x_0)} &\leq \abs{f(x)-f_{N+1}(x)} + \abs{f_{N+1}(x)-f_{N+1}(x_0)} + \abs{f_{N+1}(x_0)-f(x_0)} \\
            &< \varepsilon/3 + \varepsilon/3 + \varepsilon/3 = \varepsilon
            \end{align*}
            从而 \(f\) 是连续的. 
        \end{ins}

        \begin{thm}
            [banach-iff-closed]
            {Banach 空间的子空间是 Banach 空间的充要条件}
            \(X\) 是 Banach 空间, \(Y\) 是 \(X\) 的一个线性子空间. 则 \(Y\) 是 Banach 空间当且仅当 \(Y\) 是 \(X\) 的闭子空间.
        \end{thm}

        \begin{rmk}
            [Dedicatia]
            1. 这里的子空间既是线性子空间, 也是拓扑子空间, 二者是相容的. \\
            2. 线性子空间不一定是闭的. 只有在有限维空间中, 线性子空间才是闭的.
        \end{rmk}

        \begin{prf}
            "(\(\implies\))": 设 \(Y\) 是 Banach 空间, 设 \(y_n\in Y\) 是 \(Y\) 中收敛到 \(y\) 的序列. 则 \(y_n\) 也是 \(X\) 中收敛到 \(y\) 的序列. 因为 \(Y\) 是 Banach 空间, 所以 \(y\in Y\). 从而 \(Y\) 是 \(X\) 的闭子空间. \\
            "(\(\impliedby\))": 设 \(Y\) 是 \(X\) 的闭子空间, 设 \(y_n\in Y\) 是 \(Y\) 中的 Cauchy 列. 因为 \(Y\) 是 \(X\) 的子空间, 所以 \(y_n\) 也是 \(X\) 中的 Cauchy 列. 因为 \(X\) 是 Banach 空间, 所以 \(y_n\) 在 \(X\) 中收敛于某个元素 \(y\). 因为 \(Y\) 是 \(X\) 的闭子空间, 所以 \(y\in Y\). 从而 \(Y\) 是 Banach 空间.
        \end{prf}

        \begin{cxmp}
            {多项式空间在连续函数空间中不是闭的}
            多项式显然都是连续函数, 且多项式是 \(C_0\) 中的子空间, 但一列多项式的极限可能不是多项式, 如:
            \[p_n(x)=\sum_{k=0}^n\frac{x^k}{k!}\]
            其一致收敛到 \(\exp x\), 但 \(\exp x\) 不是多项式. 从而多项式空间在 \(C_0\) 中不是闭的.
        \end{cxmp}

        \continued

    \subsection{线性算子}
    
        \begin{dfn}
            [blo]
            {有界线性算子}
            [Bounded Linear Operator]
            设 \((X, \norm{\cdot}_X)\) 和 \((Y, \norm{\cdot}_Y)\) 是两个赋范线性空间. 定义线性映射 \(T:X\to Y\) 是一个\textbf{有界线性算子}, 当且仅当存在 \(M\geq 0\) 使得
            \[\forall x\in X, \norm{Tx}_Y\leq M\norm{x}_X\]
        \end{dfn}

        \begin{dfn}
            [blf]
            {有界线性泛函}
            [Bounded Linear Functional]
            设 \((X, \norm{\cdot})\) 是一个赋范线性空间. 定义线性映射 \(f:X\to\R\) 是一个\textbf{有界线性泛函}, 当且 \(f\) 是一个有界线性算子.
        \end{dfn}

        \begin{rmk}
            [Dedicatia]
            直观地理解, 有界线性算子将 \(Y\) 的单位闭球拉回 \(X\) 中半径不超过 \(M\) 的单位闭球中.
        \end{rmk}

        \begin{dfn}
            {有界线性算子的范数}
            [Norm of Bounded Linear Operator]
            设 \((X, \norm{\cdot}_X)\) 和 \((Y, \norm{\cdot}_Y)\) 是两个赋范线性空间, \(T:X\to Y\) 是一个有界线性算子. 定义 \(T\) 的范数为
            \[\norm{T}:=\inf\{M\geq 0 : \forall x\in X, \norm{Tx}_Y\leq M\norm{x}_X\}\]
        \end{dfn}

        \begin{rmk}
            [Dedicatia]
            \(T\) 的范数 \(\norm{T}\) 是满足 \(\forall x\in X, \norm{Tx}_Y\leq M\norm{x}_X\) 的最小的 \(M\).
        \end{rmk}

        \begin{ppt}
            {有界线性算子范数的等价定义}
            [Equivalent Definition of Norm of Bounded Linear Operator]
            设 \((X, \norm{\cdot}_X)\) 和 \((Y, \norm{\cdot}_Y)\) 是两个赋范线性空间, \(T:X\to Y\) 是一个有界线性算子. 则 \(T\) 的范数为
            \[\norm{T}=\sup_{x\neq 0}\frac{\norm{Tx}_Y}{\norm{x}_X}=\sup_{\norm{x}_X=1}\norm{Tx}_Y\]
        \end{ppt}

        \begin{rmk}
            [Dedicatia]
            为什么可以这样等价, 是因为 \(X,Y\) 都是线性空间, 可以进行标量乘法. \\
            所以对于任意 \(x\in X\), 当 \(x\neq 0\) 时, 可以将 \(x\) 写成 \(\norm{x}_X\cdot\frac{x}{\norm{x}_X}\). 从而 \(\norm{Tx}_Y=\norm{T(\norm{x}_X\cdot\frac{x}{\norm{x}_X})}_Y=\norm{x}_X\norm{T(\frac{x}{\norm{x}_X})}_Y\). 因此 \(\frac{\norm{Tx}_Y}{\norm{x}_X}=\norm{T(\frac{x}{\norm{x}_X})}_Y\). 当 \(\norm{x}_X=1\) 时, \(\frac{\norm{Tx}_Y}{\norm{x}_X}=\norm{Tx}_Y\).
        \end{rmk}

        \begin{ppt}
            {有界线性算子范数的性质}
            [Properties of Norm of Bounded Linear Operator]
            设 \(T\) 是有界线性算子, 则 \(\norm{Tx}\leq\norm{T}\norm{x}\).
        \end{ppt}

        \begin{ppt}
            {有界线性算子形成赋范线性空间}
            [Bounded Linear Operators Form a Normed Vector Space]
            设 \((X, \norm{\cdot}_X)\) 和 \((Y, \norm{\cdot}_Y)\) 是两个赋范线性空间. 定义 \(\blo(X,Y)\) 是从 \(X\) 到 \(Y\) 的所有有界线性算子组成的集合. 则 \(\blo(X,Y)\) 是赋范线性空间.
        \end{ppt}

        \begin{thm}
            [blo-banach-iff-codom-banach]
            {有界线性算子空间形成 Banach 空间的充要条件}
            设 \(X,Y\) 是两个赋范线性空间. 则 \(\blo(X,Y)\) 是 Banach 空间当且仅当 \(Y\) 是 Banach 空间.
        \end{thm}

        \begin{prf}
            设 \(\{T_n\}\) 是一列 \(X\) 到 \(Y\) 的有界线性算子, 使得 \(\{T_n\}\) 是 \(\blo(X,Y)\) 中的 Cauchy 列. 则对于任意 \(x\in X\), \(\{T_n x\}\) 是 \(Y\) 中的 Cauchy 列. 因为 \(Y\) 是完备的（Banach 空间）, 所以存在 \(y\in Y\) 使得 \(T_n x\to y\). 定义 \(T:X\to Y\) 使得 \(Tx=y\). \\\
            验证: (1) \(T\) 是有界线性算子: 线性性是显然的; 有界性是由于
            \[\norm{Tx}\leq\norm{T_N x}+\lim_{n\to\infty}\norm{T_n x-T_N x}\leq\norm{T_N x}+\epsilon\norm{x}\leq(\norm{T_N}+\epsilon)\norm{x}\]
            (2) \(T_n\) 依范数收敛到 \(T\): 
            \[\norm{T_n-T}=\sup_{\norm{x}=1}\norm{T_n x-Tx}=\sup_{\norm{x}=1}\lim_{k\to\infty}\norm{T_n x-T_k x}=0.\]
        \end{prf}

        \begin{crl}
            {依范数收敛蕴含范数收敛}
            若 \(T_n\to T\), 则 \(\norm{T_n}\to\norm{T}\) 作为实数序列收敛.
        \end{crl}

        \begin{crl}
            {依范数收敛蕴含逐点收敛}
            若 \(T_n\to T\), 则 \(\forall\,x\in X, T_n x\to Tx\).
        \end{crl}

        \begin{thm}
            [bounded-iff-continuous]
            {有界等价于连续}
            [Bounded iff Continuous]
            设 \(X,Y\) 是两个赋范线性空间, \(T:X\to Y\) 是一个线性映射. 则 \(T\) 是有界的当且仅当 \(T\) 是连续的.
        \end{thm}

        \begin{prf}
            仅作提示:\\
            "(\(\implies\))": 设 \(T\) 是有界的, 则存在 \(M\geq 0\) 使得 \(\forall x\in X, \norm{Tx}_Y\leq M\norm{x}_X\). 则对于任意 \(\epsilon>0\), 当 \(\norm{x-x_0}_X<\frac{\epsilon}{M}\) 时, \(\norm{Tx-Tx_0}_Y=\norm{T(x-x_0)}_Y\leq M\norm{x-x_0}_X<\epsilon\). 从而 \(T\) 在 \(x_0\) 处连续. 因为 \(x_0\) 是任意的, 所以 \(T\) 是连续的.\\
            "(\(\impliedby\))": 设 \(T\) 是连续的, 则 \(T\) 在 \(0\) 处连续. 则存在 \(\delta>0\) 使得当 \(\norm{x}_X<\delta\) 时, \(\norm{Tx}_Y<1\). 根据赋范线性空间的平移不变性, 对于任意 \(x\in X\), 当 \(\norm{x}_X<\delta\) 时, \(\norm{Tx}_Y<1\). 从而对于任意 \(x\in X\), \(\norm{Tx}_Y<\frac{1}{\delta}\norm{x}_X\). 因此 \(T\) 是有界的.
        \end{prf}

        \begin{rmk}
            [Dedicatia]
            在 Banach 空间中, 算子的有界性与连续性等价, 这是因为线性空间有数乘结构, 可以将单位球缩放到任意小.\\
            事实上有界线性算子不仅是连续的, 而且是 Lipschitz 连续的. 也就是说
            \[\norm{Tx-Ty}_Y\leq\norm{T}\norm{x-y}_X\]
        \end{rmk}

        \begin{thm}
            {有界线性算子的复合是有界线性算子}
            [Composition of Bounded Linear Operators is Bounded]
            设 \(X,Y,Z\) 是三个赋范线性空间, \(T:X\to Y\) 和 \(S:Y\to Z\) 是有界线性算子. 则 \(S\circ T: X\to Z\) 也是有界线性算子. 且满足:
            \[\norm{S\circ T}\leq\norm{S}\norm{T}\]
        \end{thm}

        \begin{prf}
            令 \(\norm{x}=1\), 则
            \[\norm{S\circ T}=\sup\norm{S(T x)}=\sup\norm{S(y)}\leq\norm{S}\sup\norm{T x}=\norm{S}\norm{T}.\]
        \end{prf}

        \begin{ins}
            {全矩阵空间是 Banach 空间}
            全体矩阵 \(M_{m\times n}(\R)\cong\blo(\R^n, \R^m)\) 是 Banach 空间, 其范数定义为
            \[\norm{A}=\sup_{\norm{x}=1}\norm{Ax}\]
            且满足 \(\dim M_{m\times n}(\R)=mn\).
        \end{ins}

        \begin{ins}
            {积分算子是有界线性算子}
            定义算子 \(T:\Cont([0,1])\to\Cont([0,1])\) 使得
            \[Tf(x)=\int_0^x f(t)dt\]
            则 \(T\) 是一个有界线性算子, 且 \(\norm{T}=1\).
        \end{ins}

        \begin{cxmp}
            {微分算子不是有界线性算子}
            定义算子 \(D:\Cont^1([0,1])\to\Cont([0,1])\) 使得
            \[Df(x)=f'(x)\]
            则 \(D\) 不是一个有界线性算子. 这是因为考虑函数列 \(f_n=x^n\), 则 \(\norm{f_n}_\infty=1\), 但 \(\norm{Df_n}_\infty=n\to\infty\).
            \tcbline
            这好像很违背直觉. 但是从另一个角度看, 微分算子也不是连续的.
        \end{cxmp}

        \begin{vardfn}
            {Banach 代数}
            [Banach Algebra]
            设 \(X\) 是一个 Banach 空间, 且 \(X\) 上定义了一个二元运算 \(\cdot\), 满足:
            \begin{enumerate}[label=(\roman*)]
                \item \((X, \cdot)\) 是代数, 满足代数运算公理;
                \item 有界性: \(\exists M\geq 0\) 使得 \(\forall x,y\in X, \norm{x\cdot y}\leq M\norm{x}\norm{y}\).
            \end{enumerate}
            则 \(X\) 是一个\textbf{Banach 代数}.
        \end{vardfn}

    \subsection{有限维 Banach 空间}

        \begin{intrormk}
            [Dedicatia]
            在有限维的情形下 Banach 空间的结构非常规整, 也有很多漂亮的性质.
        \end{intrormk}

        \begin{dfn}
            [dual-space]
            {对偶空间}
            [Dual Space]
            设 \((X, \norm{\cdot})\) 是一个赋范线性空间. 定义 \(X\) 的\textbf{对偶空间} \(X^*\) 是 \(\blo(X, \R)\).
        \end{dfn}

        \begin{rmk}
            [Dedicatia]
            \(X^*\) 是 \(X\) 上的所有有界线性泛函组成的集合. 对偶空间中的元素都是线性泛函. \\
            需要注意的是, 这和作为线性空间的对偶空间的定义是不同的. 作为线性空间的对偶空间是所有线性泛函组成的集合, 而作为赋范线性空间的对偶空间是所有\textbf{有界线性泛函}组成的集合.\\
            只有在有限维线性空间中这两种定义才是相同的. 因为有限维的线性泛函一定是有界的.
        \end{rmk}

        \begin{thm}
            {有限维赋范线性空间的对偶空间同构于原空间}
            设 \(V\) 是赋范线性空间, \(\dim V<\infty\). 则 \(V^*\cong V\).
        \end{thm}

        \begin{prf}
            提示: 注意到二者的维数相同, 取定一组 \(V\) 的基 \(\{v_i\}_i^n\). 定义映射
            \[T:\qquad\begin{matrix}
                V & \to & V^* \\
                v_i & \mapsto & f_i
            \end{matrix}\]
            其中 \(f_i\) 是 \(V\) 上的线性泛函, 满足 \(f_i(v_j)=\delta_{ij}\). \\
            或者直接使用下面的同构定理.
        \end{prf}

        \begin{thm}
            [riesz]
            {Riesz 引理}
            [Riesz's Lemma]
            设 \(X\) 是一个赋范线性空间, \(Y\) 是 \(X\) 的一个闭子空间, 且 \(Y\neq X\). 则对于任意 \(0<\epsilon<1\), 存在 \(x_0\in X\) 使得 \(\norm{x_0}=1\) 且 
            \[\inf_{x\in Y}\norm{x_0-x}\geq \epsilon\]
        \end{thm}

        \begin{prf}
            任取 \(a\in X\setminus Y\), 则根据 \(Y\) 的闭性, \(d:=\inf_{y\in Y}\norm{a-y}>0\). 从 \(Y\) 中取 \(y\) s.t. \(\norm{a-y}<\frac{d}{\epsilon}\), 规范化:
            \[x_0=\frac{a-y}{\norm{a-y}}\]
            那么 \(\norm{x_0}=1\), 且对于任意 \(x\in Y\),
            \[\norm{x_0-x}=\norm{\frac{a-y}{\norm{a-y}}-x}=\frac{1}{\norm{a-y}}\norm{a-(y+\norm{a-y}x)}\geq\frac{d}{\norm{a-y}}>\epsilon.\]
        \end{prf}

        \begin{dfn}
            [heine-borel]
            {Heine-Borel 性质}
            [Heine-Borel Property]
            设 \(X\) 是一个赋范线性空间. 定义 \(X\) 满足\textbf{Heine-Borel 性质}, 当且仅当 \(X\) 中的每个有界闭集都是\hyperref[dfn:compact]{紧致集}.
        \end{dfn}

        \begin{thm}
            [riesz-2]
            {无限维赋范线性空间的单位闭球不是紧的}
            设 \(X\) 是一个赋范线性空间, \(\dim X=\infty\). 则 \(X\) 中的单位闭球 \(B=\{x\in X:\norm{x}\leq 1\}\) 不是紧致集.
        \end{thm}

        \begin{rmk}
            [Dedicatia]
            由于赋范线性空间一定是度量空间, 而在度量空间中, \hyperref[dfn:compact]{紧致性}与\hyperref[dfn:sequential-compact]{列紧性}是等价的, 我们将不再区分这二者, 统称为紧性. 实际上, 下面的证法就是在证明单位闭球不是列紧的.
        \end{rmk}

        \begin{prf}
            任取 \(x_1\in X\), \(\norm{x_1}=1\), \(x_1\) 张成了子空间 \(X_1:=\{\alpha x_1: \alpha \in \R\}\), 根据 \hyperref[thm:riesz]{Riesz 引理}, 存在 \(x_2\in X\) 使得 \(\norm{x_2}=1\) 且 \(\inf_{x\in X_1}\norm{x_2-x}\geq\frac{1}{2}\). \\
            继续: \(x_1, x_2\) 张成了子空间 \(X_2:=\{\alpha x_1+\beta x_2: \alpha, \beta \in \R\}\), 根据 \hyperref[thm:riesz]{Riesz 引理}, 存在 \(x_3\in X\) 使得 \(\norm{x_3}=1\) 且 \(\inf_{x\in X_2}\norm{x_3-x}\geq\frac{1}{2}\). \\
            以此类推, 可以得到一列 \(\{x_n\}\) 使得 \(\norm{x_n}=1\) 且 \(\inf_{x\in X_{n-1}}\norm{x_n-x}\geq\frac{1}{2}\), 其中 \(X_{n-1}=\{\sum_{i=1}^{n-1}\alpha_i x_i: \alpha_i \in \R\}\). 则对于任意 \(m<n\), \(\norm{x_n-x_m}=\inf_{x\in X_{n-1}}\norm{x_n-x}\geq\frac{1}{2}\). 从而 \(\{x_n\}\) 中不存在收敛子列. 因此单位闭球不是紧的.
        \end{prf}

        \begin{cxmp}
            {连续函数空间的单位闭球不是紧的}
            定义 \(\Cont([0,1])\) 上的函数列为
            \[f_n(x) = x^n\]
            则可以验证 \(\norm{f_n}_\infty=1\), 这是一个有界函数列, 它确实在单位闭球中, 但 \(f_n\) 的极限函数是
            \[f(x) = \begin{cases}
                0, & x\in[0,1) \\
                1, & x=1
            \end{cases}\]
            这说明 \(\{f_n\}\) 中不存在收敛子列. 从而 \(\Cont([0,1])\) 的单位闭球不是\hyperref[dfn:sequential-compact]{序列紧的}, 从而不是紧的.\\
            这个函数列不是 Cauchy 列, 也不具有一致收敛极限. 因此它不能用来说明 \(\Cont([0,1])\) 不是 Banach 空间.
        \end{cxmp}

        \begin{thm}
            {有限维赋范线性空间同构定理}
            设 \(V\) 是赋范线性空间, \(n=\dim V<\infty\). 则 \(V\cong\R^n\).
        \end{thm}

        \begin{rmk}
            [Dedicatia]
            这里我们说的同构指的是作为赋范线性空间的同构, 也就是说存在一个双射 \(T:V\to\R^n\) 使得 \(T\) 和 \(T^{-1}\) 都是有界线性算子. 
        \end{rmk}

        \begin{prf}
            取定一组 \(V\) 的基 \(\{v_i\}_i^n\). 定义映射
            \[T:\qquad\begin{matrix}
                \R^n & \to & V \\
                (x_1, x_2, \cdots, x_n) & \mapsto & \sum_{i=1}^n x_i v_i
            \end{matrix}\]
            验证:\\
            (1) 线性性显然成立;\\
            (2) \(T\) 是双射, 根据有限维线性空间的同构, 这也是显然的;\\
            (3) 考虑 \(T\) 的有界性:
            \[\norm{T(x_1, x_2, \cdots, x_n)}=\norm{\sum_{i=1}^n x_i v_i}\leq\sum_{i=1}^n|x_i|\norm{v_i}\leq\qty(\sum_{i=1}^n\abs{x_i})\qty(\sum_{i=1}^n\norm{v_i})  \]
            这说明
            \[\norm{T}\leq\sum_{i=1}^n\norm{v_i}<\infty\]
            (4) 考虑 \(T^{-1}\) 的有界性: 反证: 假如不是有界的, 那么不存在这样的 \(M\). 也就是说存在一列 \(\{x_n\}\) 使得
            \[\norm{T^{-1}(x_n)}>n\norm{x_n}\]
            根据齐次性, 则可以取定规范的 \(\{x_n\}\) s.t. \(\norm{x_n}=1\). 则 \(\norm{T^{-1}(x_n)}>n\). 令
            \[a_n = \frac{T^{-1}(x_n)}{\norm{T^{-1}(x_n)}}\]
            那么 \(\norm{a_n}=1\) 对所有 \(n\) 成立. 由于 \(V\) 是有限维的, 所以 \(V\) 中的单位球是紧的. 从而 \(\{a_n\}\) 中存在一个收敛子列 \(\{a_{n_k}\}\), 设其极限为 \(a\). 那么
            \[\norm{T^{-1}a} = \lim_{k \to \infty} \norm{T^{-1}a_{n_k}} = \lim_{k \to \infty} \frac{T^{-1}(x_{n_k})}{\norm{T^{-1}x_{n_k}}}\]
            由于 \(\norm{T^{-1}x_{n_k}}>n_k\), 所以 \(\norm{T^{-1}a}=0\). 从而 \(a=0\). 这与 \(\norm{a}=1\) 矛盾. 因此 \(T^{-1}\) 是有界的.
        \end{prf}

        \begin{thm}
            {有限维空间的范数等价}
            [Norms on Finite-Dimensional Spaces are Equivalent]
            设 \(V\) 是一个有限维的线性空间, \(\norm{\cdot}_1\) 和 \(\norm{\cdot}_2\) 是 \(V\) 上的两个范数. 则存在常数 \(C_1, C_2>0\) 使得
            \[\forall x\in V, C_1\norm{x}_1\leq\norm{x}_2\leq C_2\norm{x}_1\]
            或者等价地写作: 存在常数 \(C>0\) 使得
            \[\forall x\in V, \frac{1}{C}\norm{x}_1\leq\norm{x}_2\leq C\norm{x}_1.\]
        \end{thm}

        \begin{rmk}
            [Dedicatia]
            我们已经证明了决定有限维赋范线性空间的结构的唯一因素就是维数. 因此只要维度一样, 不论是什么范数一定可以建立同构. 自然也就是等价的.
        \end{rmk}

        \begin{dfn}
            [bidual-space]
            {双对偶空间}
            [Bidual Space]
            定义 \(X\) 的\textbf{双对偶 / 二次对偶空间} \(X^{**}\) 是 \(X^*\) 的对偶空间, 即 \(X^{**}=(X^*)^*=\blo(X^*, \R)\).
        \end{dfn}

        \begin{dfn}
            [canonical-embedding]
            {典范嵌入}
            [Canonical Embedding]
            定义 \(X\) 的\textbf{典范嵌入} \(J_X:X\to X^{**}\) 是一个线性映射:
            \[J_X(f)(x) = f(x).\]
            根据定义, 这是一个单射, 记作 \(J_X:X\hookrightarrow X^{**}\).
        \end{dfn}

        \begin{dfn}
            [reflexive]
            {自反空间}
            [Reflexive Space]
            设 \(X\) 是一个赋范线性空间. 定义 \(X\) 是\textbf{自反的}, 当且仅当 \(X\) 的\hyperref[dfn:canonical-embedding]{典范嵌入}是双射.
        \end{dfn}

        \begin{rmk}
            注意: 这里的\hyperref[dfn:reflexive]{自反空间}是利用\hyperref[dfn:canonical-embedding]{典范嵌入}定义的. 有可能出现如下的情况:
            \begin{itemize}
                \item \(X\cong X^{**}\), 但 \(J_X\) 不是双射; 这时按定义 \(X\) 不是自反空间;
            \end{itemize}
            即使 \(X\) 同构于 \(X^{**}\) 也不能说明它们是通过 \(J_X\) 建立的同构.
        \end{rmk}

        \begin{ppt}
            {有限维赋范线性空间是自反空间}
        \end{ppt}

        \begin{prf}
            根据维数相同, 显然.
        \end{prf}

        \begin{dfn}
            [separable]
            {可分 / 可析}
            [Separable]
            设 \(X\) 是一个赋范线性空间. 定义 \(X\) 是\textbf{可分的}, 当且仅当 \(X\) 中存在一个可数的稠密子集.
        \end{dfn}

        \begin{rmk}
            [Dedicatia]
            可分空间的含义是可以用空间中的一列元素来近似空间中的任意元素. \\
            ``Separable'' 翻译为``可分''. 在《实变函数论与泛函分析》中翻译为了``可析''.
        \end{rmk}

        \begin{varthm}
            {可分的赋范线性空间的原对偶空间也是可分的}
            设 \(X\) 是赋范线性空间, 则如果 \(X^*\) 是\hyperref[dfn:separable]{可分}的, 则 \(X\) 也是可分的.
        \end{varthm}

        \begin{prf}
            由于 \(X^*\) 是可分的, 存在一列 \(\{f_n\}\) 在单位球中稠密. 对每个 \(f_n\), 由于 \(\sup_{\norm{x}=1}|f_n(x)|=\norm{f_n}>\frac{1}{2}\), 则必然存在 \(x_n\) 使得 
            \[|f_n(x_n)|>\frac{1}{2}\]
            这一列 \(\{x_n\}\) 张成线性闭子空间 \(X_0\). 如果 \(X\) 不是可分的, 那么 \(X_0\neq X\), 从而在 \(X^*\) 中存在一个元素 \(f\) 使得 \(f(x)\to 0\) 但 
            \[\norm{f_n-f}\geq|f_n(x_n)-f(x_n)|=|f_n(x_n)|>\frac{1}{2}\]
            这与 \(\{f_n\}\) 在单位球中稠密矛盾. 因此 \(X\) 是可分的.
        \end{prf}

        \begin{ins}
            {序列空间 \(\ell^\infty\) 是不可分的}
            全体有界序列构成的空间 \(\ell^\infty\) 配备上确界范数不是可分的. 因为所有由 \(0,1\) 构成的序列是不可数的, 且任意两个不同的由 \(0,1\) 构成的序列之间的距离都是 \(1\). 从而 \(\ell^\infty\) 中存在一个不可数的集合, 使得集合中的任意两个元素之间的距离都大于等于 \(1\). 按定义 \(\ell^\infty\) 不是可分的.
        \end{ins}

    \subsection{Hilbert 空间}

        \begin{intrormk}
            [Dedicatia]
            Hilbert 空间被称为``最像 Euclid 空间的 Banach 空间'', 因为它具有内积结构, 可以定义角度和正交, 基本符合了我们生活经验中的直觉. 这使得 Hilbert 空间在数学和物理中都有广泛的应用.\\
            本节我们就来介绍一下 Hilbert 空间的基本概念和性质.
        \end{intrormk}

        \begin{dfn}
            [inner-product]
            {内积}
            [Inner Product]
            设 \(X\) 是 \(\F\) 上的线性空间, 定义二元运算 \(\langle\cdot,\cdot\rangle:X\times X\to\F\) 是 \(X\) 上的一个\textbf{内积 (Inner Product)}, 当且仅当满足:
            \begin{enumerate}[label=(\roman*)]
                \item \textbf{正定性}: \(\langle x,x\rangle\geq 0\), \(\forall x\in X\),  \(x=0\iff \langle x,x\rangle=0\);
                \item \textbf{共轭对称性}: \(\langle x,y\rangle=\overline{\langle y,x\rangle}\), \(\forall x,y\in X\);
                \item \textbf{线性性}: \(\langle \alpha x+y,z\rangle=\alpha\langle x,z\rangle+\langle y,z\rangle\), \(\forall \alpha\in\F, x,y,z\in X\).
            \end{enumerate}
        \end{dfn}

        \begin{dfn}
            [ips]
            {内积空间}
            [Inner Product Space]
            设 \(X\) 是 \(\F\) 上的线性空间, \(\langle\cdot,\cdot\rangle\) 是 \(X\) 上的一个内积. 则 \((X, \langle\cdot,\cdot\rangle)\) 是一个\textbf{内积空间}.
        \end{dfn}

        \begin{ppt}
            {内积可以诱导范数}
            [Inner Product Induces Norm]
            设 \((X, \langle\cdot,\cdot\rangle)\) 是一个内积空间. 定义
            \[\norm{x}=\sqrt{\langle x,x\rangle}\]
            则 \(\norm{\cdot}\) 是 \(X\) 上的一个范数.
        \end{ppt}

        \begin{dfn}
            [hs]
            {Hilbert 空间}
            [Hilbert Space]
            设 \((X, \langle\cdot,\cdot\rangle)\) 是一个内积空间, \(\norm{\cdot}\) 是由内积诱导的范数. 则如果 \((X, \norm{\cdot})\) 是一个 Banach 空间, 则称 \(X\) 是一个\textbf{Hilbert 空间}.
        \end{dfn}

        \begin{rmk}
            [Dedicatia]
            完备的内积空间就是 Hilbert 空间.
        \end{rmk}

        \begin{thm}
            [cauchy-schwarz]
            {Cauchy-Schwarz 不等式}
            [Cauchy-Schwarz Inequality]
            设 \((X, \langle\cdot,\cdot\rangle)\) 是一个内积空间, \(\norm{\cdot}\) 是由内积诱导的范数. 则对于任意 \(x,y\in X\), 都有
            \[\abs{\langle x,y\rangle}\leq\norm{x}\norm{y}.\]
        \end{thm}

        \begin{prf}
            构造二次函数 
            \[f(\alpha)=\norm{\alpha x+y}^2=\langle \alpha x+y, \alpha x+y\rangle=\abs{\alpha}^2\norm{x}^2+2\Re(\alpha\langle x,y\rangle)+\norm{y}^2.\]
            考虑其最小值.
        \end{prf}

        \begin{thm}
            [parallelogram-law]
            {范数可以诱导内积的条件 / 平行四边形法则}
            [Condition for Norm to Induce Inner Product / Parallelogram Law]
            设 \((X, \norm{\cdot})\) 是一个赋范线性空间. 则存在一个内积使得 \(\norm{\cdot}\) 是由该内积诱导的范数当且仅当
            \[\forall x,y\in X, \quad \norm{x+y}^2+\norm{x-y}^2=2\norm{x}^2+2\norm{y}^2.\]
        \end{thm}

        \begin{rmk}
            [Dedicatia]
            平行四边形法则是很强的限制. 这说明了多数范数都不来源于内积.
        \end{rmk}

        \begin{thm}
            {赋范线性空间的内积的显式表示 / 极化恒等式}
            [Explicit Representation of Inner Product / Polarization Identity]
            设 \((X, \norm{\cdot})\) 是一个赋范线性空间, \(\norm{\cdot}\) 满足平行四边形法则. 则对于任意 \(x,y\in X\), 内积可以通过范数来表示:
            \[\langle x,y\rangle=\frac{1}{4}\qty(\norm{x+y}^2-\norm{x-y}^2+ \ii\norm{x+\ii y}^2 - \ii\norm{x-\ii y}^2).\]
            如果 \(X\) 是实数域上的内积空间, 则内积的表示式可以简化为
            \[\langle x,y\rangle=\frac{1}{4}\qty(\norm{x+y}^2-\norm{x-y}^2).\]
        \end{thm}

        \begin{dfn}
            [orthogonality]
            {正交}
            [Orthogonality]
            设 \((X, \langle\cdot,\cdot\rangle)\) 是一个内积空间. 定义 \(x,y\in X\) 是\textbf{正交的}, 当且仅当 \(\langle x,y\rangle=0\). 记作 \(x\perp y\).
        \end{dfn}

        \begin{dfn}
            [orthogonal-complement]
            {正交补}
            [Orthogonal Complement]
            设 \((X, \langle\cdot,\cdot\rangle)\) 是一个内积空间, \(Y\) 是 \(X\) 的一个子集. 定义 \(Y\) 的\textbf{正交补} \(Y^\perp\) 是
            \[Y^\perp=\{x\in X: \langle x,y\rangle=0, \forall y\in Y\}.\]
        \end{dfn}

        \begin{ppt}
            {正交补是闭子空间}
            [Orthogonal Complement is a Closed Subspace]
            设 \((X, \langle\cdot,\cdot\rangle)\) 是一个内积空间, \(Y\) 是 \(X\) 的一个子集. 则 \(Y^\perp\) 是 \(X\) 的一个闭子空间.
        \end{ppt}

        \begin{dfn}
            [orthogonal-projection]
            {正交投影}
            [Orthogonal Projection]
            设 \((X, \langle\cdot,\cdot\rangle)\) 是一个内积空间, \(Y\) 是 \(X\) 的一个子空间. 定义 \(P_Y:X\to Y\) 是 \(X\) 到 \(Y\) 的一个映射, 满足: 
            \[\forall x\in X, P_Y x=\inf_{y\in Y} \norm{x-y}.\]
            且满足
            \[\forall x\in X, x-P_Y x\in Y^\perp.\]
            则称 \(P_Y\) 是从 \(X\) 到 \(Y\) 的一个\textbf{正交投影}.
        \end{dfn}

        \begin{thm}
            {正交投影的存在唯一性}
            [Existence and Uniqueness of Orthogonal Projection]
            设 \((X, \langle\cdot,\cdot\rangle)\) 是一个内积空间, \(Y\) 是 \(X\) 的一个完备的线性子空间. 则 \(\forall\,x\in X\), \(x\) 在 \(Y\) 中有唯一的\hyperref[dfn:orthogonal-projection]{正交投影}.
        \end{thm}

        \begin{prf}
            取一列 \(x_n\in Y\) 使得
            \[\lim_{n\to\infty}\norm{x-x_n}=\inf_{y\in Y} \norm{x-y}=:d.\]
            根据内积与范数的关系, 我们有
            \[\norm{x_m-x_n}^2=2(\norm{x-x_m}^2+\norm{x-x_n}^2)-4\norm{\frac{x_m+x_n}{2}-x}^2\leq 2(\norm{x-x_m}^2+\norm{x-x_n}^2)-4d^2 \]
            当 \(m,n\to 0\) 时, \(\norm{x_m-x_n}\to 0\). 从而 \(\{x_n\}\) 是 \(Y\) 中的一个 Cauchy 列. 根据 \(Y\) 的完备性, \(\{x_n\}\) 收敛于 \(Y\) 中的唯一的元素 \(y_0\). 那么
            \[\norm{x-y_0}=\inf_{y\in Y} \norm{x-y}=d.\]
            接下来证明正交性: 任意取 \(z\in Y\), 有
            \[d^2\leq\norm{x-y_0-\lambda z}^2=\norm{x-y_0}^2-2\lambda\Re\langle x-y_0, z\rangle+|\lambda|^2\norm{z}^2.\]
            取
            \[\lambda=\frac{\langle x-y_0, z\rangle}{\norm{z}^2}\]
            则
            \[d^2\leq d^2-\frac{|\langle x-y_0, z\rangle|^2}{\norm{z}^2}\]
            因此只能有 \(\langle x-y_0, z\rangle=0\). 从而 \(x-y_0\in Y^\perp\).
        \end{prf}

        \begin{ppt}
            {子集的正交补可以扩展到子空间上}
            \(X\) 是内积空间, \(S\subseteq X\). $S^\perp=\overline{\spn(S)}^\perp$.
        \end{ppt}

        \begin{ppt}
            [orthogonal-direct-sum]
            {正交直和}
            [Orthogonal Direct Sum]
            设 \((X, \langle\cdot,\cdot\rangle)\) 是一个内积空间, \(Y\) 是 \(X\) 的一个子空间. 则 \(X=Y\oplus Y^\perp\).
        \end{ppt}

        \begin{ppt}
            [pythagorean-equality]
            {勾股定理}
            [Pythagorean Equality]
            设 \((X, \langle\cdot,\cdot\rangle)\) 是一个内积空间, \(x,y\in X\) 满足 \(x\perp y\). 则
            \[\norm{x+y}^2=\norm{x}^2+\norm{y}^2.\]
        \end{ppt}

        \begin{dfn}
            [orthonormal-set]
            {标准正交组 / 标准正交集}
            [Orthonormal Set]
            设 \((X, \langle\cdot,\cdot\rangle)\) 是一个内积空间. 定义 \(S\subseteq X\) 是一个\textbf{标准正交组}, 当且仅当对于任意 \(x,y\in S\), 都有
            \[\langle x,y\rangle=\delta_{xy}=\begin{cases}
                1, & x=y; \\
                0, & x\neq y.
            \end{cases}\]
        \end{dfn}

        \begin{dfn}
            [onb]
            {标准正交基}
            [Orthonormal Basis]
            设 \((X, \langle\cdot,\cdot\rangle)\) 是一个内积空间. 定义 \(S\subseteq X\) 是一个\textbf{标准正交基}, 当且仅当 \(S\) 是 \(X\) 的一个标准正交组, 且 \(\overline{\spn(S)}=X\).
        \end{dfn}

        \begin{rmk}
            [Dedicatia]
            也就是说标准正交基张成的线性空间在 \(X\) 中稠密, 也就是说可以用标准正交基的线性组合来逼近 \(X\) 中的任意元素.\\
            按照我们这里的定义, 标准正交基满足完备性. 有些参考材料将其定义为满足 \(\spn(S)=X\) 的标准正交组. 需要注意细致区分. 
        \end{rmk}

        \begin{rmk}
            [Dedicatia]
            如果 \(\dim X=\infty\), 那么 \(X\) 的 \hyperref[dfn:hamel-basis]{Hamel 基}一定是不可数的, 但是 \(X\) 的标准正交基可以是可数的. 这说明了 Hamel 基与 Hilbert 空间的标准正交基的区别.\\
            事实上, 标准正交基是一种 Schauder 基, 与 Hamel 基完全不同. 
        \end{rmk}

        \begin{thm}
            [bessel]
            {Bessel 不等式}
            [Bessel's Inequality]
            设 \((X, \langle\cdot,\cdot\rangle)\) 是一个内积空间, \(S\subseteq X\) 是一个标准正交组. 则对于任意 \(x\in X\), 都有
            \[\sum_{s\in S} \abs{\langle x,s\rangle}^2\leq\norm{x}^2.\]
        \end{thm}

        \begin{thm}
            [parseval]
            {Parseval 等式}
            [Parseval's Identity]
            设 \((X, \langle\cdot,\cdot\rangle)\) 是一个内积空间, \(S\subseteq X\) 是一个（完备的）标准正交基. 则对于任意 \(x\in X\), 都有
            \[\sum_{s\in S} \abs{\langle x,s\rangle}^2=\norm{x}^2.\]
        \end{thm}

        \begin{ins}
            [fourier-series]
            {Fourier 级数}
            [Fourier Series]
            设 \(L^2(\mathbb{T}[0,2\pi])\) 是定义在 \([0,2\pi]\) 上的满足 \(f(0)=f(2\pi)\) 平方可积函数空间. 取标准正交基
            \[\qty{e_n:=\frac{1}{\sqrt{2\pi}}\exp(\ii nx)}_{n\in\mathbb{Z}}\]
            那么对于任意 \(f\in L^2(\mathbb{T}[0,2\pi])\), 都有
            \[f=\sum_{n=-\infty}^{+\infty} \langle f, e_n\rangle e_n= \sum_{n=-\infty}^{+\infty} a_n \exp(\ii nx).\]
            其中 \(a_n\) 是 Fourier 系数, 满足
            \[a_n=\frac{1}{\sqrt{2\pi}}\int_0^{2\pi} f(x)\exp(-\ii nx) \dd x.\]
            用 Fourier 变换来表示, 就是
            \[\mathscr{F}:\qquad\begin{matrix}
                L^2(\mathbb{T}[0,2\pi]) & \to & \ell^2(\mathbb{Z}) \\
                f & \mapsto & (a_n)_{n\in\mathbb{Z}}:=\langle f, e_n\rangle_{n\in\mathbb{Z}}
            \end{matrix}\]
            则 \(\mathscr{F}\) 是 Hilbert 空间之间的\textbf{等距}(保持内积)\textbf{同构}(保持线性运算).
        \end{ins}

        \begin{thm}
            {具有可数标准正交基的 Hilbert 空间的充要条件}
            设 \(X\) 是 Hilbert 空间, 以下条件等价:
            \begin{enumerate}
                \item \(X\) 有可数的标准正交基;
                \item \(X\) 与 \(\ell^2\) 等距同构;
                \item \(X\) 上存在 Fourier 级数;
                \item \(X\) 是\hyperref[dfn:separable]{可分的}.
            \end{enumerate}
        \end{thm}

        \begin{prf}
            这是很直观的性质, 我们这里只介绍思路:\\
            取 \(X\) 的一个可数的标准正交基 \(\{e_n\}\), 定义映射
            \[T:\qquad\begin{matrix}
                X & \to & \ell^2 \\
                x & \mapsto & \qty{\langle x, e_n\rangle}_{n\in\mathbb{N}}
            \end{matrix}\]
            是一个等距同构; 这也正是 Fourier 级数的定义; Fourier 级数张成的线性空间在 \(X\) 中稠密, 也就是说可以用标准正交基的线性组合来逼近 \(X\) 中的任意元素.\\
            若存在可数的稠密子集, 则可以利用正交化过程得到可数的标准正交基.
        \end{prf}

        \begin{thm}
            [riesz-hilbert]
            {Hilbert 空间的 Riesz 表示定理}
            [Riesz Representation Theorem for Hilbert Spaces]
            \(X\) 是 Hilbert 空间, 则 \(X\cong X^*\). 其中的等距同构为:\\
            给定 \(f\in X^*\), 存在唯一的 \(x\in X\) 使得
            \[\forall y\in X, f(y)=\langle y,x\rangle\]
            且
            \[\norm{f}_{X^*}=\norm{x}_X.\]
        \end{thm}

        \begin{prf}
            我们仅考虑实的情况, 复数域上的情况类似.
            若 \(f=0\), 则取 \(x=0\) 即可. \\
            否则, \(Z:=\{y\in X: f(y)=0\}\), 则 \(Z=\ker f\) 是一个闭子空间. Hilbert 空间 \(X=Z\oplus Z^\perp\). 取 \(x_0\in Z^\perp\), 则 \(\forall\,y\in X\):
            \[y-\frac{f(y)}{f(x_0)x_0}\in Z\implies\left\langle y-\frac{f(y)}{f(x_0)}, x_0\right\rangle =0\implies f(y)=\frac{f(x_0)}{\norm{x_0}^2}\left\langle y,x_0\right\rangle  \]
            设 \(x=\frac{f(x_0)}{\norm{x_0}^2}x_0\), 则
            \[\forall y\in X, f(y)=\langle y,x\rangle.\]
            下面考虑等距性
            \[|f(y)|=|\left\langle y,x\right\rangle|\leq\norm{y}\norm{x} \implies \norm{f}_{X^*}\leq\norm{x}_X.\]
            等号取到的原因是范数的定义
            \[\norm{f}_{X^*}\geq\frac{|f(x)|}{\norm{x}}=\frac{|\langle x,x\rangle|}{\norm{x}}=\norm{x}.\]
        \end{prf}

\section{线性泛函分析基础}

    \subsection{Hahn-Banach 定理}

        \begin{dfn}
            [seminorm]
            {半范数}
            [Seminorm]
            数域 \(\F\) 上的线性空间 \(X\) 上的一个函数 \(p:X\to\R\) 是 \(X\) 上的一个\textbf{半范数}, 当且仅当:
            \begin{enumerate}[label=(\roman*)]
                \item \textbf{非负性}: \(p(x)\geq 0\), \(\forall x\in X\);
                \item \textbf{绝对齐次性}: \(p(\alpha x)=\abs{\alpha}p(x)\), \(\forall \alpha\in\F, x\in X\);
                \item \textbf{次可加性 / 三角不等式}: \(p(x+y)\leq p(x)+p(y)\), \(\forall x,y\in X\).
            \end{enumerate}
        \end{dfn}

        \begin{rmk}
            [Dedicatia]
            半范数与范数的区别在于, 半范数允许 \(p(x)=0\) 对某些非零的 \(x\) 成立. 也就是说, 半范数不要求满足正定性.
        \end{rmk}

        \begin{ppt}
            {半范数的零空间是线性子空间}
            [Kernel of Seminorm is a Linear Subspace]
            设 \(p\) 是 \(X\) 上的一个半范数. 定义 \(\ker p=\{x\in X: p(x)=0\}\). 则 \(\ker p\) 是 \(X\) 的一个线性子空间.
        \end{ppt}

        \begin{ppt}
            {半范数可以诱导范数}
            [Seminorms Induce Norms]
            设 \(p\) 是 \(X\) 上的一个半范数. 定义
            \(\norm{x + \ker p} = p(x)\)
            则 \(\norm{\cdot + \ker p}\) 是商空间 \(X/\ker p\) 上的一个范数.
        \end{ppt}

        \begin{ins}
            {\(\Lp\) 空间上的 \(\Lp\) 半范数诱导的范数}
            取等价类为几乎处处相等的函数, 则 $\Lp$ 半范数在 $\Lp$ 空间上成为范数. 这也是我们一般认为的 $\Lp$ 空间的定义. 严格上认为 \(\Lp\) 空间中的元素是函数的等价类而不是函数本身.
        \end{ins}

        \begin{dfn}
            [sublinear-functional]
            {次线性泛函}
            [Sublinear Functional]
            设 \(X\) 是一个线性空间. 定义 \(p:X\to\R\) 是一个\textbf{次线性泛函}, 当且仅当:
            \begin{enumerate}[label=(\roman*)]
                \item \textbf{正齐次性}: \(p(\alpha x)=\alpha p(x)\), \(\forall \alpha\geq 0, x\in X\);
                \item \textbf{次可加性 / 三角不等式}: \(p(x+y)\leq p(x)+p(y)\), \(\forall x,y\in X\).
            \end{enumerate}
        \end{dfn}

        \begin{thm}
            [hahn-banach]
            {Hahn-Banach 定理}
            [Hahn-Banach Theorem]
            设 \(X\) 是一个线性赋范空间, \(p\) 是 \(X\) 上的一个次线性泛函, \(Y\) 是 \(X\) 的一个线性子空间, \(f:Y\to\R\) 是一个线性泛函且满足
            \[f(y)\leq p(y), \quad \forall y\in Y.\]
            则存在一个线性泛函 \(F:X\to\R\) 使得
            \[F(y)=f(y), \quad \forall y\in Y,\]
            且
            \[F(x)\leq p(x), \quad \forall x\in X.\]
        \end{thm}

        \begin{rmk}
            [Dedicatia]
            这个定理的意义在于, 在满足某些条件的情况下, 可以将一个定义在子空间上的线性函数(i)扩展到整个空间上, (ii)同时保持被某个次线性函数的控制.
        \end{rmk}

        \begin{prf}
            先考虑一步延拓, 设 \(x_0\in X\setminus Y\). 我们要把 \(f\) 从 \(Y\) 延拓到 \(Y\oplus\spn(x_0)\) 上. 为此考虑
            \[\tilde{f} (y+\alpha x_0)=f(y)+\alpha\cdot \tilde{f}(x_0)\]
            我们只需要确定 \(\tilde{f}(x_0)\) 的值. 由于 \(\tilde{f}\) 需要满足 \(\tilde{f}(y+\alpha x_0)\leq p(y+\alpha x_0)\), 则对于任意 \(y\in Y\) 和 \(\alpha\in\R\), 都有:\\
            (1) 如果 \(\alpha=0\), 显然成立;\\
            (2) 如果 \(\alpha>0\), 则
            \[f(x_0)\leq\frac{p(y+\alpha x_0)-f(y)}{\alpha}=\frac{p(y+\alpha x_0)}{\alpha}-\frac{f(y)}{\alpha} \]
            等价于
            \[\tilde{f}(x_0)\leq p(z+x_0)-f(z)\]
            其中 \(z=\frac{y}{\alpha}\);\\
            (3) 如果 \(\alpha<0\), 则同理可得
            \[\tilde{f}(x_0)\geq p(z+x_0)-f(z)\]
            这就说明了
            \[\sup_{z\in Y}\{-f(z)+p(z+x_0)\}\leq \tilde{f}(x_0)\leq\inf_{z\in Y}\{p(z+x_0)-f(z)\}\]
            这等价于
            \[\forall y_1,y_2\in Y, \quad -f(y_2)+p(y_2+x_0)\leq p(y_1+x_0)-f(y_1)\]
            只要证明
            \[\sup_{z\in Y}\{-f(z)+p(z+x_0)\}\leq\inf_{z\in Y}\{p(z+x_0)-f(z)\}\]
            就说明了 \(\tilde{f}(x_0)\) 的取值可以被夹出来. 由于
            \[p(y_1+x_0)+p(y_2+x_0)-f(y_1)-f(y_2)\geq p(y_1+y_2)-f(y_1+y_2)\geq 0\]
            从而该延拓是一定存在的. \\
            构造集合 \(\mathcal{S}:=\{(Y_i, g_i)\}\). 其中 \(Y_i\) 是 \(X\) 的某个子空间, \(g_i\) 是 \(Y_i\) 上的一个线性泛函, 满足 \(Y\subseteq Y_i\), \(g_i|_Y=f\), 且 \(g_i(x)\leq p(x)\) 对所有 \(x\in Y_i\) 成立. \\
            这是一个偏序集, 其中的偏序定义为: \((Y_i, g_i)\leq (Y_j, g_j)\) 当且仅当 \(Y_i\subseteq Y_j\) 且 \(g_j|_{Y_i}=g_i\). 设 \(\mathcal{C}\) 是 \(\mathcal{S}\) 中的一个全序子集. 则 \(\mathcal{C}\) 中的元素可以被包含在一个更大的元素中. 也就是说, 定义
            \[Y_0=\bigcup_{(Y_i, g_i)\in\mathcal{C}} Y_i, \quad g_0:Y_0\to\R, g_0(x)=g_i(x) \text{ if } x\in Y_i\]
            则 \((Y_0, g_0)\) 是 \(\mathcal{S}\) 中的一个元素, 且包含了 \(\mathcal{C}\) 中的所有元素. 根据 Zorn 引理, \(\mathcal{S}\) 中存在一个极大元 \((Y_0, g_0)\). 如果 \(Y_0\neq X\), 则可以通过上面的一步延拓来得到一个更大的元素, 矛盾. 因此 \(Y_0=X\), \(g_0\) 就是我们要找的 \(F\). 
        \end{prf}

        \begin{rmk}
            [Dedicatia]
            Hahn-Banach 定理只适用于对泛函的延拓, 不适用于对任意线性算子的延拓.
        \end{rmk}

        \begin{thm}
            {赋范延拓}
            [Norming Extension]
            设 $X$ 是赋范线性空间, $Y$ 是 $X$ 的子空间, $f$ 是定义在 $Y$ 上的线性泛函, $\exists F \in \blo (X, \R)$ 使得 $F$ 是 $f$ 的延拓, 且 $\norm{F}=\norm{f}$.
        \end{thm}

        \begin{prf}
            定义
            \[p(x)=\norm{f}\norm{x}, \quad \forall x\in X.\]
            则 \(p\) 是 \(X\) 上的一个次线性泛函, 且对于任意 \(y\in Y\), \(f(y)\leq\norm{f}\norm{y}=p(y)\). 根据 Hahn-Banach 定理, 存在一个线性泛函 \(F:X\to\R\) 使得 \(F|_Y=f\) 且 \(F(x)\leq p(x)\) 对所有 \(x\in X\) 成立. 从而 \(\norm{F}\leq\norm{f}\). \\
            另一方面, 因为 \(F|_Y=f\), 所以 \(\norm{f}\leq\norm{F}\). 因此 \(\norm{F}=\norm{f}\).
        \end{prf}

        \begin{thm}
            [norming-functional]
            {赋范泛函}
            [Norming Functional]
            设 \(X\) 是赋范线性空间, \(\forall x\neq 0\), \(\exists\,f\in X^*\st f(x)=\norm{x}\) 且 \(\norm{f}=1\).
        \end{thm}

        \begin{prf}
            设 \(f\) 定义在 \(\spn(x)\) 上且满足 \(f(\alpha x)=\alpha\norm{x}\). 则根据 Hahn-Banach 定理, 可以将其延拓到整个空间上且保持范数不变. 从而存在 \(F\in X^*\) 使得 \(F(x)=\norm{x}\) 且 \(\norm{F}=\norm{f}=1\).
        \end{prf}

        \begin{rmk}
            [Dedicatia]
            这个定理保证了 $X$ 的对偶空间足够丰富, 以至于每个非零向量都可以通过某个对偶空间中的元素来“测量”其范数. 
        \end{rmk}

        \begin{crl}
            {范数的等价定义}
            [Equivalent Definition of Norm]
            设 \(X\) 是一个线性空间, \(\norm{\cdot}\) 是 \(X\) 上的一个函数. 则 \(\norm{\cdot}\) 是 \(X\) 上的一个范数, 当且仅当
            \[\forall x\in X, \norm{x}=\sup_{f\in X^*, \norm{f}=1} |f(x)|.\]
        \end{crl}

        \begin{prf}
            \(|f(x)|\leq\norm{f}\norm{x}\) 成立, 所以 \(\sup_{f\in X^*, \norm{f}=1} |f(x)|\leq\norm{x}\). 另一方面, 根据上一个定理, 存在 \(f\in X^*\) 使得 \(f(x)=\norm{x}\) 且 \(\norm{f}=1\). 从而等号取到, \(\sup_{f\in X^*, \norm{f}=1} |f(x)|=\norm{x}\).
        \end{prf}

    \subsection{【集合论】序结构与 Zorn 引理}

        \begin{dfn}
            {偏序结构 / 预序结构}
            [Partial Order / Preorder]
            设 \(X\) 是一个集合. 定义 \(X\) 上的一个\textbf{偏序结构 / 预序结构} 是一个二元关系 \(\leq\), 满足:
            \begin{enumerate}[label=(\roman*)]
                \item \textbf{自反性}: \(x\leq x\), \(\forall x\in X\);
                \item \textbf{反对称性}: \(x\leq y\land y\leq x\implies x=y\), \(\forall x,y\in X\);
                \item \textbf{传递性}: \(x\leq y\land y\leq z\implies x\leq z\), \(\forall x,y,z\in X\).
            \end{enumerate}
        \end{dfn}

        \begin{dfn}
            {偏序集}
            [Partially Ordered Set]
            设 \(X\) 是一个集合, \(\leq\) 是 \(X\) 上的一个偏序结构. 则 \((X, \leq)\) 是一个\textbf{偏序集}.
        \end{dfn}

        \begin{dfn}
            {全序结构 / 线性序结构}
            [Total Order / Linear Order]
            设 \(X\) 是一个集合. 定义 \(X\) 上的偏序结构成为\textbf{全序结构 / 线性序结构}, 当且仅当对于任意 \(x,y\in X\), \(x\leq y\) 或 \(y\leq x\) 至少有一个成立. 
        \end{dfn}

        \begin{dfn}
            {全序集}
            [Totally Ordered Set]
            设 \(X\) 是一个集合, \(\leq\) 是 \(X\) 上的一个全序结构. 则 \((X, \leq)\) 是一个\textbf{全序集}.
        \end{dfn}

        \begin{rmk}
            [Dedicatia]
            直观地理解:
            \begin{itemize}
                \item 偏序结构: 只是元素之间的二元关系;
                \item 全序结构: 任何两个元素之间都有二元关系, 也就是说, 任意两个元素之间可以比较大小.
            \end{itemize}
        \end{rmk}

        \begin{dfn}
            {上界 / 下界}
            [Upper Bound / Lower Bound]
            设 \((X, \leq)\) 是一个偏序集, \(A\subseteq X\). 定义 \(x\in X\) 是 \(A\) 的一个\textbf{上界}, 当且仅当 \(\forall a\in A, a\leq x\). \\
            定义 \(x\in X\) 是 \(A\) 的一个\textbf{下界}, 当且仅当 \(\forall a\in A, x\leq a\).
        \end{dfn}

        \begin{dfn}
            {极大元 / 极小元}
            [Maximal Element / Minimal Element]
            设 \((X, \leq)\) 是一个偏序集, \(A\subseteq X\). 定义 \(x\in A\) 是 \(A\) 的一个\textbf{极大元}, 当且仅当 \(\forall y\in A, x\leq y\implies x=y\). \\
            定义 \(x\in A\) 是 \(A\) 的一个\textbf{极小元}, 当且仅当 \(\forall y\in A, y\leq x\implies x=y\).
        \end{dfn}

        \begin{rmk}
            [Dedicatia]
            注意, 这定义并不要求
            \[\exists\,x\in A\st \forall y\in A, x\leq y\]
            这是因为偏序并不保证任意元素之间都有二元关系. 所以可能存在多个极大元.
        \end{rmk}

        \begin{thm}
            [zorn]
            {Zorn 引理}
            [Zorn's Lemma]
            设 \((X, \leq)\) 是一个\textit{偏序集}. 如果 \(X\) 中的每个\textit{全序子集}都有\textit{上界}, 则 \(X\) 中存在一个\textit{极大元}.
        \end{thm}

        \begin{thm}
            {良序定理}
            [Well-Ordering Theorem]
            任意集合都可以被全序化. \\
            设 \(X\) 是一个集合. 则存在一个全序结构 \(\leq\) 使得 \((X, \leq)\) 是一个全序集.
        \end{thm}

        \begin{rmk}
            [Dedicatia]
            Zorn 引理和良序定理都是等价于选择公理的. 
        \end{rmk}

        \begin{thm}
            {所有向量空间都有基}
            设 \(V\) 是任意向量空间, 则 \(V\) 有一组基.
        \end{thm}

        \begin{prf}
            设 \(\mathcal{S}\) 是 \(V\) 中所有线性无关的集合的集合, 则 \(\mathcal{S}\) 是偏序集, 其中的偏序关系是 \(\subseteq\). \\
            根据 Zorn 引理, \(\mathcal{S}\) 中存在全序子集 \(\{L_i\}_{i\in I}\). 不妨设
            \[L:=\bigcup_{\alpha\in I}L_\alpha\]
            则我们要证明 \(L\) 是线性无关的. 从 \(L\) 中取有限个元素 \(\{v_1,v_2, \cdots v_n\}\):
            \[\forall i\in\{1,\cdots,n\},\exists \alpha_i\in I,\st v_i\in L_{\alpha_i}\]
            所以
            \[\exists j \in \{1, 2, \cdots, n\}, \st \alpha_j \geq \alpha_i, \forall i \in \{1, 2, \cdots, n\}. \]
            所以 $v_1, v_2, \cdots, v_n \in L_{\alpha_j}$. 这说明它们是线性无关的. \(L\) 的任意有限子集都是线性无关的, 从而 \(L\) 是线性无关的. \\
            因此 \(L\) 是 \(\mathcal{S}\) 中的一个上界. 根据 Zorn 引理, \(\mathcal{S}\) 中存在一个极大元 \(B\). 则 \(B\) 是 \(V\) 的一组基.
        \end{prf}

        \begin{dfn}
            [hamel-basis]
            {Hamel 基}
            [Hamel Basis]
            设 \(V\) 是一个向量空间. 定义 \(V\) 的\textbf{Hamel 基} 是 \(V\) 的依赖 Zorn 引理得到的一组基.
        \end{dfn}

        \begin{rmk}
            [Dedicatia]
            Hamel 基的定义不能给出构造性的、具体的形式, 因为 Zorn 引理是非构造性的. 这也是为什么我们很少使用 Hamel 基去研究分析学.
        \end{rmk}

        \begin{ins}
            {存在满足可加性但不满足齐次性的函数}
            在 \(\R\) 中, \(1\) 和 \(\sqrt{2}\) 是线性无关的. 我们可以由此得到一组 Hamel 基 \(\{e_i\}_{i\in I}\). 其中 \(e_1=1\), \(e_2=\sqrt{2}\). 定义函数 \(f:\R\to\R\):
            \[f(1)=\sqrt{2}, f(\sqrt{2})=1, f(e_i)=e_i, \forall i \in I\setminus\{1,\sqrt{2}\}. \]
            则 \(f\) 满足可加性: 只需按分量定义 \(f\) 的值即可. 但是 \(f\) 不满足齐次性. 
        \end{ins}

        \begin{thm}
            {无限维 Banach 空间的 Hamel 基是不可数的}
            设 \(X\) 是一个无限维的 Banach 空间, 则 \(X\) 的任何 Hamel 基都是不可数的.
        \end{thm}

        \begin{prf}
            设 \(X\) 是 Banach 空间, \(\dim X=\infty\), \(X\) 具有可数的 Hamel 基 \(\{e_i\}_{i=1}^{\infty}\). \\
            那么 \(X=\spn(\{e_i\})\). 
            \[X=\bigcup_{i=1}^{\infty}V_n\]
            \(V_n\) 作为 \(X\) 的闭子空间显然不可能是无处稠密的, 设 $V_i \supset B(r,x_0)$. 则 $V_i \supseteq\{x-y: x,y\in B(r,x_0)\} \supset B(0,r)$. 得到
            \[V_i\supset\bigcup_{a\in\R}\{ax: x\in B(0,r)\} \]
            这说明 \(V_i=X\). 矛盾.
        \end{prf}

        \begin{xmp}
            {序列空间 \(\ell_{00}\) 不是 Banach 空间}
            只有有限项非零项的序列组成的空间 \(\ell_{00}\) 不是 Banach 空间. 因为 \(\ell_{00}\) 的 Hamel 基是可数的, 但是 \(\ell_{00}\) 是无限维的线性空间.
        \end{xmp}

        \begin{rmk}
            这说明在 Hamel 基的意义下 Banach 空间要么是有限维的, 要么是不可数维的. 这是强完备性的体现.
        \end{rmk}

        \begin{dfn}
            [ordinal]
            {序数}
            [Ordinals]
            序数是所有按序同构的良序集的等价类. 记作 \(\ordinal\).
        \end{dfn}
        
        \begin{rmk}
            [Dedicatia]
            序数的定义从直观上说就是一个元素的位置的编号. 但是为了描述无限集合的结构, 有一些特殊的编号. 这也正是我们定义序数的动机.
        \end{rmk}

        \begin{ppt}
            {序数集合是全序的}
        \end{ppt}

        \begin{dfn}
            {前驱序数 / 后继序数}
            [Predecessor / Successor Ordinal]
            设 \(\alpha\) 是一个序数. 定义 \(\alpha\) 的\textbf{前驱序数}是 \(\alpha-1\), 定义 \(\alpha\) 的\textbf{后继序数}是 \(\alpha+1\).
        \end{dfn}

        \begin{dfn}
            {极限序数}
            设 \(\alpha\in\ordinal\) 且 \(\alpha\neq 0\), 定义 \(\alpha\) 是一个\textbf{极限序数}, 当且仅当 \(\alpha\) 不是任何序数的后继. \\
            也就是说, \(\alpha\) 是一个极限序数当且仅当 \(\alpha\) 可以表示为一个小于 \(\alpha\) 的序数的集合的并.
        \end{dfn}

        \begin{ins}
            {第一个极限序数 \(\omega\)}
            设 \(\omega\) 是所有自然数的集合. 则 \(\omega\) 是第一个极限序数.
        \end{ins}

        \begin{rmk}
            [Dedicatia]
            事实上, 还有更多的从 \(\omega\) 得到的极限序数, 例如 \(\omega+\omega\), \(\omega^2\) 等等. 这些极限序数的结构非常复杂. \\
            由于 \(\N\) 是可数的, 我们还可以在不可数集上定义序数, 这就是不可数序数 \(\omega_1, \omega_2, \cdots\).
        \end{rmk}

        \begin{vardfn}
            [cardinal]
            {基数}
            [Cardinal]
            基数是所有按双射同构的集合的等价类.
        \end{vardfn}

        \begin{ins}
            {可数无穷基数}
            设 \(\aleph_0\) 是所有可以与自然数 \(\N\) 建立双射关系的集合. \(\aleph_0\) 是第一个可数无穷基数.
        \end{ins}

        \begin{ins}
            {不可数无穷基数 / 连续统基数}
            设 \(\mathfrak{c}\) 是所有可以与实数 \(\R\) 建立双射关系的集合. \(\mathfrak{c}\) 是\textbf{不可数无穷基数}, 也叫\textbf{连续统基数(Continuum Cardinal)}.
        \end{ins}

        \begin{rmk}
            [Dedicatia]
            \textbf{连续统假设}\; 不存在一个基数 \(\kappa\) 使得 \(\aleph_0 < \kappa < \mathfrak{c}\). \\
            也就是说, 没有一个集合的基数介于可数无穷和连续统之间. 这个假设是独立于 ZFC 公理系统的, 也就是说, 无论我们接受还是拒绝这个假设, 都不会导致矛盾. Godel 证明了连续统假设在 ZFC 中是无法被否定的, Cohen 证明了连续统假设在 ZFC 中是无法被证明的.
        \end{rmk}

        \begin{thm}
            {超限归纳法}
            [Transfinite Induction]
            设 \(P(\alpha)\) 是一个关于序数 \(\alpha\) 的命题. 如果满足以下条件:
            \begin{enumerate}[label=(\roman*)]
                \item \(P(0)\) 成立;
                \item 对于任意序数 \(\alpha\), 如果 \(P(\beta)\) 对所有 \(\beta<\alpha\) 都成立, 则 \(P(\alpha)\) 也成立.
            \end{enumerate}
            则 \(P(\alpha)\) 对所有序数 \(\alpha\) 都成立.
        \end{thm}

        \begin{rmk}
            [Dedicatia]
            \textbf{超限归纳法对于序数 \(\ordinal\) 的关系就好比普通数学归纳法对于自然数 \(\N\) 的关系}.\\
            超限归纳法是对普通数学归纳法的推广. 可以用于任意序数. 这使得我们可以在更广泛的范围内进行归纳证明.
        \end{rmk}

        \begin{ins}
            {Bernstein 集合}
            \(\R\) 上存在集合 \(A\) 使得 \(A\) 和 \(A^c\) 都满足以下条件:
            \begin{enumerate}[label=(\roman*)]
                \item 与 \(\R\) 的任意完全集交集非空;
                \item 不包含任何完全集.
            \end{enumerate}
            这是基于以下两条事实:
            \begin{itemize}
                \item \(\R\) 中完全集的个数是 \(\mathfrak{c}\);
                \item 每一个完全集中也有 \(\mathfrak{c}\) 个元素.
            \end{itemize}
            证明时, 将 \(\R\) 中的所有完全集赋予序数, 然后使用超限归纳: 在序数为 \(\alpha\) 步的时候, 处理完全集 \(P_\alpha\):\\
            1. 从第 \(\alpha\) 个完全集中取一个元素加入 \(A\);\\
            2. 从第 \(\alpha\) 个完全集中取一个元素加入 \(A^c\).\\
            事实上我们不用担心在某一步的时候完全集中的元素已经被之前的步骤取完了, 因为每个完全集中都有 \(\mathfrak{c}\) 个元素, 而我们在每一步只取一个元素. 因此我们总是可以从第 \(\alpha\) 个完全集中取出一个元素加入 \(A\) 和 \(A^c\). 重复直到所有完全集都被处理完了. 最后得到的集合 \(A\) 和 \(A^c\) 就满足要求了:\\
            (1) 对任意完全集 \(P_\alpha\), \(A\) 和 \(A^c\) 都与 \(P_\alpha\) 有交集;\\
            (2) \(A\) 和 \(A^c\) 都不包含任何完全集, 因为每个完全集都必然有元素在 \(A\) 中, 也有元素在 \(A^c\) 中. 如果包含了完整的完全集就会导致矛盾.\\
            这就得到了 \(A\). \(A\) 被称为 \textbf{Bernstein 集合}.
            \tcbline
            无论是分析还是拓扑, Bernstein 集合都是非常奇怪的: 它有如下性质:
            \begin{itemize}
                \item \(A\) 和 \(A^c\) 都在 \(\R\) 中稠密, 它的所有元素都是极限点, 但 \(A\) 和 \(A^c\) 都不包含任意区间;
                \item \(A\) 和 \(A^c\) 都具有零内测度, 这是因为如果 \(A\) 或 \(A^c\) 有正的内测度, 那么它就包含一个完全集, 这与 \(A\) 和 \(A^c\) 都不包含任何完全集的性质矛盾. 因而 \(A\) 与 \(A^c\) 都是不可测集;
                \item \(A\) 和 \(A^c\) 都既不是第一纲集, 也不是余纲集.
            \end{itemize}
        \end{ins}

    \subsection{凸集分离定理}

        \begin{intrormk}
            [Dedicatia]
            本节讲述的是 Hahn-Banach 定理的几何形式, 也就是凸集分离定理.在这之前我们需要一些概念和定义.
        \end{intrormk}

        \begin{ppt}
            {线性泛函可以精准区分向量}
            [Linear Functionals Separate Points]
            设 \(X\) 是 Banach 空间, 则
            \[\forall\,x,y\in X, x\neq y\implies\exists\,f\in X^*\st f(x)\neq f(y).\]
        \end{ppt}

        \begin{prf}
            只需证明 \(\forall x\in X, x\neq 0\implies\exists f\in X^*\st f(x)\neq 0\). 存在赋范泛函 \(f\in X^*\) 使得 \(f(x)=\norm{x}\neq 0\).
        \end{prf}

        \begin{ppt}
            {线性泛函可以精准区分点和闭子空间}
            [Linear Functionals Separate Points and Closed Subspaces]
            设 \(X\) 是 Banach 空间, \(Y\) 是 \(X\) 的一个闭子空间, \(x\in X\setminus Y\), 则
            \[\exists f\in X^*\st \norm{f}=1, f(x)\neq 0, f(y)=0, \forall y\in Y.\]
        \end{ppt}

        \begin{prf}
            设 
            \[d=\dist(x_0, Y)=\inf_{y\in Y}\norm{x_0-y}.\]
            设
            \[g:Y \oplus \spn (x_0)=: Z \to \R, g(y+\alpha x_0) = \alpha d\]
            则 \(g\) 是 \(Z\) 上的一个线性泛函, 且 \(\forall y\in Y, \alpha\in\R\), 都有
            \[g(y+\alpha x_0)=\alpha d\leq \alpha\norm{x_0-y}\leq\norm{y+\alpha x_0}.\]
            根据 Hahn-Banach 定理, 存在 \(f:X\to\R\) 使得 \(\forall y\in Y, \alpha\in\R\), 都有 \(f(y+\alpha x_0)\leq\alpha\norm{y+\alpha x_0} \), \(f(y)=0\), \(f(x_0)=d>0\). 这就是我们要的线性泛函.\\
            考虑 \(\forall\epsilon>0\), 设 \(y\in Y\) s.t. \(\norm{x-y}\leq\eta+d\), 其中 \(\eta=\frac{d\epsilon}{1-\epsilon}\), 这样的 \(y\) 是可以找到的. 则
            \[\frac{|f(x_0-y)|}{\norm{x_0-y}}=\frac{d}{\norm{x_0-y}}>\frac{d}{d+\eta}\geq 1-\epsilon\] 
            由于 \(\epsilon\) 是任意的, 所以 \(\norm{f}\geq 1\). 另一方面因为 \(f(y+\alpha x_0)\leq\alpha\norm{y+\alpha x_0}\), 所以 \(\norm{f}\leq 1\). 因此 \(\norm{f}=1\).
        \end{prf}

        \begin{dfn}
            [absorbing]
            {吸收集}
            [Absorbing Set]
            设 \(X\) 是一个线性空间. 定义 \(A\subseteq X\) 是一个\textbf{吸收集}, 当且仅当
            \[\forall\,x\in X, \exists\,\lambda>0\st x\in\lambda A:=\left\{\lambda a : a\in A\right\}\]
        \end{dfn}

        \begin{dfn}
            [minkowski-functional]
            {Minkowski 泛函}
            [Minkowski Functional]
            设 \(X\) 是一个线性空间, \(A\subseteq X\) 是一个吸收集. 定义 \(A\) 的 Minkowski 泛函 \(\mu_A:X\to\R\) 为
            \[\mu_A(x)=\inf\{\lambda>0 : x\in\lambda A\}.\]
        \end{dfn}

        \begin{dfn}
            [balanced]
            {平衡集 / 均衡集}
            [Balanced Set]
            设 \(X\) 是一个线性空间. 定义 \(A\subseteq X\) 是一个\textbf{平衡集}, 当且仅当
            \[\forall\,x\in A, \forall\,\lambda\in\R, |\lambda|\leq 1\implies \lambda x\in A.\]
        \end{dfn}

        \begin{dfn}
            [convex]
            {凸集}
            [Convex Set]
            设 \(X\) 是一个线性空间. 定义 \(A\subseteq X\) 是一个\textbf{凸集}, 当且仅当
            \[\forall\,x,y\in A, \forall\,t\in[0,1], tx+(1-t)y\in A.\]
        \end{dfn}

        \begin{rmk}
            [Dedicatia]
            直观上理解, 这三类集合的定义是这样的:
            \begin{itemize}
                \item 吸收集是指对于空间中的任意向量, 都存在一个足够大的标量使得该向量被包含在这个标量倍的集合中. 这意味着吸收集可以通过标量缩放覆盖整个空间.
                \item 平衡集是指对于集合中的任意元素其乘以一个绝对值不超过1的标量时, 结果仍然在集合中. 这意味着均衡集中的点不会因为缩放而离开集合.
                \item 凸集是指对于集合中的任意两个元素, 连接它们的线段也完全包含在集合中, 凸集中的点之间的线性组合仍然在集合中.
            \end{itemize}
            需要注意的是, 这三类集合都是定义在线性空间上的. 如果没有线性空间结构就无法定义这些集合.
        \end{rmk}

        \begin{dfn}
            [absolutely-convex]
            {绝对凸集}
            [Absolute Convex Set]
            设 \(X\) 是一个线性空间. 定义 \(A\subseteq X\) 是一个\textbf{绝对凸集}, 当且仅当 \(A\) 同时是一个吸收集、平衡集、凸集.
        \end{dfn}

        \begin{thm}
            {半范数诱导的单位开球是绝对凸的}
            [Unit Open Ball Induced by Seminorm is Absolutely Convex]
            设 $p$ 是赋范线性空间 $X$ 上的半范数. 则单位开球 $B = \{x \in X: p(x) < 1\}$ 是一个绝对凸集. 并且半范数等于 Minkowski 泛函: $p = \mu_B$.
        \end{thm}

        \begin{prf}
            1. \(B\) 是平衡的: \(x\in B\implies p(x)<1\), \(\forall t\in\R\) 其中 \(|t|\leq 1\), \(p(tx)=|t|p(x)\leq p(x)<1\), 说明 \(tx\in B\).\\
            2. \(B\) 是凸的: \(\forall\,x,y\in B, t\in (0,1)\), \(p(tx+(1-t)y)\leq p(tx)+(1-t)p(y)<1\), 说明 \(tx+(1-t)y\in B\).\\
            3. \(B\) 是吸收的: 取 \(s>p(x)\). \(p(\frac{x}{s})=\frac{1}{s}p(x)<1\), 这说明 \(\frac{x}{s}\in B\), 所以 \(x\in sB\).\\
            4. \(p=\mu_B\): 之前我们已经证明 \(\mu_B\leq p\). 考虑如果 \(a\leq\mu_B(x)\), 那么 \(p(\frac{x}{a})\geq 1\), 这说明 \(x\notin aB\implies a\leq\mu_B(x)\). 因此 \(p=\mu_B\).
        \end{prf}

        \begin{thm}
            {吸收的凸集的 Minkowski 泛函是次线性泛函}
            [Minkowski Functional of Absorbing Convex Set is Sublinear]
            设 \(X\) 是一个线性空间, \(A\subseteq X\) 是一个吸收的凸集. 则 \(A\) 的 Minkowski 泛函 \(\mu_A:X\to\R\) 是一个次线性泛函.
        \end{thm}

        \begin{prf}
            只证明次可加性: \(\mu_A(x+y)\leq\mu_A(x)+\mu_A(y)\). 设 \(a=\mu_A(x)+\epsilon, b=\mu_A(y)+\epsilon\). 则 \(x\in aA, y\in bA\). 所以
            \[\frac{x+y}{a+b}=\frac{x}{a}\frac{a}{a+b}+\frac{y}{b}\frac{b}{a+b} \in\frac{a}{a+b}A+\frac{b}{a+b}A\subseteq A.\]
            这说明 \(x+y\in (a+b)A\). 所以 \(\mu_A(x+y)\leq a+b=\mu_A(x)+\mu_A(y)+2\epsilon\). 由于 \(\epsilon\) 是任意的, 所以 \(\mu_A(x+y)\leq\mu_A(x)+\mu_A(y)\).
        \end{prf}

        \begin{thm}
            [convex-set-separation]
            {凸集分离定理}
            [Convex Set Separation Theorem]
            设 \(X\) 是 \(\R\) 上的赋范线性空间. \(A, B\) 是两个非空的凸集, \(A\cap B=\varnothing\). \(A\) 是开集, 那么 
            \[\exists\,f\in X^*, \alpha\in\R \st f(a)<\alpha\leq f(b), \forall\,a\in A, b\in B.\]
        \end{thm}

        \begin{prf}
            固定 \(a_0\in A, b_0\in B\). 那么令 \(x_0=b_0-a_0, x_0\neq 0\). 定义集合 \(O=A-B+x_0\). 那么 \(x_0\notin O\), \(O\) 是一个凸的开集且是原点的邻域.\\
            则 \(O\) 包含了一个原点的开球, 所以 \(O\) 是吸收的, 于是 \(\mu_O\) 是一个次线性泛函, \(\mu_O(x_0)\geq 1\), 因为 \(x_0\notin O\). \\
            考虑
            \[f_0: \spn(\{x_0\})\to\R, f(\lambda x_0)=\lambda  \]
            则 \(f_0(x_0)=1\leq \mu_O(x_0)\). 使用 Hahn-Banach 定理进行延拓, \(f_0\) 被延拓到整个空间上得到 \(f\). 则 \(\forall a\in A, b\in B\), 
            \[f(a)-f(b)= f(a-b+x_0)-f(0)= f(a-b+x_0)-1\leq \mu_O(a-b+x_0)-1< 0\]
            从而 \(f(a)<f(b)\). 这就是我们要的线性泛函.
        \end{prf}

        \begin{thm}
            {紧凸集分离定理}
            [Compact Convex Set Separation Theorem]
            设 \(X\) 是 \(\R\) 上的赋范线性空间. \(A, B\) 是两个非空的凸集, \(A\cap B=\varnothing\). \(A\) 是紧集, \(B\) 是闭集, 那么
            \[\exists\,f\in X^*, \alpha\in\R \st f(a)<\alpha< f(b), \forall\,a\in A, b\in B.\] 
        \end{thm}

        \begin{prf}
            由于二者都是闭的, 可以定义距离
            \[\dist(A,B)=\delta\]
            则 \(\delta>0\). 取
            \[\tilde{A}:=\{x\in X: \dist(x,A)<\delta/2\}\]
            使用一般的凸集分离定理证明.
        \end{prf}

    \subsection{【点集拓扑】Baire 纲定理}

        \begin{dfn}
            [metric]
            {度量}
            [Metric]
            设 \(X\) 是一个集合. 定义 \(X\) 上的一个\textbf{度量}是一个函数 \(d:X\times X\to\R\), 满足:
            \begin{enumerate}[label=(\roman*)]
                \item \textbf{正定性}: \(d(x,y)\geq 0\), \(d(x,y)=0\iff x=y\);
                \item \textbf{对称性}: \(d(x,y)=d(y,x)\), \(\forall x,y\in X\);
                \item \textbf{次可加性 / 三角不等式}: \(d(x,z)\leq d(x,y)+d(y,z)\), \(\forall x,y,z\in X\).
            \end{enumerate}
        \end{dfn}

        \begin{dfn}
            [ms]
            {度量空间}
            [Metric Space]
            设 \(X\) 是一个集合, \(d\) 是 \(X\) 上的一个\hyperref[dfn:metric]{度量}. 则 \((X, d)\) 是一个\textbf{度量空间(Metric Space)}.
        \end{dfn}

        \begin{rmk}
            [Dedicatia]
            可以证明, 由度量可以定义拓扑, 即度量意义下的开球是度量空间的拓扑基, 二者是相容的. 因此我们将不再重复拓扑空间中的各种概念的定义.
        \end{rmk}

        \begin{ppt}
            {赋范线性空间是度量空间}
            [Normed Vector Space is a Metric Space]
            范数可以诱导度量: 具体来说, 对于赋范线性空间 \((X, \norm{\cdot})\), \(x,y\in X\), 定义
            \[d(x,y)=\norm{x-y}\]
            是 \(X\) 上的一个度量函数, 从而 \(X\) 也是一个度量空间.
        \end{ppt}

        \begin{rmk}
            [Dedicatia]
            一般的度量不一定能诱导范数.
        \end{rmk}

        \begin{dfn}
            [dense]
            {稠密集}
            [Dense Set]
            设 \(X\) 是一个拓扑空间, \(A\subseteq X\). 定义 \(A\) 是 \(X\) 中的一个\textbf{稠密集}, 当且仅当 \(\overline{A}=X\).  \\
            也就是说, \(A\) 的闭包是整个空间.
        \end{dfn}

        \begin{rmk}
            [Dedicatia]
            等价地理解, \(A\) 在 \(X\) 中稠密, 即是 \(X\) 中的任意点都可以被 \(A\) 中的点任意接近. 这也是为什么叫稠密集的原因.
        \end{rmk}

        \begin{dfn}
            [nowhere-dense]
            {无处稠密}
            [Nowhere Dense]
            设 \(X\) 是一个拓扑空间, \(A\subseteq X\). 定义 \(A\) 是 \(X\) 中的一个\textbf{无处稠密}, 当且仅当 \(\interior(\overline{A})=\varnothing\). 
        \end{dfn}

        \begin{thm}
            {完备化定理}
            对任意度量空间 \(X\), 都存在一个\hyperref[dfn:complete]{完备}度量空间 \(Y\) 使得 \(X\) 是 \(Y\) 的一个稠密度量子空间.
        \end{thm}

        \begin{ins}
            {\(Q\) 在 \(\R\) 中稠密}
            这是因为对任意实数, 都能找到一列有理数收敛到它. 从而 \(\overline{\Q}=\R\).\\
            事实上, \(\R\) 就是 \(\Q\) 的完备化度量空间.
        \end{ins}

        \begin{dfn}
            [meager]
            {Baire 第一纲集}
            [Baire 1st Category Set / Meager Set]
            设 \(X\) 是一个拓扑空间. 定义 \(X\) 中的一个子集 \(A\) 是一个\textbf{Baire 第一纲集}, 当且仅当 \(A\) 可以表示为可数个无处稠密集的并集:
            \[A=\bigcup_{n=1}^\infty A_n, \interior(\overline{A_n})=\varnothing.\]
        \end{dfn}

        \begin{dfn}
            [residual]
            {余纲集}
            [Residual Set / Comeager Set]
            设 \(X\) 是一个拓扑空间. 定义 \(X\) 中的一个子集 \(A\) 是一个\textbf{余纲集}, 当且仅当 \(A^c\) 是\hyperref[dfn:meager]{第一纲集}.
        \end{dfn}

        \begin{dfn}
            [second-category]
            {Baire 第二纲集}
            [Baire 2nd Category Set]
            设 \(X\) 是一个拓扑空间. 定义 \(X\) 中的一个子集 \(A\) 是一个\textbf{Baire 第二纲集}, 当且仅当 \(A\) 不是\hyperref[dfn:meager]{第一纲集}.
        \end{dfn}

        \begin{rmk}
            [Dedicatia]
            需要注意的是, \colora{第二纲集和余纲集是不同的概念}. 存在既不是\hyperref[dfn:meager]{第一纲集}也不是\hyperref[dfn:residual]{余纲集}的集合. 但一个集合一定要么是第一纲集, 要么是第二纲集.\\
            由于 Baire 纲集是对任意拓扑空间都可以定义的, 所以在某些性质奇怪的空间中, 也有可能存在不是第二纲集的余纲集. 所以这是两个完全无关的概念, 请务必严格区分. 
        \end{rmk}

        \begin{ins}
            {\(\Q\) 是 \(\R\) 中的第一纲集}
            这是因为有理数是可数的, 而单点集的闭包是自己, 内部是空集.
        \end{ins}

        \begin{ins}
            {Cantor 集是 \([0,1]\) 中的第一纲集}
            因为它本身就是无处稠密的.
        \end{ins}

        \begin{ppt}
            {第一纲集的任意子集仍然是第一纲集}
        \end{ppt}

        \begin{ppt}
            {第一纲集的可数并仍然是第一纲集}
            其对偶结论为: \hyperref[dfn:residual]{余纲集}的可数交仍然是余纲集.
        \end{ppt}

        \begin{ppt}
            {可数个稠密开集的交是余纲集}
            这是第一纲集的定义的对偶结论.
        \end{ppt}

        \begin{rmk}
            由于
            \[X\setminus\interior A=\overline{X\setminus A}\]
            \[X\setminus\overline{B}=\interior(X\setminus B) \]
            仅证明第一个:\\
            设 \(x\in X\setminus\interior A\), 那么 \(x\) 的任意半径为 \(\epsilon\) 的开邻域都与 \(A\) 有交集. 从而 \(x\) 是 \(X\setminus A\) 的一个极限点, 也就是说 \(x\in\overline{X\setminus A}\). \\
            另一方面, 设 \(x\in \overline{X\setminus A}\), 存在一列 \(x_n\in X\setminus A\) 收敛到 \(x\). 但 \(x\) 的任意半径为 \(\epsilon\) 的开邻域都包含 \(x_n\) 中的某个元素, 从而 \(x\) 不可能是 \(\interior A\) 中的元素. 因此 \(x\in X\setminus\interior A\).
        \end{rmk}

        \begin{rmk}
             设 \(\{O_n\}_{n=1}^{\infty}\) 是 \(X\) 中的一列稠密开集, 那么
            \[\bigcap_{n=1}^{\infty} (X\setminus O_n)\]
            是 \(X\) 中的第一纲集. 考虑 \(X\setminus O_n\) 是闭集, 因为 \(O_n\) 是开集. 由于 \(O_n\) 是稠密的, 所以 \(\overline{O_n}=X\). 从而
            \[\interior(\overline{X\setminus O_n})=\interior(X\setminus O_n)=X\setminus\overline{O_n}=\varnothing\]
            这说明 \(X\setminus O_n\) 是 \(X\) 中的一个无处稠密集. 从而 \(\bigcap_{n=1}^{\infty} (X\setminus O_n)\) 是 \(X\) 中的第一纲集.
        \end{rmk}

        \begin{dfn}
            [baire-space]
            {Baire 空间}
            [Baire Space]
            设 \(X\) 是一个拓扑空间. 定义 \(X\) 是一个\textbf{Baire 空间}, 当且仅当 \(X\) 中任意可数个稠密开集的交集也是稠密的.
        \end{dfn}

        \begin{thm}
            {完备度量空间一个是 Baire 空间}
            设 \(X\) 是一个完备度量空间, 则 \(X\) 是一个 Baire 空间.
        \end{thm}

        \begin{thm}
            [bct]
            {Baire 纲定理}
            [Baire Category Theorem]
            Baire 空间中的开集一定是\hyperref[dfn:second-category]{第二纲的}.
        \end{thm}

        \begin{prf}
            设 \(\{O_n\}_{n=1}^{\infty}\) 是 \(X\) 中的一列稠密开集, 则我们要证明 \(\bigcap_{n=1}^{\infty} O_n\) 是稠密的, 等价于它与任意非空开集都有交集. 设 \(B:=\{y:|y-x|<r\}\) 是一个非空开球, 则我们要证明 \(B\cap \bigcap_{n=1}^{\infty} O_n\neq \varnothing\). \\
            从 \(O_1\) 开始, 由于 \(O_1\) 是稠密的, \(B(x,r)\) 与 \(O_1\) 有交集. 从而存在 \(x_1\in O_1\cap B\), \(r_1<\frac{r}{2}\) 使得 \(B(x_1,r_1)\subseteq O_1\cap B\). 从 \(O_2\) 开始, 由于 \(O_2\) 是稠密的, \(B(x_1,r_1)\) 与 \(O_2\) 有交集. 从而存在 \(x_2\in O_2\cap B(x_1,r_1)\), \(r_2<\frac{r}{4}\) 使得 \(B(x_2,r_2)\subseteq O_2\cap B(x_1,r_1)\). 重复这个过程, 可以得到一列点 \(\{x_n\}_{n=1}^{\infty}\) 和一列半径 \(\{r_n\}_{n=1}^{\infty}\), 使得
            \[|x_n-x_{n+p}|<r_p\]
            这说明 \(\{x_n\}_{n=1}^{\infty}\) 是 \(X\) 中的一个 Cauchy 列. 由于 \(X\) 是完备的, \(\{x_n\}_{n=1}^{\infty}\) 收敛到 \(X\) 中的一个点 \(x_0\). 
            \[x_0\in \bigcap_{n=1}^{\infty} O_n\]
            因此 \(B\cap \bigcap_{n=1}^{\infty} O_n\neq \varnothing\). 证毕.
            \tcbline
            该性质的对偶结论为: Baire 空间中可数个无处稠密的闭集的并有空的内部. 如果某个 Baire 空间是第一纲的, 那么它就可以表示为可数个无处稠密的闭集的并, 从而它的内部是空的. 这是不可能发生的.
        \end{prf}

        \begin{ppt}
            {余纲集的等价定义}
            Baire 空间中的 \(X\) 中的集合 \(A\) 是\hyperref[dfn:residual]{余纲集}当且仅当存在 G-delta 集 \(G\) 使得 \(G\subseteq A\), \(\overline{G}=A\).
        \end{ppt}

        \begin{rmk}
            [Dedicatia]
            \textbf{F-sigma 集是可数闭集的并集, G-delta 集是可数开集的交集}.\\
            定义一个集合 \(A\) 是一个 F-sigma 集当且仅当 \(A\) 可以表示为可数个闭集的并集. 定义一个集合 \(B\) 是一个 G-delta 集当且仅当 \(B\) 可以表示为可数个开集的交集.\\
            理解符号的含义有助于记忆:
            \begin{itemize}
                \item F-sigma 集中的 F 代表法语单词 ``Fermé'', 意为 ``闭集'', sigma 代表法语单词 ``Somme'', 意为``并'', 因此 F-sigma 集是可数闭集的并集.
                \item G-delta 集中的 G 代表德语单词 ``Gebiet'', 意为 ``区域(开集)'', delta 代表德语单词 ``Durchschnitt'', 意为``交'', 因此 G-delta 集是可数开集的交集.
            \end{itemize}
        \end{rmk}

        \begin{prf}
            "(\(\impliedby\))": 设 \(E\supseteq G:=\bigcap_{n=1}^{\infty} O_n\). 则 \(O_n\) 必定是一列稠密开集, 否则 \(G\) 就不是稠密的了. 从而根据上面的性质, \(E\) 是余纲集. \\
            "(\(\implies\))": 设 \(E\) 是余纲集. 则 \(X\setminus E\) 是第一纲集. 即
            \[X\setminus E=\bigcup_{n=1}^{\infty} A_n, \interior(\overline{A_n})=\varnothing.\]
            那么
            \[E=X\setminus\qty(\bigcup_{n=1}^{\infty}A_n)=\bigcap_{i=1}^{\infty}(X\setminus A_n)\supseteq\bigcap_{n=1}^{\infty}(X\setminus\overline{A_n} )=:G \]
            这就证明了 \(E\) 包含了稠密的 G-delta 集 \(G\).
        \end{prf}

        \begin{ppt}
            {Baire 空间中的余纲集是第二纲集}
            显然它不是第一纲集.
        \end{ppt}

        \begin{ins}
            {连续函数空间中处处不可微的函数是一个余纲集}
            装备了上确界范数的连续函数空间 \(C[0,1], \norm{\cdot}_{\infty}\) 中的处处不可微的函数集合是余纲集.\\
            我们考虑至少在某一点可微的函数集 \(Y\):
            \[Y\subset\bigcup_{n=1}^{\infty} \bigcup_{m=1}^{\infty} \qty{f:\exists\,x\in[0,1], 0\leq |x-y|<\frac{1}{n}, |f(x)-f(y)|<m|x-y| }=:\mathcal{D} \]
            事实上, 只要给 \(f\) 加上高频震荡波, 就可以使得 \(f\) 范数不变但导数变化很剧烈. 从而 \(\mathcal{D}\) 是第一纲集.
        \end{ins}

        \begin{ins}
            {存在零测的余纲集}
            将 \([0,1]\) 的有理数枚举为 \(q_i\) 并定义区间
            \[I_{i,j}:=\qty(q_i-\frac{1}{i+j},q_i+\frac{1}{i+j})\]
            定义
            \[G_j:=\bigcup_{i=1}^{\infty}I_{i,j}\]
            \[B:=\bigcap_{j=1}^{\infty}G_j\]
            则 \(G\) 有以下性质:
            \begin{enumerate}[label=(\arabic*)]
                \item \(G\) 是 G-delta 集, 因为 \(G\) 是可数个开集的交集;
                \item \(G\) 是余纲集, 因为它是可数的稠密开集 \(G_j\) 的交集;
                \item \(G\) 的外测度为零.
            \end{enumerate}
        \end{ins}

        \begin{rmk}
            [Dedicatia]
            这说明测度意义下的 ``大'' 和 Baire 纲意义下的 ``大'' 是不同的概念. 例如上面的集合 \(G\) 在测度意义下是一个零测集, 但在 Baire 纲意义下却是一个\hyperref[dfn:residual]{余纲集}.\\
            四类集合都有例子:
            \begin{center}
                \begin{tabular}{c|cc}
                    & 零测集 & 正测度的 Lebesgue 可测集 \\
                    \hline
                    第一纲集 & \(\Q\) & \([0,1]\setminus G\) \\
                    余纲集 & \(G\) & \([0,1]\)
                \end{tabular}
            \end{center}
        \end{rmk}

        \begin{ins}
            {不存在仅在有理点连续的函数}
            设 \(f:[0,1]\to\R\). 则我们知道存在仅在无理点连续的函数 \(f\), 例如 Dirichlet 函数. 但是不存在仅在有理点连续的函数. \\
            更一般地, 设度量空间 $X,Y$, 函数 $f:X \to Y$, 则 $f$ 的连续点集是一个 G-delta 集. \\
            因此, 仅在有理数点连续的函数不存在, 因为 $\Q$ 不是 G-delta 集. 但仅在无理数点连续的函数存在, 因为 $\R \setminus \Q$ 是 G-delta 集.
        \end{ins}

        \begin{ins}
            {无穷可微函数的性质}
            设 \(f\) 是无穷阶可微的实函数, 并且 \(\forall\,n\in\N, \exists\,x\in\R\) s.t. \(f^{(n)}(x)=0\), 则 \(f\) 是一个多项式函数.
        \end{ins}

    \subsection{线性算子的性质定理}

        \begin{dfn}
            [open-mapping]
            {开映射}
            [Open Mapping]
            设 \(X, Y\) 是两个拓扑空间, \(f:X\to Y\) 是一个映射. 定义 \(f\) 是一个\textbf{开映射 (Open Mapping)}, 当且仅当 \(f\) 将 \(X\) 中的任意开集映到 \(Y\) 中的开集.
        \end{dfn}

        \begin{rmk}
            开映射与连续映射恰好是相对的概念: 开映射将开集推前, 连续映射将开集拉回. \\
            一个同时是开映射、连续映射、双射的映射是\textbf{同胚映射 (Homeomorphism)}.
        \end{rmk}

        \begin{thm}
            [omt]
            {Banach 空间的开映射定理}
            [Open Mapping Theorem for Banach Spaces]
            设 \(X,Y\) 是两个 Banach 空间, \(T\in\blo(X,Y)\), 如果 \(T\) 是满射, 则 \(T\) 是一个\hyperref[dfn:open-mapping]{开映射}.
        \end{thm}

        \begin{prf}
            由于 \(X\) 可以分成一列开球的并
            \[X=\bigcup_{n=1}^{\infty} B_X(0,n)\]
            而 \(T\) 是满射
            \[Y=T(X)=T\qty(\bigcup_{n=1}^{\infty} B_X(0,n))=\bigcup_{n=1}^{\infty} T(B_X(0,n))\]
            根据 \(T\) 的线性性, 
            \[Y=\bigcup_{n=1}^{\infty} T(B_X(0,n))=\bigcup_{n=1}^{\infty} nT(B_X(0,1))\]
            由于 \(Y\) 是完备的度量空间, 所以是 Baire 空间, 满足 Baire 纲定理, 它一定是第二纲的. 这说明上面的一列集合中至少有一个集合的闭包是有内部的.\\
            也就是 \(\overline{T(B_X(0,1))}\) 必然包含开球. 设为 \(B_Y(y_0,r)\subseteq\overline{T(B_X(0,1))} \). 则
            \[y=y_0+y-y_0\in\overline{T(B_X(0,1))}-\overline{T(B_X(0,1))}\subseteq\overline{T(B_X(0,2))}  \]
            根据线性性可以把 \(2\) 用标量乘法换成任意 \(\epsilon>0\). 所以我们得到了
            \[\exists\,\eta>0, \st B_Y(0,\epsilon\eta)\subseteq\overline{T(B_X(0,\epsilon))} \]
            特别地可以取 \(\epsilon=1\):
            \[\exists\,\eta>0, \st B_Y(0,\eta)\subseteq\overline{T(B_X(0,1))} \]
            接下来应该想办法把闭包去掉, 考虑证明 \(B_Y(0,\frac{\eta}{2})\subseteq T(B_X(0,1))\). 根据上面的结论, 令 \(\epsilon=\frac{1}{2}\), 得到
            \[\exists\,\eta>0, \st B_Y(0,\frac{\eta}{2})\subseteq\overline{T(B_X(0,\frac{1}{2}))} \]
            则存在 \(x_1\in B_X(0,\frac{1}{2})\) 使得 \(\norm{T(x_1)-y}<\frac{\eta}{4}\). \\
            设 \(y_1=y-T(x_1)\), 再令 \(\epsilon=\frac{1}{4}\), 得到
            \[\exists\,\eta>0, \st B_Y(0,\frac{\eta}{4})\subseteq\overline{T(B_X(0,\frac{1}{4}))} \]
            同样的, 存在 \(x_2\in B_X(0,\frac{1}{4})\) 使得 \(\norm{T(x_2)-y_1}<\frac{\eta}{8}\). \\
            重复这个过程, 可以得到一列点 \(\{x_n\}_{n=1}^{\infty}\) 和一列数 \(\{\frac{\eta}{2^n}\}_{n=1}^{\infty}\). 按照构造过程, 满足
            \[\norm{x_n}\leq\frac{1}{2^n}\]
            这说明 \(\{x_n\}\) 是 Cauchy 列, 根据完备性, 收敛到 \(X\) 中的一个点 \(x\):
            \[\norm{x}\leq\sum_{n=1}^{\infty} \norm{x_n} \leq \sum_{n=1}^{\infty} \frac{1}{2^n} = 1\implies x\in B_X(0,1) \]
            另一方面, 由于
            \[\norm{y-T(x_1+x_2+\cdots+x_n)}\leq\frac{\eta}{2^{n+1}}\to 0\]
            且 \(T\) 是连续线性算子, 则可以交换极限和 \(T\) 的作用,
            \[\lim_{n\to\infty}T\qty(\sum_{k=1}^n x_k)=T\qty(\lim_{n\to\infty}\sum_{k=1}^n x_k)=Tx\]
            所以 \(T(x)=y\). 这说明对任意 \(B_Y(0,\frac{\eta}{2})\) 的元素都有来自 \(T(B_X(0,1))\) 的原像. 从而 \(T(B_X(0,1))\) 确实包含了 \(B_Y(0,\frac{\eta}{2})\). 这就说明了 \(T\) 是一个开映射.
        \end{prf}

        \begin{rmk}
            [Dedicatia]
            这个证明中 Banach 空间的要求是不可缺的: \\
            (i) \(Y\) 是 Banach 空间是为了保证 \(Y\) 是 Baire 空间;
            (ii) \(X\) 是 Banach 空间是为了保证最后构造的级数一定收敛.
        \end{rmk}

        \begin{thm}
            [imt]
            {Banach 空间的逆映射定理}
            [Inverse Mapping Theorem for Banach Spaces] 
            设 \(X,Y\) 是两个 Banach 空间, \(T\in\blo(X,Y)\), 如果 \(T\) 是双射, 则 \(T^{-1}\in\blo(Y,X)\).
        \end{thm}

        \begin{prf}
            当 \(T\) 是双射时, \(T\) 是开映射等价于 \(T^{-1}\) 是连续映射. 根据开映射定理即可得到.
        \end{prf}

        \begin{thm}
            {Banach 空间的范数等价定理}
            [Norm Equivalence Theorem for Banach Spaces]
            设 \(X\) 是一个线性空间, \(\norm{\cdot}_1, \norm{\cdot}_2\) 是 \(X\) 上的两个范数, 使得 \((X,\norm{\cdot}_1)\) 和 \((X,\norm{\cdot}_2)\) 都是 Banach 空间, 则 \(\norm{\cdot}_1\) 和 \(\norm{\cdot}_2\) 是等价的.
        \end{thm}

        \begin{prf}
            对两个范数诱导的 Banach 空间使用逆映射定理.
        \end{prf}

        \begin{dfn}
            {闭线性算子}
            [Closed Linear Operator]
            设 \(X,Y\) 是两个 Banach 空间, \(T:X\to Y\) 是一个线性映射. 定义 \(T\) 是一个\textbf{闭线性算子}, 当且仅当 \(T\) 的图像 \(\Gamma(T)=\{(x,T(x)):x\in X\}\) 是积拓扑空间 \(X\times Y\) 中的一个闭集. 这个闭集称为\textbf{闭图像(Closed Graph)}.
        \end{dfn}

        \begin{thm}
            [cgt]
            {Banach 空间的闭图像定理}
            [Closed Graph Theorem for Banach Spaces]
            设 \(X,Y\) 是两个 Banach 空间, \(T:E\subseteq X\to Y\) 是一个线性映射, \(E\) 是闭子空间. 如果 \(T\) 是闭线性算子, 则 \(T\in\blo(X,Y)\).
        \end{thm}

        \begin{prf}
            设 \(X\) 上的范数为 \(\norm{\cdot}_X\), \(Y\) 上的范数为 \(\norm{\cdot}_Y\). 定义 \(X\times Y\) 上的范数为
            \[\norm{(x,y)}_{X\times Y}=\norm{x}_X+\norm{y}_Y.\]
            则 \(X\times Y\) 是一个 Banach 空间. 由于 \(T\) 是闭线性算子, \(\Gamma(T)\) 是 \(X\times Y\) 中的一个闭子空间. 从而 \(\Gamma(T)\) 是一个 Banach 空间. 定义映射 \(S:\Gamma(T)\to X\), \((x,T(x))\mapsto x\). 则 \(S\) 是一个双射, 根据逆映射定理, \(S^{-1}\in\blo(X,\Gamma(T))\). 从而 \(T=\pi_Y\circ S^{-1}\in\blo(X,Y)\).
        \end{prf}

        \begin{cxmp}
            {微分算子是一个闭线性算子}
            在配备上确界范数的 \(\Cont^1[0,1]\) 函数空间上定义的微分算子 \(D(f)=f'\) 是闭线性算子. 也就是说, 如果函数列 \(\{f_n\}\) 一致收敛到 \(f\), 且 \(D(f_n)\) 一致收敛到 \(g\), 则 \(g=D(f)\). \\
            但是我们之前已经说过微分算子不是有界线性算子. 这不与闭图像定理矛盾, 因为此时的 \(\Cont^1[0,1]\) 不是一个 Banach 空间.\\
            \(\Cont^1[0,1]\) 在配备 C1 范数下才是完备的. C1 范数定义为
            \[\norm{f}_{C1}=\norm{f}_{\infty}+\norm{f'}_{\infty}\]
            或者, 需要固定 \(x_0\in[0,1]\) 使得 \[\norm{f}_{C1}=\norm{f(x_0)}+\norm{f'}_{\infty}\]
            这时微分算子 \(D\) 就是一个有界线性算子, 自然也满足了闭图像定理. 
        \end{cxmp}

        \begin{thm}
            [banach-steinhaus]
            {一致有界性定理 / Banach-Steinhaus 定理 / 共鸣定理}
            [Uniform Boundedness Principle / Banach-Steinhaus Theorem]
            设 \(X\) 是 Banach 空间, \(Y\) 是赋范线性空间, 一族有界线性算子满足
            \[T_i\in\blo(X, Y), \qquad \forall i\in I,\forall x\in X, \sup_{i\in I} \norm{T_i x} < \infty.\]
            则对任意 \(x\in X\), 有
            \[\sup_{i\in I} \norm{T_i(x)} < \infty.\]
        \end{thm}

        \begin{prf}
            在 \(X\) 上定义范数 
            \[\norm{x}_M:=\max\{\norm{x}, \sup_{i\in I} \norm{T_i x}\}\]
            定义映射
            \[J: (X, \norm{\cdot}_M) \to (X, \norm{\cdot}_X)\]
            则 \(J\) 是双射. 根据逆映射定理, \(J^{-1}\in\blo(X, \norm{\cdot}_M)\). 即
            \[\norm{x}_M = \norm{J^{-1}(x)}_M \leq \norm{x}_X\norm{J^{-1}}.\]
            这说明了 \(\sup_{i\in I} \norm{T_i x} \leq \norm{x}_X\norm{J^{-1}}\). 证毕.
        \end{prf}

        \begin{rmk}
            [Dedicatia]
            之所以叫一致有界性定理, 是因为它能根据逐点的有界推出算子范数的有界, 并且所有的算子都是等价地有界.\\
            其背后的核心原因就是 Banach 空间的完备性: 根据 Baire 空间的性质, 如果有一个坏点被有界线性算子拉到了无穷大, 那么全空间中应该到处都是. 这显然违背了逐点的有界性.
            \tcbline
            这个定理的一个应用是证明了存在连续函数, 它的 Fourier 级数在某些点不收敛. 这是因为 Fourier 级数的部分和算子是一个有界线性算子, 但它们的范数却不受限制地增长. 自然也找不到一致的上界.
        \end{rmk}

\section{弱拓扑}

    \subsection{【点集拓扑】度量空间的性质}

        \begin{intrormk}
            [Dedicatia]
            本节先介绍拓扑学中的几个一般概念及其在一般的\hyperref[dfn:ms]{度量空间}中的体现.
        \end{intrormk}

        \begin{vardfn}
            {开覆盖}
            [Open Cover]
            设 \(X\) 是一个拓扑空间, \(A\subseteq X\). 定义 \(A\) 的一个\textbf{开覆盖}是 \(X\) 中的一族开集 \(\{O_\alpha\}_{\alpha\in I}\), 使得 
            \[A\subseteq\bigcup_{\alpha\in I} O_\alpha.\]
            这时也称 \(\{O_\alpha\}_{\alpha\in I}\) 覆盖住了 \(A\).
        \end{vardfn}

        \begin{vardfn}
            [compact]
            {紧致集 / 紧集}
            [Compact Set]
            设 \(X\) 是一个拓扑空间, \(A\subseteq X\). 定义 \(A\) 是一个\textbf{紧致集}, 当且仅当 \(A\) 的任意开覆盖都存在一个有限子覆盖. \\
            也就是说, 对任意 \(A\) 的开覆盖 \(\{O_\alpha\}_{\alpha\in I}\), 都存在 \(\alpha_1,\dots,\alpha_n\in I\) 使得
            \[A\subseteq\bigcup_{i=1}^n O_{\alpha_i}.\]
        \end{vardfn}

        \begin{vardfn}
            {有限交性}
            [Finite Intersection Property]
            设 \(X\) 是拓扑空间, \(\{A_\alpha\}_{\alpha\in I}\) 是 \(X\) 中的一族集合. 定义 \(\{A_\alpha\}_{\alpha\in I}\) 满足有限交性当且仅当
            \[\bigcap_{i=1}^n A_{\alpha_i} \neq \varnothing.\]
            也就是说, 任意有限个集合的交集都非空.
        \end{vardfn}

        \begin{varthm}
            [compact]
            {紧集的有限交性定义}
            \(A\) 是 \(X\) 中的\hyperref[dfn:compact]{紧集}当且仅当对 \(X\) 中的任意一族闭集 \(\{F_\alpha\}_{\alpha\in I}\) 都满足:
            若 \(\{F_\alpha\cap A\}_{\alpha\in I}\) 满足有限交性, 则 \(\bigcap_{\alpha\in I} F_\alpha\cap A\neq\varnothing.\)
        \end{varthm}

        \begin{rmk}
            [Dedicatia]
            这个等价定义实际上是使用开集定义的紧集的对偶版本. 只需利用闭集是开集的补、以及 De Morgan 定律就可以证明了.
        \end{rmk}

        \begin{varthm}
            [closed-subset-of-compact]
            {紧集的闭子集是紧的}
            [Closed Subset of Compact Set is Compact]
            设 \(A\) 是 \(X\) 中的一个\hyperref[dfn:compact]{紧集}, 则 \(A\) 的任意闭子集 \(B\) 也是 \(X\) 中的一个紧集.
        \end{varthm}

        \begin{prf}
            设 \(B\) 是紧集 \(A\) 的闭子集. 则 \(X\) 中的一族闭集 \(\{F_\alpha\}_{\alpha\in I}\) 满足 \(\{F_\alpha\cap B\}_{\alpha\in I} = \{(F_\alpha\cap B)\cap A\}_{\alpha\in I}\) 满足有限交性. 利用 \(A\) 的紧性, 
            \[\bigcap_{\alpha\in I} (F_\alpha\cap B) = \bigcap_{\alpha\in I} ((F_\alpha\cap B)\cap A) \neq \varnothing.\]
            这就说明了 \(B\) 是紧集.
        \end{prf}

        \begin{varthm}
            {连续映射保持紧致性}
            设 \(f:X\to Y\) 是拓扑空间 \(X\) 到 \(Y\) 的一个连续映射, \(A\subseteq X\) 是 \(X\) 中的一个\hyperref[dfn:compact]{紧集}, 则 \(f(A)\) 是 \(Y\) 中的一个紧集.
        \end{varthm}

        \begin{prf}
            设 \(\{O_\alpha\}_{\alpha\in I}\) 是 \(Y\) 中的一族开集, 使得 \(f(A)\subseteq\bigcup_{\alpha\in I} O_\alpha\). 则 \(\{f^{-1}(O_\alpha)\}_{\alpha\in I}\) 是 \(X\) 中的一族开集, 使得
            \[A\subseteq f^{-1}(f(A)) \subseteq \bigcup_{\alpha\in I} f^{-1}(O_\alpha).\]
            根据 \(A\) 的紧性, 存在 \(\alpha_1,\dots,\alpha_n\in I\) 使得
            \[A\subseteq \bigcup_{i=1}^n f^{-1}(O_{\alpha_i}).\]
            从而
            \[f(A)\subseteq \bigcup_{i=1}^n O_{\alpha_i}.\]
            这就说明了 \(f(A)\) 是 \(Y\) 中的一个紧集.
        \end{prf}

        \begin{varthm}
            [tychonoff]
            {Tychonoff 定理}
            设 \(X_\alpha\) 是一族\hyperref[dfn:compact]{紧拓扑空间}, \(I\) 是任意一个索引集. 则积拓扑空间
            \[\prod_{\alpha\in I} X_\alpha\]
            是一个紧拓扑空间.
        \end{varthm}

        \begin{prf}
            证明该定理需要承认 \hyperref[thm:zorn]{Zorn 引理}. 设 \(\mathscr{F}\) 是积拓扑空间 \(\prod\limits_{\alpha\in I} X_\alpha\) 的一族满足有限交性质的闭集. 根据 Zorn 引理, 设 \(\mathscr{P}\) 是任意包含了 \(\mathscr{F}\) 的具有有限交性的闭子集族, 其中偏序关系是 \(\subseteq\), 对于 \(\mathscr{P}\) 的任意全序子集 \(\mathscr{C}\), 都存在一个上界 \(\bigcup\mathscr{C}\). 从而存在一个包含 \(\mathscr{F}\) 的最大的满足有限交性的闭集族 \(\mathscr{E}\).\\
            令 \(p_\alpha\) 是 积拓扑空间 \(\prod\limits_{\alpha\in I} X_\alpha\) 上的第 \(\alpha\) 个坐标投影. 对任意 \(E\in\mathscr{E}\), 记 \(E_\alpha = p_\alpha(E)\). 则 \(\{E_\alpha:E\in\mathscr{E}\}\) 也满足有限交性, 也就是
            \[\bigcap_{i=1}^n E_{\alpha_i} = \bigcap_{i=1}^n p_{\alpha_i}(E) = p_{\alpha_1}\qty(\bigcap_{i=1}^n E_i)\neq\varnothing.\]
            \bysorry
        \end{prf}

        \begin{vardfn}
            [sequential-compact]
            {列紧性 / 序列紧性}
            [Sequential Compactness]
            设 \(X\) 是\hyperref[dfn:ms]{度量空间}, 定义 \(X\) 是列紧的, 当且仅当
            \[\forall\,\{x_n\}\subset X, \exists\,\{x_{n_k}\}\subseteq\{x_n\} \st \lim_{k \to \infty} x_{n_k} \in X\]
            也就是, \(X\) 中的任意序列都存在一个收敛子序列. \\
            对子集也可以定义列紧性. 设 \(A\subseteq X\), 定义 \(A\) 是列紧的, 当且仅当 \(A\) 中的任意序列都存在一个收敛子序列, 且收敛子序列的极限也在 \(A\) 中.
        \end{vardfn}

        \begin{rmk}
            需要注意的是, 列紧性的定义就是其中的任意序列都有收敛的子列. 而 \(\R\) 上的 Bolzano-Weierstrass 定理则是其中的任意有界序列都有收敛的子列. 这两者是不同的概念. \\
            但换句话表述就是, \(\R\) 中的任意有界闭集都是列紧的. 结合 \(\R\) 是度量空间的事实, 这正是 Heine-Borel 定理的内容.
        \end{rmk}

        \begin{thm}
            {度量空间中的紧性与列紧性是等价的}
            设 \(X\) 是一个\hyperref[dfn:ms]{度量空间}, 则 \(X\) 是\hyperref[dfn:compact]{紧集}当且仅当 \(X\) 是\hyperref[dfn:sequential-compact]{列紧的}.
        \end{thm}

        \begin{vardfn}
            {邻域基}
            [Neighborhood Base]
            设 \(X\) 是一个拓扑空间, \(x\in X\). 定义 \(X\) 中的一个集合族 \(\B_x\) 是 \(x\) 的一个\textbf{邻域基}, 当且仅当:
            \(\B_x\) 是 \(x\) 的所有邻域的子集, 且对于 \(x\) 的任意一个邻域 \(V\), 都存在 \(B\in\B_x\) 使得 \(B\subseteq V\).\\
            \(x\) 的邻域基记作 \(\B_x\).
        \end{vardfn}

        \begin{vardfn}
            [first-countable]
            {第一可数性质}
            [First Countability]
            设 \(X\) 是一个拓扑空间. 定义 \(X\) 满足\textbf{第一可数性质}, 当且仅当
            \[\forall\,x\in X, \exists\,\B_x \st \card \B_x \leq \aleph_0\]
            也就是说, \(X\) 中的任意一个点都存在可数的邻域基. \\
            满足第一可数性质的空间也称为\textbf{第一可数空间}.
        \end{vardfn}

        \begin{thm}
            {度量空间是第一可数空间}
            \(X\) 是\hyperref[dfn:ms]{度量空间}, 则 \(X\) 满足\hyperref[dfn:first-countable]{第一可数性质}.
        \end{thm}

        \begin{prf}
            \(\forall x\in X\), 开球族 \(\qty{B\qty(x,\frac{1}{n}):n\in\N}\) 是 \(x\) 的一个可数邻域基.
        \end{prf}

        \begin{vardfn}
            [second-countable]
            {第二可数性质}
            [Second Countability]
            设 \(X\) 是拓扑空间, 定义 \(X\) 满足\textbf{第二可数性质}, 当且仅当 \(X\) 有可数的拓扑基. \\
            满足第二可数性质的空间也称为\textbf{第二可数空间}.
        \end{vardfn}

        \begin{varthm}
            {第二可数蕴含第一可数}
            显然.
        \end{varthm}

        \begin{varthm}
            {可分的度量空间是第二可数空间}
            设 \(X\) 是一个\hyperref[dfn:ms]{度量空间}, 如果 \(X\) 是\hyperref[dfn:separable]{可分的}, 则 \(X\) 满足\hyperref[dfn:second-countable]{第二可数性质}.
        \end{varthm}

        \begin{prf}
            设 \(D\) 是 \(X\) 中的一个可数稠密子集, 则 \(\B=\{B(x,r):x\in D, r\in\Q^+\}\) 是 \(X\) 的一个可数拓扑基.
        \end{prf}

        \begin{rmk}
            [Dedicatia]
            事实上, 第二可数性质比\hyperref[dfn:separable]{可分性}强. 但是在度量空间中二者等价.
        \end{rmk}

    \subsection{弱拓扑与弱*拓扑}

        \begin{intrormk}
            [Dedicatia]
            本节介绍弱拓扑和弱*拓扑. 这是两种在线性空间上定义的比范数拓扑更弱的拓扑. 这两种拓扑在泛函分析中非常重要, 因为它们允许我们在 Banach 空间中讨论收敛性和连续性时使用更宽松的条件.
        \end{intrormk}

        \begin{dfn}
            {弱邻域}
            [Weak Neighborhood]
            设 \(X\) 是一个赋范线性空间, \(x\in X\). 定义 \(X\) 中的一个集合 \(V\) 是 \(x\) 的一个\textbf{弱邻域}, 当且仅当存在 \(f_1,\dots,f_n\in X^*\), \(\epsilon>0\) 使得
            \[V\supseteq\bigcap_{i=1}^n \{y\in X: |f_i(y-x)|<\epsilon\}.\]
        \end{dfn}

        \begin{rmk}
            需要注意的是, 弱邻域的定义中要求泛函的数量是有限个. \\
            这带来的影响就是弱邻域非常的大, 只限制了有限个方向, 剩余的无数个泛函的值都没有限制. 进而, 我们将会看到它是不可度量的.
        \end{rmk}

        \begin{dfn}
            [weak-topology]
            {弱拓扑}
            [Weak Topology]
            设 \(X\) 是一个赋范线性空间. 定义 \(X\) 上的\textbf{弱拓扑}是由 \(X\) 中的所有弱邻域构成的拓扑. 记作 $(X,w_X)$.
        \end{dfn}

        \begin{thm}
            {弱拓扑的等价定义}
            设 \(X\) 是赋范线性空间, \(X\) 上的弱拓扑可以定义为:
            \begin{enumerate}[label=(\alph*)]
                \item \(f_i\) 是一族能分离点的线性泛函, 弱拓扑是使得所有 \(f_i\) 都连续的最弱的拓扑;
                \item 由半范数族 \[\{p_f\}_{f\in X^*}, p_f (x) = |f(x)|\] 生成的拓扑: 弱邻域为 \[\{x\in X: p_{f_i}(x-x_0) < \epsilon,\forall i = 1,\cdots,n\}\]
            \end{enumerate}
        \end{thm}

        \begin{prf}
            这里我们只介绍思路: \\
            弱拓扑是使得所有线性泛函都连续的最弱的拓扑, 根据拓扑空间连续映射的定义, 泛函映入的标量域的开集的原像必须也是开的. 而根据拓扑空间的定义, 有限多个开集的交是开集, 所以弱邻域的定义就自然地是\textbf{有限多个连续线性泛函的原像的交集}:
            \[\bigcap_{i=1}^n \{y\in X: |f_i(y-x)|<\epsilon\}\]
            接下来只需验证这是一族半范数:\begin{enumerate}[label=(\roman*)]
                \item 非负性: \(p_f(x)=|f(x)|\geq 0\);
                \item 绝对齐次性: \(p_f(\lambda x)=|\lambda||f(x)|=|\lambda|p_f(x)\);
                \item 次可加性: \(p_f(x+y)=|f(x+y)|\leq |f(x)|+|f(y)|=p_f(x)+p_f(y)\).
            \end{enumerate}
            所以这不过是等价转述了弱邻域的定义而已.
        \end{prf}

        \begin{ppt}
            {弱拓扑空间是 Hausdorff 空间}
        \end{ppt}

        \begin{prf}
            设 \(x,y\in X, x\neq y\). 根据 Hahn-Banach 定理, \(\exists\,f \in X^*,\alpha\in\R, f(x)>\alpha>f(y)\). 从而 \(f\) 是一个连续线性泛函, \(f^{-1}((\alpha,+\infty))\) 和 \(f^{-1}((-\infty,\alpha))\) 是 \(X\) 中的两个弱开集, 分别包含 \(x\) 和 \(y\), 从而 \(X\) 是 Hausdorff 空间.
        \end{prf}

        \begin{dfn}
            {弱*邻域}
            [Weak-star Neighborhood]
            设 \(X\) 是一个赋范线性空间, \(f\in X^*\). 定义 \(X^*\) 中的一个集合 \(V\) 是 \(f\) 的一个\textbf{弱*邻域}, 当且仅当存在 \(x_1,\dots,x_n\in X\), \(\epsilon>0\) 使得
            \[V\supseteq\bigcap_{i=1}^n \{g\in X^*: |g(x_i)-f(x_i)|<\epsilon\}.\]
        \end{dfn}

        \begin{dfn}
            {弱*拓扑}
            [Weak-star Topology]
            设 \(X\) 是一个赋范线性空间. 定义 \(X^*\) 上的\textbf{弱*拓扑}是由 \(X^*\) 中的所有弱*邻域构成的拓扑. 记作 \((X^*, w^*)\).
        \end{dfn}

        \begin{rmk}
            [Dedicatia]
            \(X^{**}\) 上的弱*拓扑不一定同胚于 \(X\) 上的弱拓扑. 只有当 \(X\) 是一个\hyperref[dfn:reflexive]{自反空间}时, \(X^{**}\) 上的弱*拓扑才同胚于 \(X\) 上的弱拓扑.
        \end{rmk}

        \begin{thm}
            {有限维空间的弱拓扑、弱*拓扑与范数拓扑一致}
            设 \(X\) 是赋范线性空间, \(\dim X<+\infty\). 则 \((X, \tau_{\norm{\cdot}})\cong (X, w_X) \cong (X^*,w^*)\)
        \end{thm}

        \begin{prf}
            我们仅证明对弱*拓扑的情况. 设 \(O\subseteq X^*\) 是弱*开集, \(f_0\in O\). 设 \(\{e_i\}_{i=1}^n\) 是 \(X\) 的一组基. 定义
            \[p(f) = \max_{1\leq i\leq n} |f(e_i)|\]
            则 \(p\) 是 \(X^*\) 上的范数, 且与原先任意一个范数等价. 它诱导的开集
            \[O'=\{f\in X^*: p(f-f_0)<\epsilon\} = \bigcap_{i=1}^n \{f\in X^*: |f(e_i)-f_0(e_i)|<\epsilon\}\]
            是一个弱*开集. 这就说明了 \(O'\subseteq O\). 另一方面, 设 \(O''\) 是 \(X^*\) 上的一个弱*开集, \(f_0\in O''\). 则存在 \(x_1,\dots,x_m\in X\), \(\epsilon>0\) 使得
            \[O''\supseteq\bigcap_{j=1}^m \{f\in X^*: |f(x_j)-f_0(x_j)|<\epsilon\}\]
            设 \(x_j=\sum_{i=1}^n a_{ij} e_i\). 则
            \[|f(x_j)-f_0(x_j)| = \left|\sum_{i=1}^n a_{ij} (f(e_i)-f_0(e_i))\right| \leq \sum_{i=1}^n |a_{ij}| |f(e_i)-f_0(e_i)|\]
            从而
            \[\bigcap_{j=1}^m \{f\in X^*: |f(x_j)-f_0(x_j)|<\epsilon\} \supseteq \bigcap_{i=1}^n \qty{f\in X^*: |f(e_i)-f_0(e_i)|<\frac{\epsilon}{\sum_{j=1}^m |a_{ij}|}} = O'''\]
            这就说明了 \(O'''\subseteq O''\). 综上, \(O'''\subseteq O''\subseteq O\). 这就说明了 \(X^*\) 上的任意一个弱*开集都包含了一个范数开集. 从而 \(X^*\) 上的弱*拓扑与范数拓扑一致.
        \end{prf}

        \begin{ppt}
            {弱*拓扑空间是 Hausdorff 空间}
            设 \(X\) 是一个赋范线性空间, 则 \((X^*, w^*)\) 是一个 Hausdorff 空间.
        \end{ppt}

        \begin{ppt}
            {双对偶空间中的弱*连续性}
            设 \(X\) 是 Banach 空间, \(F\in X^{**}\). 如果 \(F\) 是弱*连续的, 则 \(\exists\,x\in X, \st F=J_X(x)\). \(J_X\) 是\hyperref[dfn:canonical-embedding]{典范嵌入}.
        \end{ppt}

        \begin{prf}
            \(F\) 是弱*连续的, 也就是它映入的标量域的开集的原像是弱*开集. 设为 \(O\):
            \[O_\epsilon=\bigcap_{i=1}^{n}\{f\in X^*, |f(x_i)|<\epsilon \}\]
            换句话说, \(\forall\,f\in O_\epsilon\), \(|F(f)|<M\), 不妨取 \(M=1\). 考虑
            \[\mathcal{N}:=\bigcap_{i=1}^n \qty{f\in X^*, f(x_i)=0}=\bigcap_{i=1}^n\ker J_X(x_i) \]
            则 \(\mathcal{N}\) 是 \(X^*\) 中的一个线性子空间且使得: \(\forall \lambda\in\R, |F(\lambda f)|\leq 1\), 所以 \(F(f)=0\) 对任意 \(f\in\mathcal{N}\) 都成立. 所以
            \[F = \sum_{i=1}^n \alpha_i J_X(x_i).\]
        \end{prf}

        \begin{thm}
            {无限维赋范线性空间的弱拓扑是不可度量化的}
            [Weak Topology on Infinite-dimensional Normed Linear Spaces Is Not Metrizable]
            设 \(X\) 是一个赋范线性空间, \(\dim X = +\infty\). 则 \(X\) 上的弱拓扑不可度量化. 即不存在度量 \(d\) 使得 \((X, w_X)\) 是一个度量空间.
        \end{thm}

        \begin{prf}
            提示: 弱拓扑不满足\hyperref[dfn:first-countable]{第一可数性质}. 因为弱邻域很大, 依旧是不可数的.
        \end{prf}

        \begin{dfn}
            [weak-convergence]
            {弱收敛}
            [Weak Convergence]
            设 \(X\) 是一个赋范线性空间, \(\{x_n\}_{n=1}^{\infty}\) 是 \(X\) 中的一列元素, \(x\in X\). 定义 \(\{x_n\}_{n=1}^{\infty}\) \textbf{弱收敛}到 \(x\), 当且仅当 \(\forall\,f\in X^*, f(x_n)\to f(x)\).\\
            记作 \(x_n\wlim x\).
        \end{dfn}

        \begin{dfn}
            [weak-star-convergence]
            {弱*收敛}
            [Weak-star Convergence]
            设 \(X\) 是一个赋范线性空间, \(\{f_n\}_{n=1}^{\infty}\) 是 \(X^*\) 中的一列元素, \(f\in X^*\). 定义 \(\{f_n\}_{n=1}^{\infty}\) \textbf{弱*收敛}到 \(f\), 当且仅当 \(\forall\,x\in X, f_n(x)\to f(x)\).\\
            记作 \(f_n\wslim f\).
        \end{dfn}

        \begin{rmk}
            [Dedicatia]
            弱收敛/弱*收敛就是在弱拓扑/弱*拓扑下的收敛方式.
        \end{rmk}

        \begin{ins}
            {逐点收敛}
            [Pointwise Convergence]
            弱*收敛通俗地说也就是逐点收敛 / 点态收敛. \(f_n\wslim f\) 就是 \(\forall x \in X, f_n(x) \to f(x)\).\\
            我们在【数学分析】中就已经知道逐点收敛是很弱的:
            \begin{itemize}
                \item 逐点收敛的连续函数列极限不一定是连续函数;
                \item 逐点收敛的 Riemann 可积函数列极限不一定是 Riemann 可积函数;
                \item 逐点收敛的可微函数列的极限不一定是可微函数; 即使加上导数逐点收敛也不能保证;
                \item 逐点收敛的 Lebesgue 可积函数列的极限不一定是 Lebesgue 可积函数, 除非是根据 Egorov 定理, 逐点收敛的可测函数列在有限测度下差不多一致收敛.
            \end{itemize}
            相比之下, 一致收敛拓扑可以保持绝大多数良好性质:
            \begin{itemize}
                \item 一致收敛的连续函数列极限是连续函数;
                \item 一致收敛的 Riemann 可积函数列极限是 Riemann 可积函数;
                \item 一致收敛的可微函数列极限不一定是可微函数, 但是如果赋予 C1 范数, 即改为要求逐点收敛的可微函数列的导数也一致收敛, 则极限函数是可微函数;
                \item 一致收敛的 Lebesgue 可积函数列极限是 Lebesgue 可积函数.
            \end{itemize}
            这个例子可以看出范数拓扑与弱*拓扑的差别.
        \end{ins}

        \begin{thm}
            {任意可分的 Hilbert 空间的标准正交基弱收敛到 0}
            设 \(H\) 是一个\hyperref[dfn:separable]{可分}的 Hilbert 空间, \(\{e_n\}_{n=1}^{\infty}\) 是 \(H\) 的一个标准正交基. 则 \(\{e_n\}_{n=1}^{\infty}\wlim 0\).
        \end{thm}

        \begin{prf}
            我们仅对 \(\ell^2\) 空间证明. 根据 \hyperref[thm:riesz-hilbert]{Riesz 表示定理}, 存在唯一 \(a\in\ell^2\) 使得
            \[T e_n = \langle a, e_n \rangle\]
            当 \(n\to\infty\), \(Te_n\to 0\), 从而 \(\{e_n\}_{n=1}^{\infty}\wlim 0\).
        \end{prf}

        \begin{rmk}
            [Dedicatia]
            令 \(H=\Lp[2]\), 取标准正交基为 \(\{1, \cos x, \sin x, \cdots, \cos nx, \sin nx, \cdots\}\), 那么这个结论就是著名的 Riemann-Lebesgue 引理:
            \[\int_0^{2\pi} \cos nx \cdot g(x) \dd{x} \to 0, \forall\,g\in\Lp[2]([0,2\pi]),\]
            \[\int_0^{2\pi} \sin nx \cdot g(x) \dd{x} \to 0, \forall\,g\in\Lp[2]([0,2\pi]).\]
        \end{rmk}

        \begin{cxmp}
            {弱收敛不蕴含范数收敛}
            考虑 \(\Lp[2]([0,1])\) 中的函数列 \(\{\sin 2n\pi x\}\), 则根据 Riemann-Lebesgue 引理我们有
            \[\int_0^1 \sin 2n\pi x \cdot g(x) \dd{x} \to 0, \forall\,g\in\Lp[2]([0,1]).\]
            从而 \(\{\sin 2n\pi x\}\wlim 0\). 但是 \(\norm{\sin 2n\pi x}_2 = \frac{1}{\sqrt{2}}\) 不趋于零, 从而 \(\{\sin 2n\pi x\}\) 弱收敛但不范数收敛.
        \end{cxmp}

        \begin{ins}
            {\(\ell^1\) 空间上的弱收敛与范数收敛等价}
        \end{ins}

        \begin{thm}
            [banach-alaoglu]
            {Banach-Alaoglu 定理}
            [Banach-Alaoglu Theorem]
            设 \(X\) 是 Banach 空间, 则 \(X^*\) 的单位闭球 \(B=\overline{B}(0,1):=\{f\in X^*: \norm{f}\leq 1\}\) 在弱*拓扑下是\hyperref[dfn:compact]{紧的}.
        \end{thm}

        \begin{prf}
            对任意 \(x\in X\), 定义
            \[D_x:=\{a\in\R, |a|\leq\norm{x}\}\]
            显然这是一个紧致集. 那么积拓扑空间 
            \[\prod_{x\in X} D_x=:D\]
            根据 \hyperref[thm:tychonoff]{Tychonoff 定理}, \(D\) 是一个紧空间. 这时, 因为所有的 \(a\leq\norm{x}\) 可以看成是某个泛函 \(\phi\) 的值域, 所以 \(D\) 中的元素是所有将 \(x\) 映到 \(\R\) 的函数 \(\phi\), 且满足 \(|\phi(x)|\leq\norm{x}\). \\
            并且, 这里的积拓扑是逐点的, 恰好与弱*拓扑一致, 所以干脆将包含了所有有界线性泛函的 \(X^*\) 的单位闭球 \(B\) 看作是 \(D\) 的一个子集.
            我们要证明 \(B\) 是紧的, 根据定理\ref{thm:closed-subset-of-compact}, 只需证明 \(B\) 是 \(D\) 的一个闭子空间, 也就是要证明线性泛函在一族有界泛函中逐点收敛到线性泛函, 这是显而易见的. \\
            从而 \(B\) 是 \(D\) 的一个闭子空间, 也是一个紧空间.
        \end{prf}

        \begin{rmk}
            [Dedicatia]
            实际上, 这里定理证明不需要完备性. 该定理对一般赋范线性空间也成立.
        \end{rmk}

        \begin{varthm}
            {弱*闭集是紧的当且仅当其在范数拓扑下有界}
            设 \(X\) 是赋范线性空间, 则弱*闭集 \(S\subseteq X^*\) 是弱*紧集当且仅当其在范数拓扑下有界.
        \end{varthm}

        \begin{prf}
            (\(\impliedby\)): 若 \(S\) 是有界的, 那么 \(S\subseteq B_r\) 包含在了某个有界闭球之内. 根据 \hyperref[thm:banach-alaoglu]{Banach-Alaoglu 定理}, \(B_r\) 是弱*紧的, 紧集的闭子集也是紧的, 所以 \(S\) 也是弱*紧的.\\
            (\(\implies\)): 设 \(S\) 是弱*紧的, 那么 \(S\) 中的泛函 \(f\) 对 \(x\in X\) 的像 \(f(x)\) 也是紧的, 因为连续映射将紧集映到紧集. \(\R\) 上的紧集是有界闭的, 所以 \(\forall\, x\in X\), \(f(x)\) 逐点有界. 根据\hyperref[thm:banach-steinhaus]{一致有界性定理}, \(\norm{f}\) 一致有界. 所以 \(S\) 在范数拓扑下有界.
        \end{prf}

        \begin{crl}
            {自反空间的单位闭球是弱紧的}
            [Weak Compactness of the Unit Ball in Hilbert Spaces]
            这是因为原空间到其\hyperref[dfn:dual-space]{对偶空间}存在同构, 从而自反空间的单位闭球在弱拓扑下也是\hyperref[dfn:compact]{紧的}. \\
            特别地, Hilbert 空间的单位闭球是弱紧的.
        \end{crl}

        \begin{varthm}
            {弱收敛的点列是一致有界的}
            设 $\{x_n\}$ 是一列赋范线性空间 \(X\) 上的弱收敛序列, 则 \(\{x_n\}\) 是一致有界的.
        \end{varthm}

        \begin{prf}
            根据\hyperref[thm:banach-steinhaus]{一致有界性定理}, 对一列 \(\{x_n\}\), 若对任意 \(f\in X^*\) 都有 
            \[|f(x_n)|\leq M\norm{f}, \forall x_n\]
            那么 \(\{x_n\}\) 有一致的界. \(\exists\,c\in\R\) s.t. \(\norm{x_n}\leq c,\forall x_n\).
        \end{prf}

        \begin{varthm}
            {弱收敛极限的下半连续性}
            设 \(X\) 是赋范线性空间, \(\{x_n\}\subset X\), \(x_n\wlim x\). 那么有
            \[\norm{x}\leq\liminf\norm{x_n}. \]
        \end{varthm}

        \begin{prf}
            根据定理\ref{thm:norming-functional}, \(\exists\,f\in X^*\) s.t. \(\norm{f}=1\) \(\land\) \(\norm{x}=\abs{f(x)}\). 则
            \[f(x)=\lim_{n\to\infty}f(x_n), \abs{f(x_n)}\leq\norm{f}\norm{x_n}=\norm{x_n} \]
            所以下半连续性成立.
        \end{prf}

        \begin{rmk}
            [Dedicatia]
            这些定理同样有对弱*收敛的表述. 不再赘述.
        \end{rmk}

        \begin{thm}
            [mazur]
            {Mazur 定理}
            [Mazur's Theorem]
            Banach 空间 \(X\) 上的闭凸集 \(K\) 一定是弱闭的.
        \end{thm}

        \begin{prf}
            我们尝试证明对 \(z\in X\), 如果 \(z\notin K\), 那么 \(z\) 也不在 \(K\) 在弱拓扑的闭包之中. \\
            由于 \(K\) 是范数拓扑下的闭集, 所以存在 \(z\) 的开邻域 \(B(z,r)\) 使得 \(B(z,r)\cap K=\varnothing\). 根据\hyperref[thm:convex-set-separation]{凸集分离定理}, \(\exists f\in X^*\) \(c\in\R\), s.t.
            \[f(u)\leq c\leq f(v), \forall\,u\in K, v\in B(z,r). \]
            对于 \(v\in B(z,r)\), \(v=z+x\), \(\norm{x}\leq r\). 那么我们有
            \[\inf_{v\in B(z,r)}f(v)=f(z)+\inf_{|x|\leq R}f(x)=f(z)-R\norm{f}\]
            这说明
            \[f(z)\geq c\]
            所以 \(z\in\{x: f(x)\geq c\}\), 这正是一个弱开集. 但是根据分离性, 这个集合中没有任何 \(K\) 中的元素. 由此得到了证明.
        \end{prf}

        \begin{thm}
            {Goldstine 定理}
            [Goldstine's Theorem]
            设 \(X\) 是赋范线性空间, 则 \(X\) 的单位闭球 \(\overline{B}_X\) 的典范嵌入的像 \(J_X(\overline{B}_X)\) 在 \(X^{**}\) 的单位闭球 \(\overline{B}_{X^{**}}\) 中是弱*稠密的. 也就是
            \[\overline{J(\overline{B}_{X})}^{w^*}=\overline{B}_{X^{**}}.\]
            即: \(\forall\,f\in X^*, \exists\,\{x_\alpha\}\subset\overline{B}_X\st f(x_\alpha)\to f(0)\).
        \end{thm}

        \begin{prf}
            (\(\implies\)): 由于 \(J\) 是等距的嵌入, 所以自然有 
            \[J_X(\overline{B}_X)\subseteq \overline{B}_{X^{**}}.\]
            根据 \hyperref[thm:banach-alaoglu]{Banach-Alaoglu 定理}, 单位闭球是弱*紧的, 自然也是闭的. 所以取闭包仍然包含在 \(B_{X^{**}}\) 内. 即
            \[\overline{J_X(\overline{B}_X)}\subseteq \overline{B}_{X^{**}}.\]
            %\[J_X(\overline{B}_X)\subseteq \overline{B}_{X^{**}}\implies J_X(\overline{B}_X)^{w^*}\subseteq (\overline{B}_{X^{**}})^{w^*}=\overline{B}_{X^{**}}.\]
            (\(\impliedby\)): 设存在 \(\xi\in\overline{B}_{X^{**}}\setminus\overline{J_X(\overline{B}_X)}^{w^*} \). 则根据 \hyperref[thm:hahn-banach]{Hahn-Banach 定理}, 存在一个将装备了弱*拓扑的 \(X^{**}\) 映到 \(\R\) 的连续线性泛函 \(F\), 和 \(a\in\R\) 使得 
            \[F(\xi)>a\geq\sup_{y\in J_X(\overline{B}_X)} F(y) \]
            而装备了弱*拓扑的 \(X^{**}\) 的对偶空间正好是 \(X^*\) (这一点验证比较繁琐, 此处从略). 这说明 \(\forall\eta\in X^{**}\), \(F(\eta)=\eta(f)\). 那么分离条件变为
            \[\xi(f)>a\geq\sup_{\norm{x}\leq 1} J_X(x)(f)=\sup_{\norm{x}\leq 1} f(x)\]
            这时, \(\xi(f)>\norm{f}\). 但是
            \[\abs{\xi(f)}\leq\norm{\xi}\norm{f}\leq\norm{f} \]
            这就导致了矛盾. 这说明这样的 \(\xi\) 不存在. 所以 \(\overline{B}_{X^{**}}\subseteq\overline{J(\overline{B}_{X})}^{w^*} \).
            \iffalse
            这二者可以通过开集分离. 设 \(U\) 是弱*开集, \(U\cap (J_X(\overline{B}_X))^{w^*} =\varnothing\), \(\xi\in U\). 则
            \[U = \bigcap_{i=1}^n \{F\in X^{**}: |F(x_i)-\xi(x_i)|<\epsilon\}\]
            其中 \(x_1,\dots,x_n\in X\), \(\epsilon>0\). 定义
            \[V = \bigcap_{i=1}^n \{x\in X: |x_i(x)-\xi(x_i)|<\epsilon\}\]
            则 \(V\) 是 \(X\) 中的一个弱邻域, 且 \(V\cap \overline{B}_X = \varnothing\). 这就说明了 \(0\in X\) 和 \(\overline{B}_X\) 可以通过开集分离. 这显然是矛盾的. 从而 \(\xi\in (J_X(\overline{B}_X))^{w^*}\).
            \fi
        \end{prf}

        \begin{thm}
            {可分与可度量的关系}
            设 \(X\) 是赋范线性空间, 则:
            \begin{enumerate}
                \item \(X\) 在范数拓扑下可分 \(\iff\) \(X^*\) 在弱*拓扑的单位闭球是可度量的;
                \item \(X^*\) 在范数拓扑下可分 \(\iff\) \(X\) 在弱拓扑的单位闭球是可度量的.
            \end{enumerate}
        \end{thm}

        \begin{prf}
            只证明 (1). \\
            \(\implies\): 设 \(X\) 是可分的, 则存在可数的稠密子集
            \[\{x_n\}_{n=1}^{\infty}\subseteq\overline{B}_X \]
            由于 \(X^*\) 上的弱*拓扑就是逐点的拓扑, 假使 \(f,g\in X^*\) 的接近程度可以使用这些点 \(\{x_n\}\) 来刻画, 那我们的目的也就达成了. 事实上直接三段控制即可: 只要 \(f\to g\) 在 \(x_n\) 上都成立, 就有
            \[|f(x)-g(x)|\leq|f(x)-f(x_k)|+|f(x_k)-g(x_k)|+|g(x)-g(x_k)| =3\varepsilon \]
            另外两项成立是因为
            \[|f(x)-f(x_k)|\leq\norm{f}\norm{x-x_k} \]
            由于 \(\{x_k\}\) 是稠密的, 所以 \(x_k\) 可以任意接近 \(x\). \(f,g\) 在单位闭球中, 所以是有界的, 这两项也就成为无穷小.
            这表明, 可以认为 \(f\to g\) 在 \(\overline{B}_{X^*}\) 上成立. 所以一个度量可以选用
            \[d(f,g)=\sum_{n=1}^{\infty}2^{-n}\frac{|f(x_n)-g(x_n)|}{1+|f(x_n)-g(x_n)|}.\]
            \(\impliedby\): 设 \(X^*\) 的单位闭球在弱拓扑下可度量, 则满足第一可数性质, \(0\in \overline{B}_{X^*}\) 存在可数的邻域基, 记作 \(\{U_n\}_{n=1}^{\infty}\). \\
            按照弱*邻域定义, \(0\) 的一个弱*邻域应该是有限个向量就可以表达的:
            \[V_n:=\qty{f\in \overline{B}_{X^*}: |f(x_{n,1})|<\epsilon_n, \cdots, |f(x_{n,m_n})|<\epsilon_n }\]
            可数个由有限个向量描述的拓扑应该也可以被有限个向量描述. 也就是说, 所有邻域中的所有用来确定该邻域的向量收集起来得到
            \[D:=\{x_{n,j}, n\in\N, 1\leq j\leq m_n\}\]
            是可数的.\\
            现在考虑证明 \(D\) 确实可以线性张成 \(X\) 上的稠密子集. 假如不是, 那么存在 \(x_0\in X\setminus\overline{\spn(D)}\), 根据 Hahn-Banach 定理, \(\exists\,f\in X^*, f\neq 0\) 使得 \(f\) 限制在 \(\overline{\spn(D)}\) 上是 \(0\). \\
            \(f(x)=0, \forall\,x\in D\), 那么 \(f\) 在 \(0\in X^*\) 的所有邻域中. 虽然弱*拓扑非常弱, 但仍旧是 Hausdorff 空间. 所以 \(f=0\), 这就导致了矛盾. \\
            所以, 若 \(f\neq g\) 在 \(D\) 上, 那么 \(f\neq g\) 在 \(X\) 上都成立. \(D\) 就是一个可数的稠密子集. 所以 \(X\) 是可分的.
            \iffalse
            只证明 (2). \\
            \((\implies)\): 设 \(X^*\) 是可分的, 则存在可数的稠密子集
            \[\{f_n\}_{n=1}^{\infty}\subseteq X^*\]
            且满足 \(\norm{f_n}=1\), 这是可以办到的, 若不然, 只需对 \(f_n\) 利用标量乘法单位化即可.\\
            定义
            \[d(x,y) = \sum_{n=1}^{\infty} \frac{1}{2^n} |f_n(x-y)|\]
            则 \(d\) 是 \(X\) 上的一个度量, 且它诱导的拓扑与 \(X\) 上的弱拓扑一致. 从而 \(X\) 在弱拓扑的单位闭球是可度量的. \\
            (\(\impliedby\)): 设 \(X\) 的单位闭球在弱拓扑下可度量, 则满足第一可数性质, \(0\in \overline{B}_X\) 存在可数的邻域基, 记作 \(\{V_m\}_{m=1}^{\infty}\). 按照定义,
            \[\forall\,m\in\N, \exists\{f_{m,i}\}_{i=1}^{N}\subset X^*, \forall\epsilon_m>0, U_m:=\{x\in\overline{B}_X: |f_{m,i}(x)| < \epsilon_m \} \subseteq V_m\]
            定义
            \[F_m:= \qty{f_{m,1}, f_{m,2}, \ldots, f_{m,N}}\]
            那么
            \[M:=\bigcup_{m=1}^{\infty}F_m\]
            是可数个有限集的并, 也是可数的. 令 \(Y=\overline{\spn(M)}\). 则 \(Y\) 是可分的. \\
            假设 \(Y\neq X^*\). 那么存在 \(\phi\in X^{**}\) 使得 \(\norm{\phi}=1\), \(\forall\,f\in M\), \(\phi(f)=0\). \\
            根据 Goldstine 定理, \(\exists\,\{x_\alpha\}\subset\overline{B}_X\) 使得 \(x_\alpha\wslim\phi\). 即: \(\forall\,f\in X^*\), \(f(x_\alpha)\to \phi(f)\): \(f\in M\) 时:
            \[f(x_\alpha)\to 0\]
            但是 \(\{V_m\}\) 是原点的邻域基, \(x_\alpha\) 在 \(\alpha\) 足够大时一定会进入 \(V_m\), 这说明 \(x_\alpha\to 0\). 那么 \(\phi(g)=0, \forall\,g\in X^*\). 这只有一种可能就是 \(\phi=0\), 与 \(\norm{\phi}=1\) 矛盾. 从而 \(Y=X^*\), \(X^*\) 是可分的. 
            \fi
        \end{prf}

        \begin{rmk}
            [Dedicatia]
            我们可以列表比较:
            \begin{center}
                \begin{tabular}{c|ccc}
                    \toprule
                    拓扑 & 范数拓扑 & 全空间的弱拓扑 & 单位闭球的弱拓扑 \\
                    \midrule
                    第一可数性质 & \checkmark & \texttimes & \checkmark \\
                    可分 & \checkmark & \texttimes & \checkmark \\
                    可度量 & \checkmark & \texttimes & \checkmark \\ 
                    \bottomrule
                \end{tabular}
            \end{center}
            这体现了弱拓扑的好处, 当限制在单位闭球上时, 我们希望的一些好性质又回来了.
        \end{rmk}

        \begin{crl}
            {Banach-Alaoglu 定理在可分空间中的表述}
            设 \(X\) 是一个\hyperref[dfn:separable]{可分}的赋范线性空间, 则 \(X^*\) 的单位闭球在弱*拓扑下是\hyperref[dfn:sequential-compact]{列紧的}. 即对任意一列对偶空间中的元素 \(\{f_n\}_{n=1}^{\infty}\), 存在一列 \(f_{n_k}\subset \{f_n\}\) 使得 \(f_{n_k}\wslim f\) 对某个 \(f\in X^*\) 成立.
        \end{crl}

        \begin{ins}
            {序列空间的可度量性}
            我们考虑序列空间 \(\ell^1\) 与 \(\ell^\infty\) 的单位闭球的拓扑是否可度量化:
            \begin{center}
                \begin{tabular}{c|cc}
                    \toprule
                    单位球上的拓扑 & 等价条件 & 可度量 (\checkmark / \texttimes) \\
                    \bottomrule
                    \((\ell^1, w)\) & \(\ell^\infty\) 是否可分? & \texttimes \\
                    \((\ell^1, w^*)\) & \(c_0\) 是否可分? & \checkmark \\
                    \((\ell^\infty, w)\) & \((\ell^\infty)^*\) 是否可分? & \texttimes \\
                    \((\ell^\infty, w^*)\) & \(\ell^1\) 是否可分? & \checkmark \\
                    \bottomrule
                \end{tabular}
            \end{center}
            选择这两个空间的原因是它们都不是自反的. 并且我们有 \((c_0)^*=\ell^1\), \((\ell^1)^*=\ell^\infty\). 
        \end{ins}

        \begin{ins}
            {可积函数空间的可度量性}
            我们这次考虑可积函数空间 \(\Lp[1]([0,1])\) 和 \(\Lp[\infty]([0,1])\) 的单位闭球的拓扑是否可度量化:
            \begin{center}
                \begin{tabular}{c|cc}
                    \toprule
                    单位球上的拓扑 & 等价条件 & 可度量 (\checkmark / \texttimes) \\
                    \bottomrule
                    \((\Lp[1], w)\) & \((\Lp[1])^*=\Lp[\infty]\) 是否可分? & \texttimes \\
                    \((\Lp[\infty], w)\) & \((\Lp[\infty])^*\) 是否可分? & \texttimes \\
                    \((\Lp[\infty], w^*)\) & \(\Lp[1]\) 是否可分? & \checkmark \\
                    \bottomrule
                \end{tabular}
            \end{center}
            这里有一个关键差异, 在 Lebesgue 测度下, \(\Lp[1] ([0,1])\) 不是任何空间的对偶空间, 这一点将在例子\ref{cxmp:lp1-not-dual}看到. 所以它不存在弱*拓扑.
        \end{ins}

        \begin{dfn}
            {凸组合}
            [Convex Combination]
            设 \(X\) 是一个集合, \(x_1,\dots,x_n\in X\), \(\lambda_1,\dots,\lambda_n\geq 0\), \(\sum_{i=1}^n \lambda_i = 1\). 定义
            \[\sum_{i=1}^n \lambda_i x_i\]
            是 \(x_1,\dots,x_n\) 的一个\textbf{凸组合}.
        \end{dfn}

        \begin{dfn}
            [extreme-point]
            {凸集的端点}
            [Extreme Point of a Convex Set]
            设 \(A\) 是线性空间 \(X\) 中的凸集. 定义 \(x\in A\) 是 \(A\) 的一个\textbf{端点}, 当且仅当 \(x\) 不是 \(A\) 中任意两点的凸组合. 也就是
            \[\forall\,y,z\in A, \lambda\in(0,1), \quad x = \lambda y + (1-\lambda) z \implies x = y \lor x = z.\]
            凸集 \(A\) 的全体端点记作 \(\ext(A)\).
        \end{dfn}

        \begin{dfn}
            [convex-hull]
            {凸包}
            [Convex Hull]
            设 \(X\) 是一个线性空间, \(A\subseteq X\). 定义 \(A\) 的\textbf{凸包}是包含 \(A\) 的所有凸集的交集. 记作 \(\conv(A)\).\\
            即, \(\conv(A)\) 是 \(A\) 中所有凸组合的全体.
        \end{dfn}

        \begin{thm}
            {Krein-Milman 定理}
            [Krein-Milman Theorem]
            设 \(X\) 是一个赋范线性空间, \(A\subseteq X\) 是一个非空的紧凸集. 则 \(A = \overline{\conv(\ext(A))}\).
        \end{thm}

        \begin{cxmp}
            [lp1-not-dual]
            {\(\Lp[1]\) 空间不是任何 Banach 空间的对偶空间}
            区间 \([0,1]\) 上赋予 Lebesgue 测度下的 \(\Lp[1]([0,1])\) 函数空间的单位闭球没有任何端点. \\
            这是因为如果 \(f\in\Lp[1]([0,1])\) 是一个端点, 那么可以令 \(E\subseteq [0,1]\), 使得 \(\int_E|f|=\frac{1}{2}\int_{[0,1]}|f|\). 设 \(g = f\cdot \chi_E\), \(h = f\cdot \chi_{[0,1]\setminus E}\). 则 \(f = \frac{1}{2}(2g) + \frac{1}{2}(2h)\), 这说明了 \(f\) 不是一个端点. 从而 \(\ext(\overline{B}_{\Lp[1]})=\varnothing\). \\
            根据 Krein-Milman 定理, 一个紧凸集必须有端点. 假如 \(\Lp[1]([0,1])\) 是某个 Banach 空间的对偶空间, 那么它的单位闭球在弱*拓扑下是紧的. 但是它没有端点, 与 Krein-Milman 定理矛盾. 从而 \(\Lp[1]([0,1])\) 不是任何 Banach 空间的对偶空间. 相应, \(\Lp[1]([0,1])\) 不存在弱*拓扑.
        \end{cxmp}

    \subsection{【实分析】Lp 空间}

        \begin{intrormk}
            [Dedicatia]
            在实分析中, 最常用的具有研究价值的空间就是 \(\Lp\) 可积函数空间. 本节我们将以 \(\Lp\) 空间为例, 来介绍弱拓扑和弱*拓扑在实分析中的应用. 通过对 \(\Lp\) 空间中函数的弱收敛和弱*收敛的研究, 我们可以更深入地理解函数空间的结构和性质.
        \end{intrormk}

        \begin{ins}
            [lp]
            {Lp 空间}
            设 \((X, \mathcal{M}, m)\) 是测度空间, \(1\leq p < +\infty\). 定义 \(\Lp(X)\) 是所有满足 \(\int_X |f|^p \dd{m} < +\infty\) 的可测函数 \(f\) 的集合. 定义 \(L^\infty(X)\) 是所有满足 \(\esssup\limits_{x\in X} |f(x)| < +\infty\) 的可测函数 \(f\) 的集合, 统称为 \(\Lp\) 空间.
        \end{ins}

        \begin{varthm}
            {Lp 空间是赋范线性空间}
            \(\Lp\) 空间上定义范数
            \[\norm{f}_p = \begin{cases}
                \left(\int_X |f|^p \dd{m}\right)^{\frac{1}{p}}, & 1\leq p < +\infty;\\
                \esssup\limits_{x\in X} |f(x)|, & p = +\infty.
            \end{cases}\]
            则 \(\Lp(X)\) 是一个赋范线性空间.
        \end{varthm}

        \begin{prf}
            先证明 Holder 不等式
            \[\int_X |fg| \dd{m} \leq \norm{f}_p \norm{g}_q\]
            其中 \(1<p,q<\infty\). 设 \(y=\phi(x)\) 是一个严格递增的连续函数, \(x=\psi(y)\) 是其反函数. 则下列不等式成立:
            \[\int_0^a \phi(x)\dd{x}+\int_0^b \psi(y)\dd{y} \geq ab\]
            取等条件为 \(b=\phi(a)\). 则令 \(\phi(x)=x^{p-1}, \psi(y)=y^{q-1}\), 得到
            \[\int_0^a x^{p-1}\dd{x}+\int_0^b y^{q-1}\dd{y} = \frac{a^p}{p} + \frac{b^q}{q} \geq ab\]
            记作
            \[A^{1/p}B^{1/q} \geq \frac{A}{p}+\frac{B}{q}.\]
            现在来证明 Holder 不等式. 如果 \(\norm{f}_p\) 或 \(\norm{g}_q\) 中至少有一个为零, 则不等式显然成立. \\
            否则, 取 \(\phi(x)=\frac{f(x)}{\norm{f}_p}\), \(\psi(y)=\frac{g(y)}{\norm{g}_q}\), \(A=|\phi(x)|^p\), \(B=|\psi(x)|^q\), 则
            \[|\phi(x)\psi(x)|\leq\frac{|\phi(x)|^p}{p}+\frac{|\psi(x)|^q}{q}\]
            积分, 得
            \[\int_X |\phi\psi| \dd{m} \leq \frac{1}{p}\int_X |\phi|^p \dd{m} + \frac{1}{q}\int_X |\psi|^q \dd{m} = 1.\]
            换回 \(f,g\), 得到
            \[\int_X |fg| \dd{m} \leq \norm{f}_p \norm{g}_q.\]
            我们验证范数前两个性质:
            \begin{enumerate}[label=(\roman*)]
                \item 正定性: 由于是对非负函数积分, 所以显然是非负的. 由于将 Lp 空间中的元素视作等价类, 所以 \(\norm{f}_p=0\) 当且仅当 \(f=0\) a.e., 正定性得证;
                \item 绝对齐次性: \(\norm{\lambda f}_p = |\lambda| \norm{f}_p\) 显然成立.
            \end{enumerate}
            接下来证明 Minkowski 不等式
            \[\norm{f+g}_p \leq \norm{f}_p + \norm{g}_p.\]
            特殊情况1: 如果 \(p=1\), 则不等式显然成立. \\
            特殊情况2: 如果 \(\norm{f+g}_p=0\), 不等式也成立. \\
            令 \(1/p+1/q=1\), 由 Holder 不等式得到
            \[\int_X|f||f+g|^{p/q}\dd{m}\leq\qty(\int_X|f|^p\dd{m})^{1/p}\qty(\int_X|f+g|^q\dd{m})^{1/q}\]
            \[\int_X|g||f+g|^{p/q}\dd{m}\leq\qty(\int_X|g|^p\dd{m})^{1/p}\qty(\int_X|f+g|^q\dd{m})^{1/q}\]
            从而
            \[\int_X |f+g|^p\dd{m}=\int_X |f+g|^{1+\frac{p}{q}}\dd{m}\leq\int_X(|f|+|g|)\cdot|f+g|^{p/q}\dd{m}\]
            由 Holder 不等式, 得到
            \[\int_X |f+g|^p\dd{m} \leq (\norm{f}_p + \norm{g}_p) \norm{f+g}_p^{p-1}.\]
            除以 \(\norm{f+g}_p^{p-1}\), 得到
            \[\norm{f+g}_p \leq \norm{f}_p + \norm{g}_p.\]
            这就说明了范数满足次可加性. 所以 \(\Lp(X)\) 是一个赋范线性空间.
        \end{prf}

        \begin{rmk}
            [Dedicatia]
            习惯上称 Lp 空间中的依范数收敛为\textbf{均方收敛} / \textbf{平均收敛} / \textbf{Lp 收敛}, 以区别于其他的收敛方式, i.e. 依测度收敛和几乎处处收敛.
        \end{rmk}

        \begin{vardfn}
            {依测度收敛}
            [Convergence in Measure]
            设 \((X, \mathcal{M}, m)\) 是测度空间, \(\{f_n\}_{n=1}^{\infty}\) 是 \(\Lp(X)\) 中的可积函数列, \(f\in \Lp(X)\). 定义 \(\{f_n\}_{n=1}^{\infty}\) \textbf{依测度收敛}到 \(f\), 当且仅当 
            \[\forall\,\epsilon>0, m(\{x\in X: |f_n(x)-f(x)|>\epsilon\})\to 0.\]
            记作 \(f_n\mlim f\).
        \end{vardfn}

        \begin{varthm}
            {均方收敛蕴含依测度收敛}
            设 \(\{f_n\}\) 是 \(\Lp(X)\) 中的可积函数列, \(f\in \Lp(X)\). 则
            \[\int_X |f_n-f|^p \dd{m} \to 0 \implies f_n\mlim f.\]
        \end{varthm}

        \begin{prf}
            任取 \(\delta>0\), 
            \[\int_X|f_n-f|^p\dd{m}\geq\int_{X\cap \{x:|f_n(x)-f(x)|\geq\delta\}}|f_n-f|^p\dd{m}\geq\delta^p m\qty(\{x:|f_n(x)-f(x)|\geq\delta\}).\]
            令 \(n\to\infty\).
        \end{prf}

        \begin{rmk}
            [Dedicatia]
            依测度收敛是很弱的收敛方式. 为 \(X=[0,1]\) 配备 Lebesgue 测度, 则得到\textbf{概率空间}, 依测度收敛就是概率意义下的收敛, 这时也称为也\textbf{依概率收敛}. 它的含义是某件事情发生的概率越来越大, 但并不能阻止它不发生.
        \end{rmk}

        \begin{varthm}
            {Lp 空间是 Banach 空间}
            \(\Lp(X)\) 是一个 Banach 空间.
        \end{varthm}

        \begin{prf}
            我们这里仅证明 \(1\leq n<\infty\) 的情况:
            设 \(\{f_n\}\) 是 Cauchy 列. 对任何 \(\delta>0\), 
            \[m(\{x\in X : |f_n(x)-f_m(x)|\geq\delta\}) < \frac{1}{\delta}\qty(\int_X |f_n-f_m|^p \dd{m})^{1/p}=\frac{1}{\delta} \norm{f_n-f_m}_p\]
            这表明了 \(\{f_n\}\) 是依测度收敛的 Cauchy 列. 设 \(f\) 是 \(\{f_n\}\) 的一个依测度收敛的极限. 则根据 F.Riesz 定理, 存在一个子列 \(\{f_{n_k}\}\) 使得 \(\{f_{n_k}\}\) 几乎处处收敛到 \(f\). 取 
            \[\norm{f_n-f_{n_k}}<\epsilon\]
            根据 Fatou 引理, 固定 \(n\),
            \[\epsilon\geq\lim_{k\to\infty}\norm{f_{n_k}-f_n}_p=\lim_{k\to\infty}\qty(\int_X |f_{n_k}-f_n|^p \dd{m})^{1/p}=\qty(\lim_{k\to\infty}\int_X |f_{n_k}-f_n|^p \dd{m})^{1/p}\]
            \[\geq\qty(\int_X |f_n-f|^p \dd{m})^{1/p}=\norm{f_n-f}_p\]
            这表明了当 \(n\) 足够大的时候, \(\norm{f_n-f}_p<\epsilon\). 从而 \(\{f_n\}\) 依范数收敛到 \(f\). 这就说明了 \(\Lp(X)\) 是完备的.
        \end{prf}

        \begin{thm}
            {Lp 空间的 Riesz 表示定理}
            设 \((X, \mathcal{M}, m)\) 是测度空间, \(1\leq p < +\infty\), 则 \(\Lp(X)^* \cong L^q(X)\), 其中 \(1/p+1/q=1\). 
        \end{thm}

        \begin{prf}
            我们只需给出一个等距同构:
            定义
            \[T:\qquad\begin{matrix}
                L^q(X) & \to & \Lp(X)^*\\
                g & \mapsto & T_g:=\qty(f\mapsto\int_X fg \dd{m})
            \end{matrix}\]
            其证明与定理\ref{thm:riesz-hilbert}类似.
        \end{prf}

        \begin{ins}
            {L2 空间是 Lp 空间中唯一的 Hilbert 空间}
            \(\Lp[2](X)\) 是一个 Hilbert 空间, 内积定义为
            \[\langle f, g\rangle = \int_X f\cdot\overline{g} \dd{m}.\]
            这是因为只有 \(\Lp[2]\) 空间的对偶是自身.
        \end{ins}

        \begin{thm}
            {Lp 空间的自反性}
            Lp 空间是自反的当且仅当 \(1<p<\infty\).
        \end{thm}

        \begin{cxmp}
            {L1 空间的非自反性}
            \(\Lp[1](X)^*=\Lp[\infty](X)\) 但 \(\Lp[\infty](X)^* \neq \Lp[1](X)\). \\
            这是因为 \(\Lp[1](X)\) 是\hyperref[dfn:separable]{可分}的, 但 \(\Lp[\infty](X)\) 是不可分的: 取 \(X=[0,1]\), 则区间 \([0,\lambda]\) 的特征函数 \(\{\I_{[0,\lambda]}: \lambda\in[0,1]\}\) 是 \(\Lp[\infty](X)\) 中的一族元素, 且任意两个元素之间的距离都是 1, 从而 \(\Lp[\infty](X)\) 中存在一个不可数的 1-分离子集, 从而 \(\Lp[\infty](X)\) 是不可分的.
        \end{cxmp}

        \begin{rmk}
            [Dedicatia]
            总结下来, 我们有如下的表格:
            \begin{center}
                \begin{tabular}{c|cccc}
                    \toprule
                    空间 & \(\Lp[1]\) & \(\Lp[2]\) & 其余的 \(\Lp[p], 1<p<\infty\) & \(\Lp[\infty]\)\\
                    \midrule
                    Banach 空间 (完备的范数) & \checkmark & \checkmark & \checkmark & \checkmark\\
                    Hilbert 空间 (完备的内积) & \texttimes & \checkmark & \texttimes & \texttimes\\
                    \hyperref[dfn:separable]{可分性} & \checkmark & \checkmark & \checkmark & \texttimes\\
                    \hyperref[dfn:reflexive]{自反性} & \texttimes & \checkmark & \checkmark & \texttimes\\
                    单位球的一致凸性 & \texttimes & \checkmark & \checkmark & \texttimes\\
                    \bottomrule
                \end{tabular}
            \end{center}
        \end{rmk}

    \subsection{算子拓扑}

        \begin{intrormk}
            [Dedicatia]
            我们之前介绍的弱拓扑和弱*拓扑都是关于线性空间或线性泛函的拓扑. 如今我们将研究线性算子 \(\blo(X,Y)\) 作为元素的空间上的拓扑和收敛性.
        \end{intrormk}

        \begin{dfn}
            {一致拓扑 / 范数拓扑}
            [Uniform Topology / Norm Topology]
            设 \(X,Y\) 是 Banach 空间, 则 \(\blo(X,Y)\) 上的\textbf{一致拓扑}是由范数 \(\norm{T} = \sup_{\norm{x}\leq 1} \norm{Tx}\) 诱导的拓扑. 记作 \((\blo(X,Y), \tau_{\norm{\cdot}})\).
        \end{dfn}

        \begin{dfn}
            {强算子拓扑}
            [Strong Operator Topology]
            设 \(X,Y\) 是 Banach 空间, 则 \(\blo(X,Y)\) 上的\textbf{强算子拓扑}是由半范数 \(\norm{T}_x = \norm{Tx}\) 诱导的拓扑, 其中 \(x\in X\). 记作 \((\blo(X,Y), s)\).
        \end{dfn}

        \begin{dfn}
            {弱算子拓扑}
            [Weak Operator Topology]
            设 \(X,Y\) 是 Banach 空间, 则 \(\blo(X,Y)\) 上的\textbf{弱算子拓扑}是由半范数 \(\norm{T}_{x,f} = |f(Tx)|\) 诱导的拓扑, 其中 \(x\in X, f\in Y^*\). 记作 \((\blo(X,Y), w)\).
        \end{dfn}

        \begin{dfn}
            {一致收敛}
            [Uniform Convergence]
            设 \(X,Y\) 是 Banach 空间, \(\{T_n\}\subset \blo(X,Y)\), \(T\in \blo(X,Y)\). 定义 \(\{T_n\}\) \textbf{一致收敛}到 \(T\), 当且仅当 \(\norm{T_n-T} \to 0\). 记作 \(T_n\rightrightarrows T\).
        \end{dfn}

        \begin{dfn}
            {强收敛}
            [Strong Convergence]
            设 \(X,Y\) 是 Banach 空间, \(\{T_n\}\subset \blo(X,Y)\), \(T\in \blo(X,Y)\). 定义 \(\{T_n\}\) \textbf{强收敛}到 \(T\), 当且仅当 \(\forall\,x\in X, \norm{T_n x-Tx} \to 0\). 记作 \(T_n\slim T\).
        \end{dfn}

        \begin{dfn}
            {弱收敛}
            [Weak Convergence]
            设 \(X,Y\) 是 Banach 空间, \(\{T_n\}\subset \blo(X,Y)\), \(T\in \blo(X,Y)\). 定义 \(\{T_n\}\) \textbf{弱收敛}到 \(T\), 当且仅当 \(\forall\,x\in X, f\in Y^*, f(T_n x) \to f(Tx)\). 记作 \(T_n\wlim T\).
        \end{dfn}
        
        \begin{ppt}
            {三种算子拓扑的关系}
            \[\text{一致拓扑} \implies \text{强算子拓扑} \implies \text{弱算子拓扑}\]
            这说明了不同的算子拓扑之间的包含关系. 一致拓扑是最强的, 强算子拓扑次之, 弱算子拓扑最弱. 这也表明在一致拓扑下收敛的序列在强算子拓扑和弱算子拓扑下也收敛, 但反之不一定成立.
        \end{ppt}

        \begin{cxmp}
            {强收敛不蕴含一致收敛}
            设 \(X=Y=\ell^2\), 定义左移算子
            \[T_n(x_1,x_2,\dots) = (x_{n+1}, x_{n+2}, \dots)\]
            则 \(\forall\,x\in\ell^2\),
            \[\norm{T_n x}^2 = \sum_{k=n+1}^{\infty} |x_k|^2 \to 0\]
            这说明了 \(T_n\slim 0\). 但是 \(\norm{T_n} = 1\) 对所有 \(n\) 都成立, 这说明了算子序列 \(T_n\) 不一致收敛于 \(0\).
        \end{cxmp}

        \begin{cxmp}
            {弱收敛不蕴含强收敛}
            设 \(X=Y=\ell^2\), 定义投影算子
            \[T_nx=x_1e_n\]
            其中 \(\{e_n\}\) 是 \(\ell^2\) 中的标准正交基. 则根据 Riesz 表示定理可知, \(\forall\,x\in\ell^2, f\in\ell^2\), 
            \[f(T_n x) = f(x_1 e_n) = x_1 f(e_n) \to 0\]
            这说明了 \(T_n\wlim 0\). 但是只要 \(n\neq m\), 
            \[\norm{T_ne_1-T_me_1} = \norm{e_n-e_m} = \sqrt{2} \]
            这说明了 \(T_n\) 不强收敛于 \(0\).
        \end{cxmp}

\section{Banach 空间上的谱论}

    \subsection{紧致算子}

        \begin{dfn}
            [compact-operator]
            {紧致算子 / 紧算子}
            [Compact Operator]
            设 \(X\) 和 \(Y\) 是 Banach 空间, 则线性算子 \(T:X\to Y\) 是一个\textbf{紧致算子}, 当且仅当 \(T\) 将 \(X\) 中的有界集合映射到 \(Y\) 中的一个\hyperref[dfn:sequential-compact]{相对紧集}.\\
            记作 \(T\in\cpt(X,Y)\).
        \end{dfn}

        \begin{ppt}
            {紧算子的等价定义}
            $T \in\cpt(X,Y)$ 等价于 $T$ 将 \(X\) 中的单位闭球映射到 \(Y\) 中的一个相对紧集, 等价于 \(T\) 将 \(X\) 中的有界序列映射到 \(Y\) 中的一个有收敛子列的序列.\\
            可以这样定义的原因是赋范线性空间具有标量乘法和平移不变性; 以及 $Y$ 中的相对紧集的定义是对任意有界序列都有在 $Y$ 中收敛的子列. 这实际上是列紧性的定义, 但在度量空间中二者是等价的.
        \end{ppt}

        \begin{ppt}
            {任意有限维有界线性算子都是紧算子}
            显然. 因为有限维空间总是满足 Heine-Borel 性质.
        \end{ppt}

        \begin{cxmp}
            {无限维赋范线性空间上的恒等算子不是紧算子}
            设 \(\dim X= \infty\), 则恒等算子 \(I_X\) 不是紧算子. 因为 \(I_X\) 将 \(X\) 中的单位闭球映射到 \(X\) 中的单位闭球, 但 \(X\) 中的单位闭球不是紧的.
        \end{cxmp}

        \begin{ppt}
            {紧算子的复合也是紧算子}
            设 \(X,Y,Z\) 是 Banach 空间, \(T:\blo(X,Y)\), \(S:\blo(Y,Z)\), 若 \(T,S\) 中至少有一个紧算子, 则 \(S\circ T\in\cpt(X,Z)\).
        \end{ppt}

        \begin{rmk}
            [Dedicatia]
            这说明了在 $\blo(X,X)$ 中, $\cpt(X,X)$ 是一个双边理想.
        \end{rmk}

        \begin{ppt}
            {紧算子是有界线性算子的闭子空间}
            [Compact Operators Form a Closed Subspace of Bounded Linear Operators]
            $\cpt(X,Y)\leq\blo(X,Y)$. 
        \end{ppt}

        \begin{prf}
            设 $T_j \subset \cpt (X,Y)$ 且 $T_j \to T$ 在 $\blo (X,Y)$ 中. \\
            取一列有界序列 ${x_n}_{n=1}^{\infty}$ 在 $X$ 中. 根据对角线法证明即可得到, 存在一个子列 ${x_(n_k)}_{k=1}^{\infty}$ 使得 $T_j x_(n_k)$ 收敛于 $T x_(n_k)$ 对所有 $j$ 成立. 这说明 $T x_(n_k)$ 收敛, 记作 $M:=\sup_{n \in \N} \norm{x_n} $.
            三段控制:
            \[\norm{Tx_{n_k} - Tx_{n_j}} \leq \norm{Tx_{n_k} - T_j x_{n_k}} + \norm{T_j x_{n_k} - T_j x_{n_j}} + \norm{T_j x_{n_j} - Tx_{n_j}}\leq 2M\norm{T-T_j} + \norm{T_j (x_{n_k} - x_{n_j})} \]
            这两项都是无穷小量. 因此在 $k,j\to\infty$ 时, $\norm{Tx_{n_k} - Tx_{n_j}} \to 0$. 这是一个 Cauchy 列, 根据完备性, $T x_{n_k}$ 收敛. 这就说明了 \(T\in\cpt(X,Y)\).
        \end{prf}

        \begin{ppt}
            {紧算子将弱收敛序列映到范数收敛序列}
            设 \(T\in\cpt(X,Y)\). 则 \(x_n\wlim x\implies Tx_n\to Tx\). \\
            其原因很简单: \(\{Tx_n\}\) 是相对紧集中的序列且一定也是弱收敛的, 倘若它不按范数收敛, 那么存在子列远离弱收敛的极限; 但是这个子列也有子列按范数收敛到极限, 这就导致了矛盾.
        \end{ppt}
        
        \begin{dfn}
            [finite-rank-operator]
            {有限秩算子}
            [Finite Rank Operator]
            设 \(X\) 和 \(Y\) 是 Banach 空间, 则线性算子 \(T:X\to Y\) 是一个\textbf{有限秩算子}, 当且仅当 \(\dim\img T<\infty\).\\
            记作 \(\fnrank(X, Y)\).
        \end{dfn}

        \begin{ppt}
            {有限秩算子是紧算子}
            \(T:\fnrank(X, Y)\implies T:\cpt(X, Y)\). \\
            这是因为有限维空间的有界集合一定是相对紧的.
        \end{ppt}

        \begin{thm}
            {Hilbert 空间上的紧算子可以被有限秩算子逼近}
            若 \(X\) 是 Hilbert 空间, 则 \(\overline{\fnrank(X, X)}=\cpt(X,X)\).
        \end{thm}

        \begin{rmk}
            [Dedicatia]
            但是, 一般的 Banach 空间中的紧算子不一定能被有限秩算子逼近, 即使是可分的空间. 具备这种性质的 Banach 空间一般称为有限逼近性质的空间.
        \end{rmk}

        \begin{xmp}
            {\(\Lp[2](D)\) 上的积分算子}
            设 $D \subset \R$ 是区域. $K \in L^2(D \times D)$, 的定义
            \[T_K: L^2(D) \to L^2(D), \qquad (T_K f)(x) = \int_D K(x,y) f(y) \dd{y}\]
            其中 \(K\) 称为积分核, \(T_K\) 称为积分算子, 则 \(T_k\) 是紧算子:\\
            (1) 有界性: 
            \[\norm{T_K f}=\int_D\abs{\int_D K(x,y)f(y)\dd y}^2\dd x\]
            根据 Cauchy-Schwarz 不等式可知
            \[\norm{T_K f}\leq\int_D\qty(\int_D|K(x,y)|^2\dd y)\qty(\int_D|f(y)|^2\dd y)\dd{x}=\norm{K}\norm{f} \]
            这就说明了 $T_K$ 是有界的, 且 $\norm{T_K}\leq\norm{K}$.\\
            (2) 紧性:\\
            由于 \(\Lp[2]\) 是可分的 Hilbert 空间, 所以存在可数的标准正交基 \(\{e_n\}_{n=1}^\infty\). 定义
            \[(\phi_i,\phi_j), \qquad (i,j)\in\N\times\N \]
            这是 \(\Lp[2](D\times D)\) 上的标准正交基, 使得
            \[K=\sum_{i,j\in \N}\left\langle K,(\phi_i, \phi_j) \right\rangle (\phi_i, \phi_j)\]
            由于标准正交基是可数的, 向各个坐标做投影 $(i+j\leq n)$ 得到一列有限秩算子. 所以收敛.
        \end{xmp}

        \begin{dfn}
            {对偶算子}
            [Dual Operator]
            设 $T: X \to Y$ 是一个有界线性算子, 则 $T^*: Y^* \to X^*$ 是 $T$ 的对偶算子, 定义为
            \[ (T^* f) (x) = f (T(x)), \forall f \in Y^* \]
        \end{dfn}

        \begin{rmk}
            [Dedicatia]
            一定要注意其顺序. 括号不能随意减少.
        \end{rmk}

        \begin{thm}
            [schauder]
            {Schauder 定理}
            $T$ 是紧算子 iff $T^*$ 是紧算子.
        \end{thm}

        \begin{prf}
            我们先证明 \(T\in\cpt(X,Y)\implies T^*\in\cpt(Y^*,X^*)\).\\
            取一列 \(\{f_n\}\subset Y^*\), \(\norm{f_n}\leq C\). 考虑 \(K\). \(K\subseteq Y\) 是 \(X\) 的单位球面在 \(T\) 下的像的闭包. 这里, 取单位球面是因为赋范线性空间中总有标量乘法, 任意集合都可以缩放到单位球之内. 取闭包是因为要利用其紧性. 定义
            \[\Phi_n: K\to\R,\qquad y\mapsto f_n (y) \]
            由于 \(f_n\) 是有界的, 限制在 \(K\) 上也是有界的, 还是一致有界的. 考虑
            \[\abs{\Phi_n(y_1)-\Phi_n(y_2)}=\abs{f_n(y_1)-f_n(y_2)}\leq C\norm{y_1-y_2}\]
            这是因为有界线性泛函本身是 Lipschitz 连续的. 所以 \(\{\Phi_n\}\) 是等度连续的. 根据 Arzelà-Ascoli 定理, 存在一个子列 \(\{\Phi_{n_k}\}\) 使得 \(\{\Phi_{n_k}\}\) 在 \(K\) 上一致收敛. 设 \(\Phi\) 是它的极限. \\
            所以令 \(g_k\) 是一列线性泛函, 其将 \(X\) 的单位球面映入 \(\R\), 定义为 \(x\mapsto \Phi_{n_k} (Tx)=f_{n_k} (Tx)=(T^* f_{n_k}) (x) \). 由于 \(f_{n_k}\) 在 \(K\) 上一致收敛, 所以 \(g_k\) 也收敛. 设 \(g\) 是它的极限, 可以利用单位球面和线性性将其延拓到全空间上, 
            \[\norm{g_k-g}=\norm{T f_{n_k}-T f}\to 0\]
            这就说明了 \(T^* f_{n_k}\) 收敛, 从而 \(T^*\in\cpt(Y^*,X^*)\).\\
            反过来, 取典范嵌入 \(J_Y\), 
            \[J_Y (Tx_n) = T^{**} (J_X (x_n))\]
            这是因为 \(\forall\,f\in Y^*\), \(T^{**} (J_X (x_n))=(J_X (x_n)) (T^* f)=(T^* f) (x_n) = f(Tx_n) \in J_Y(Tx_n)\).
            已经证明了 \(T\) 是紧的可以得到 \(T^*\) 是紧的. 在这里, \(T^*\) 是紧的可以得到 \(T^{**}\) 是紧的, 又因为典范嵌入是等距的单射, 所以有
            \begin{center}
                \(\{x_n\}\) 是有界的 \(\implies\) \(\{J_X(x_n)\}\) 是有界的 \(\implies\) \(\{J_Y (Tx_n)\}\) 是相对紧的 \(\implies\) \(\{Tx_n\}\) 是相对紧的.
            \end{center}
            所以 \(T^*\in \cpt(Y^*,X^*)\implies T\in\cpt(X,Y)\).
        \end{prf}

        \begin{dfn}
            {谱}
            [Spectrum]
            设 \(T\) 是 \(X\) 上的一个线性算子. 定义 \(\lambda\) 在 \(T\) 的\textbf{谱 (Spectrum)}中, 当且仅当线性算子 \(\lambda\Id-T\) 不为双射.\\
            一个线性算子 \(T\) 的谱用 \(\spec T\) 表示.
        \end{dfn}

        \begin{rmk}
            [Dedicatia]
            需要注意的是, 谱和线性代数中的特征值很像, 但谱可以包含 0, 特征值一般不允许为 0.
        \end{rmk}

        \begin{dfn}
            {特征值}
            [Eigenvalue]
            设 \(T\) 是 \(X\) 上的一个线性算子. 定义 \(\lambda\) 在 \(T\) 的\textbf{谱 (Spectrum)}中, 当且仅当线性算子 \(\lambda\Id-T\) 不为单射, 即 \(\ker(\lambda\Id-T)\neq\{0\}\).
        \end{dfn}

        \begin{thm}
            {紧算子的 Frdholm 闭性}
            设 \(X\) 是 \(\C\) 上的 Banach 空间. \(T\in\cpt(X,X)\). 则 \(\forall\,z\in\C, \img(\lambda\Id-T)\) 是 \(X\) 的闭子空间.
        \end{thm}

        \begin{thm}
            {Riesz-Shauder 定理}
            设 \(X,Y\) 是 \(\C\) 上的 Banach 空间, \(T\in\cpt(X,Y)\). 则有:
            \begin{enumerate}
                \item \(\dim X=+\infty\) 时, \(0\in\spec T\);
                \item 除去 \(0\), \(T\) 的谱均为特征值;
                \item 若 \(\lambda\) 是特征值, \(\dim\ker(\lambda\Id-T)<+\infty\);
                \item \(\spec T\) 只可能有一个聚点 \(0\).
            \end{enumerate}
        \end{thm}

    \subsection{Hilbert 空间上的算子}

        \begin{dfn}
            {自伴随算子}
            [Self-adjoint Operator]
            设 $H$ 是 Hilbert 空间, $T: H \to H$ 是一个有界线性算子. 定义 $T$ 是一个\textbf{自伴随算子}当且仅当 \(T=T^*\).
        \end{dfn}

        \begin{dfn}
            {正规算子}
            [Normal Operator]
            设 $H$ 是 Hilbert 空间, $T: H \to H$ 是一个有界线性算子. 定义 $T$ 是一个\textbf{正规算子}当且仅当 \(T^* T = T T^*\).
        \end{dfn}

        \begin{dfn}
            {特征子空间}
            [Eigenspace]
            设 $H$ 是 Hilbert 空间, $T: H \to H$ 是一个有界线性算子. $\lambda$ 是 $T$ 的一个特征值. 定义 $\ker (T-\lambda I)$ 是 $T$ 的一个特征子空间.
        \end{dfn}

        \begin{dfn}
            {不变子空间}
            [Invariant Subspace]
            设 $H$ 是 Hilbert 空间, $T: H \to H$ 是一个有界线性算子. 定义 $M \subseteq H$ 是 $T$ 的一个不变子空间当且仅当 \(T(M) \subseteq M\).
        \end{dfn}

        \begin{thm}
            [orthogonal-complement-invariant]
            {不变子空间的正交补也是不变子空间}
            \(T\) 是 Hilbert 空间 \(H\) 上的自伴随算子, \(M\) 是 \(T\) 的一个不变子空间, 则 \(M^\perp\) 也是 \(T\) 的一个不变子空间.
        \end{thm}

        \begin{prf}
            任取 \(z\in W^\perp, x\in W\):
            \[\expval{Tz, x}=\expval{z, T^* x}=\expval{z, T x}=0.\]
        \end{prf}
        
        \begin{thm}
            {自伴随算子作用下的内积是实数}
            设 \(T\) 是 Hilbert 空间 \(H\) 上的自伴随算子, 则 \(\forall\,x\in H\), \(\expval{Tx, x}\in\R\).
        \end{thm}

        \begin{prf}
            By:
            \[\overline{\expval{Tx, x}}=\expval{x, T^* x}=\expval{x, T x}=\expval{Tx, x}.\]
            因此, \(\expval{Tx, x}\) 是实数.
        \end{prf}

        \begin{thm}
            {Hilbert 空间上的范数的等价定义}
            设 \(T\) 是 Hilbert 空间上的自伴随算子, 则
            \[\norm{T}=\sup_{\substack{\norm{x}=\norm{y}\leq 1}}\expval{Tx, y}.\]
        \end{thm}

        \begin{prf}
            (1) 
            \[\expval{Tx, y}\leq\norm{T}\norm{x}\norm{y}\leq\norm{T}.  \]
            (2) 若 \(T=0\), trivial;\\
            若 \(T\neq 0\), \(x\in H\setminus\ker T\), 则
            \[\norm{Tx}=\abs{\expval{Tx, \frac{Tx}{\norm{Tx}}}}\]
            取上确界
            \[\norm{T}=\sup_{\substack{\norm{x}\leq 1 \\ x\in H\setminus\ker T}}\norm{Tx}=\sup\abs{\expval{Tx, \frac{Tx}{\norm{Tx}}}} \]
            现在令 \(y=\frac{Tx}{\norm{Tx}}\), 则
            \[\leq\sup_{\substack{\norm{x}=\norm{y}\leq 1}}\abs{\expval{Tx, y}}.\]
            现在去掉绝对值: 因为该内积总是实数, 所以 \(\expval{Tx, y}\) 要么是正的, 要么是负的. 如果是负的, 上述证明中必然存在 \(x'=-x\) 使得 \(\expval{Tx', y}=-\expval{Tx, y}\) 是正的. 从而上确界是正的. 这就说明了
            \[\norm{T}\leq\sup_{\substack{\norm{x}=\norm{y}\leq 1}}\expval{Tx, y}.\]
        \end{prf}

        \begin{thm}
            [definition-norm-self-adjoint-2]
            {自伴随算子范数的等价定义2}
            设 \(T\) 是 Hilbert 空间上的自伴随算子, 则
            \[\norm{T}=\sup_{\norm{x}\leq 1}\abs{\expval{Tx, x}}.\]
        \end{thm}

        \begin{prf}
            设右边的值是 \(a\), 则根据自伴随算子作用, \(\expval{Tx, x}\) 是实数. 那么:
            \[\expval{T(x+y), x+y}=\expval{Tx, x}+\expval{Ty, y}+2\expval{Tx, y}.\]
            \[\expval{T(x-y), x-y}=\expval{Tx, x}+\expval{Ty, y}-2\expval{Tx, y}.\]
            两式相加:
            \[\expval{Tx, y} = \frac{1}{4}[\expval{T(x+y), x+y}-\expval{T(x-y), x-y}]\leq\frac{1}{4}\qty(\norm{x+y}^2+\norm{x-y}^2)=\frac{1}{2}\qty(\norm{x}^2+\norm{y}^2) \]
            令 \(\tilde{x}=\sqrt{a}x\), \(\tilde{y}=\frac{y}{\sqrt{a}}\):
            \[\expval{Tx,y}=\expval{T\tilde{x}, \tilde{y}} \leq \frac{1}{2}\qty(\norm{\tilde{x}}^2+\norm{\tilde{y}}^2)=\frac{1}{2}\qty(a\norm{x}^2+\frac{1}{a}\norm{y}^2).\]
            令 \(a=\frac{\norm{y}}{\norm{x}}\), 则
            \[\expval{Tx, y} \leq \norm{x}\norm{y}.\]
            取上确界, 得证.
        \end{prf}

        \begin{thm}
            {不同特征值的特征子空间是正交的}
            设 \(T\) 是 Hilbert 空间上的自伴随算子, \(\alpha\) 和 \(\beta\) 是 \(T\) 的两个不同的特征值, 则 \(E_\alpha:=\ker (T-\alpha I)\perp E_\beta:=\ker (T-\beta I)\).
        \end{thm}

        \begin{prf}
            任取 \(x\in E_\alpha\), \(y\in E_\beta\).
            \[\expval{Tx, y}=\expval{\alpha x, y}=\alpha\expval{x,y}\]
            \[\expval{x, Ty}=\expval{x, \beta y}=\beta\expval{x,y}\]
            所以
            \[(\alpha-\beta)\expval{x, y}=0\]
            由于 \(\alpha\neq\beta\), 所以 \(\expval{x, y}=0\). 从而 \(E_\alpha\perp E_\beta\).
        \end{prf}

        \begin{thm}
            {谱定理}
            [Spectral Theorem]
            设 $H$ 是 Hilbert 空间, $T$ 既是自伴随算子又是紧算子, 则:
            \begin{enumerate}[label=(\roman*)]
                \item 存在 \(H\) 的标准正交基 \(\{e_i\}\) 使得 \(e_i\) 可以对角化 \(T\), 即 \(T e_i = \lambda_i e_i\) 其中 \(\lambda_i\) 是 \(T\) 的特征值;
                \item 特征子空间 \(\ker (T-\lambda I)\) 是有限维的;
                \item 所有特征值都是实数, 最多只可能有一个极限点 0: $\forall r>0$, 存在有限多个特征值 $\lambda$ 使得 $|\lambda| > r$.
            \end{enumerate}
        \end{thm}

        \begin{prf}
            Step1: 先找到一个特征向量: 根据定理\ref{thm:definition-norm-self-adjoint-2}, 
            \[\norm{T}=\sup_{\norm{x}\leq 1}\expval{Tx, x}\]
            可以找到一列 \(\{x_n\}\), \(\norm{x_n}=1\), 使得
            \[\expval{Tx_n, x_n}\to\lambda\in\R.\]
            其中 \(|\lambda|=\norm{T}\neq 0\). 因为 \(T\) 是紧算子, \(T x_n\) 有收敛子列 \(Tx_{n_k}\to y\), 所以 \(y\neq 0\):
            \[\norm{Tx_{n_k}-\lambda x_{n_k}}^2 = \norm{Tx_{n_k}}^2+\lambda^2\norm{x_{n_k}}^2-2\lambda\expval{Tx_{n_k}, x_{n_k}}\leq 2\norm{T}^2-2\lambda\expval{Tx_{n_k}, x_{n_k}}\to 0\]
            这说明随着 \(Tx_{n_k}\to y\), \(\lambda x_{n_k}\to y\):
            \[T y = \lim_{k\to\infty} T x_{n_k} = \lim_{k\to\infty} \lambda x_{n_k} = \lambda y\]
            所以 \(y\) 是 \(T\) 的一个特征向量, \(\lambda\) 是 \(T\) 的一个特征值.\\
            Step2: 根据 Zorn 引理, 可以找到 \(H\) 的某个最大的标准正交基 \(\mathcal{B}\) 使得 \(\mathcal{B}\) 中的每个元素都是 \(T\) 的一个特征向量. 设 \(M=\spn(\mathcal{B})\), 则 \(M\) 是 \(T\) 的一个不变子空间. \\
            根据定理\ref{thm:orthogonal-complement-invariant}, \(M^\perp\) 也是 \(T\) 的一个不变子空间. 如果 \(M^\perp\neq \{0\}\), 则 \(T\) 限制在 \(M^\perp\) 上也是自伴随算子, 根据 Step1, \(M^\perp\) 中存在一个特征向量, 这与 \(\mathcal{B}\) 的最大性矛盾. 从而 \(M^\perp=\{0\}\), 即 \(\mathcal{B}\) 是 \(H\) 的一个标准正交基.\\
            所以 \(M^\perp\) 是平凡的, \(M=H\).\\
            Step3: 我们要证明唯一的可能极限点是 0. 取 \(H\) 的标准正交基 \(\{e_i\}_{i\in I}\). 假如存在某个 \(r>0\) 使得 \(\{|\lambda_i|\geq r\}\) 有无穷个, 则
            \[\norm{Te_i-Te_j}=\norm{\lambda_ie_i - \lambda_je_j}\]
            根据勾股定理:
            \[\norm{\lambda_ie_i - \lambda_je_j}^2=\lambda_i^2+\lambda_j^2\geq 2r^2 >0\]
            这说明标准正交基在 \(T\) 的像不是列紧的, 与 \(T\) 是紧算子矛盾. 从而 \(T\) 的特征值最多只有一个极限点 0. \\
            这也体现了 \(\dim E_\lambda<\infty\) 对所有 \(\lambda\neq 0\) 成立.
        \end{prf}

        \begin{ins}
            {\(\Lp[2]([0,1])\) 上的 Volterra 积分算子}
            在 \(\Lp[2]([0,1])\) 空间上令积分核为 \(K(x,t)=\I_{t\leq x}\), 则定义积分算子
            \[V:=f\mapsto\int_0^1K(x,t)f(t)\dd{t}=\int_0^xf(t)\dd{t} \]
            \(V=\) 就是 \([0,1]\) 上的 Volterra 积分算子. \\
            \(V\) 的唯一特征值是 \(0\).
        \end{ins}
        
        \begin{thm}
            {Hilbert 空间的直和分解}
            设 \(H\) 是 Hilbert 空间, \(T\) 是 \(H\) 上的一个自伴随算子, 则
            \[H=\bigoplus_{\lambda\in\sigma(T)}\ker (T-\lambda I).\]
            并且这是一个正交直和.
        \end{thm}

        \begin{thm}
            {自伴随算子可交换等价于可同时对角化}
            设 \(H\) 是 Hilbert 空间, \(\{T_i\}_{i\in I}\) 是 \(H\) 上的一族自伴随算子, 且满足交换性:
            \[T_i T_j = T_j T_i, \forall i,j\in I.\]
            则存在 \(H\) 的一个标准正交基 \(\{e_k\}_{k\in K}\) 使得 \(\{e_k\}_{k\in K}\) 同时对角化 \(\{T_i\}_{i\in I}\), 即 \(\forall i\in I, k\in K, T_i e_k = \lambda_{i,k} e_k\).
        \end{thm}

        \begin{thm}
            {正规算子可对角化}
            对于正规算子 \(S\), \(S\) 可被分解为
            \[S = A +\ii B\]
            其中 \(A\) 和 \(B\) 都是自伴随算子, 显式表示为
            \[A = \frac{1}{2}(S + S^*), \quad B = \frac{1}{2\ii}(S - S^*)\]
            且 \(A\) 和 \(B\) 交换. 从而 \(S\) 可被一个标准正交基对角化.
        \end{thm}

        \begin{vardfn}
            {酉算子 / 幺正算子}
            [Unitary Operator]
            设 \(H\) 是 Hilbert 空间, \(T: H \to H\) 是一个有界线性算子. 定义 \(T\) 是一个\textbf{酉算子}当且仅当 \(T\) 是 \(H\) 的等距同构.
        \end{vardfn}

        \begin{varthm}*
            {酉算子的性质}
            \begin{enumerate}[label=(\roman*)]
                \item 酉算子一定是自伴随的;
                \item \(U\) 是酉算子, 则 \(\expval{Ux, Uy} = \expval{x, y}\) 对所有 \(x, y \in H\) 成立;
                \item \(U\) 是酉算子, 则 \(U^{-1}\) 一定存在且也是一个酉算子; 
                \item \(U\) 是酉算子等价于 \(U^* U = I\) 等价于 \(U U^* = I\), 等价于 \(U^* = U^{-1}\).
            \end{enumerate}
        \end{varthm}

        \begin{varthm}*
            {酉算子是紧算子当且仅当空间是有限维的}
        \end{varthm}

        \begin{xmp}
            {Hilbert-Schmidt 算子}
            设 \(H\) 是 Hilbert 空间, \(\{e_i\}_{i\in I}\) 是 \(H\) 的标准正交基. 定义 \(T: H \to H\) 是一个\textbf{Hilbert-Schmidt 算子}当且仅当
            \[\sum_{i\in I} \norm{Te_i}^2 < +\infty.\]
            其 Hilbert-Schmidt 范数定义为
            \[\norm{T}_2 = \sqrt{\sum_{i\in I} \norm{Te_i}^2}.\]
        \end{xmp}

        \begin{ppt}
            {Hilbert-Schmidt 范数的定义不依赖于标准正交基的选取}
        \end{ppt}

        \begin{rmk}
            [Dedicatia]
            在有限维的时候, 这个范数还有一个熟悉的名字叫做矩阵的 Frobenius 范数.
        \end{rmk}

        \begin{ppt}
            {L2 空间上的 Hilbert-Schmidt 算子}
            在 L2 空间上, Hilbert-Schmidt 算子就是积分算子. \\
            这时的 Hilbert-Schmidt 范数就是积分核的 L2 范数:
            \[\norm{T}_2 = \sqrt{\int_D\int_D |K(x,y)|^2 \dd{x}\dd{y}}.\]
        \end{ppt}

        \begin{thm}
            {Hilbert-Schmidt 算子是紧算子}
            设 \(H\) 是 Hilbert 空间, 则 \(H\) 上的 Hilbert-Schmidt 算子都是紧算子.
        \end{thm}

        \begin{prf}
            取 \(\norm{x}=1\), \(x\in H\). 将 \(x\) 投影到标准正交基上:
            \[x=\sum_{i\in I}c_ie_i\]
            令
            \[T_nx=\sum_{i\leq n}c_iTe_i\]
            那么 \(T_n\) 是有限秩算子. 考虑余项, 利用 Cauchy-Schwarz 不等式:
            \[\norm{Tx-T_nx}^2 = \norm{\sum_{i>n}c_iTe_i}^2\leq\qty(\sum_{i\in I}|c_i|^2)\qty(\sum_{i>n}\norm{Te_i}^2)\]
            由于
            \[\sum_{i\in I}|c_i|^2=\norm{x}^2=1, \qquad \sum_{i>n}\norm{Te_i}^2\to 0\]
            所以 \(T_n\) 收敛到 \(T\). 从而 \(T\) 可以被一列有限秩算子近似, 是一个紧算子.
        \end{prf}

        \begin{thm}
            {Hilbert-Schmidt 范数总不小于算子范数}
            设 \(H\) 是 Hilbert 空间, \(T: H \to H\) 是一个 Hilbert-Schmidt 算子, 则 \(\norm{T}\leq\norm{T}_2\).
        \end{thm}

        \begin{prf}
            取 \(\norm{x}=1\), \(x\in H\). 将 \(x\) 投影到标准正交基上:
            \[x=\sum_{i\in I}c_ie_i\]
            利用 Cauchy-Schwarz 不等式:
            \[\norm{Tx}^2 = \norm{\sum_{i\in I}c_iTe_i}^2\leq\qty(\sum_{i\in I}|c_i|^2)\qty(\sum_{i\in I}\norm{Te_i}^2)=\norm{T}_2^2.\]
            取上确界, 得证.
        \end{prf}
        
        \begin{thm}
            {Hilbert-Schmidt 算子是 Hilbert 空间}
            设 \(H\) 是 Hilbert 空间, 则 \(H\) 上的 Hilbert-Schmidt 算子也是一个 Hilbert 空间, 其内积定义为
            \[\expval{T, S} = \sum_{i\in I} \expval{Te_i, Se_i}.\]
            换句话说, Hilbert-Schmidt 范数可以诱导内积, 且是完备的.
        \end{thm}

    \subsection{Fredholm 算子}

        \begin{dfn}
            [cokernel]
            {余核}
            [Cokernel]
            设 \(X\) 和 \(Y\) 是 Banach 空间, \(T:X\to Y\) 是一个有界线性算子. 定义 \(T\) 的\textbf{余核}为
            \[\coker T := Y/\img T.\]
        \end{dfn}

        \begin{dfn}
            {余维数}
            [Codimension]
            设 \(X\) 和 \(Y\) 是 Banach 空间, \(T:X\to Y\) 是一个有界线性算子. 定义 \(T\) 的像空间的\textbf{余维数}为
            \[\codim\img T := \dim\coker T.\]
            所以这两个记号本质上是一回事. 
        \end{dfn}

        \begin{rmk}
            [Dedicatia]
            如果放到有限维的情形, 这个关系可以表示为
            \[\dim X = \dim\ker T + \dim\img T, \qquad \dim Y = \dim\img T + \codim\img T.\]
            这正是\textbf{线性代数基本定理}的内容: 秩 = 零化度 + 维数. 但是在无限维的情形, 这个关系不再成立了, 因为 \(\img T\) 不一定是闭的, 这时的 \(Y/\img T\) 甚至不一定是 Hausdorff 空间, 也就没有维数的概念. 所以我们需要要求 \(\img T\) 是闭的, 从而 \(\coker T\) 是一个 Banach 空间, 才能定义维数.
        \end{rmk}

        \begin{dfn}
            [fredholm-operator]
            {Fredholm 算子}
            [Fredholm Operator]
            设 \(X\) 和 \(Y\) 是 Banach 空间, \(T:X\to Y\) 是一个有界线性算子. 定义 \(T\) 是一个\textbf{Fredholm 算子}当且仅当: \(\dim\ker T<\infty\), \(\dim\coker T<\infty\), 且 \(\img T\) 是闭的.\\
            全体 Fredholm 算子记作 \(\fred(X,Y)\).
        \end{dfn}

        \begin{dfn}
            [fredholm-index]
            {Fredholm 算子的指数}
            [Index of a Fredholm Operator]
            设 \(X\) 和 \(Y\) 是 Banach 空间, \(T:X\to Y\) 是一个 Fredholm 算子. 定义 \(T\) 的\textbf{指数}为
            \[\ind T := \dim\ker T - \dim\coker T.\]
        \end{dfn}

        \begin{rmk}
            [Dedicatia]
            一个 Fredholm 算子与一个可逆算子的行为很相似. 除了有限多个独立的 \(Tx=0\) 的解之外, 以及有限多个 \(y\in Y\) 无法表示为 \(Tx\) 的形式之外, 其他的行为都和可逆算子一样. 从这个角度来看, Fredholm 算子是可逆算子在 Banach 空间上的一个推广.
        \end{rmk}

        \begin{ins}
            {序列空间上的左移 / 右移算子是 Fredholm 算子}
            设 \(X=\ell^2\), 其上的左移算子
            \[T: x=(x_1, x_2, \ldots) \mapsto (x_2, x_3, \ldots)\]
            满足
            \[\ker T=\{(x_1, 0, \cdots), x_1\in\F \}, \dim\ker T = 1; \]
            \[\img T = \ell^2, \dim\coker T = 0.\]
            从而 \(T\) 是一个 Fredholm 算子, 且 \(\ind T = 1\). 同样的, 其上的右移算子也是一个 Fredholm 算子, 且 \(\ind T = -1\).
        \end{ins}

        \begin{thm}
            {Fredholm 算子在有界线性算子中是开集}
            [Fredholm Operators Form an Open Subset of Bounded Linear Operators]
            设 \(X\) 和 \(Y\) 是 Banach 空间, 则 \(\fred(X,Y)\) 是 \(\blo(X,Y)\) 中的一个开集. 
        \end{thm}

        \begin{prf}
            我们设 \(T\in\fred(X,Y)\). 目标是证明存在 \(\epsilon_0>0\) s.t. \(\forall\,S\in\blo(X,Y), \norm{S-T}<\epsilon_0\implies\) \(S\in\fred(X,Y)\). \\
            显然, 存在闭子空间 \(X',Y'\) 使得 \(X=X'\oplus\ker T\), \(Y=Y'\oplus\img T\). 其中 \(\dim\ker T<\infty\), \(\dim Y'<\infty\). 考虑 
            \[\tilde{T}: X'\oplus Y'\to Y, (x_1, y_1)\mapsto Tx_1+y_1\] 
            这是一个双射. 根据\hyperref[thm:omt]{开映射定理}, \(\tilde{T}\) 是作为 Banach 空间的同构. \\
            令 \(\epsilon_0=\frac{1}{\norm{T^{-1}}}\), 得到
            \[\norm{\tilde{S}-\tilde{T}}<\epsilon_0\]
            由此确定了的 \(S\) 满足
            \[Sx_1+y_1 = \tilde{S}(x_1, y_1)\implies \norm{T-S} < \epsilon_0\]
            有: \\
            (1) \(\ker S\) 不在 \(X'\) 中, 否则 \(S\) 在 \(X'\) 上不是单射. 说明 \(\dim\ker S\leq\dim\ker T < \infty\);
            (2) \(\img S\oplus Y'=Y\), 这是因为
            \[\forall\, y\in Y, y = Sx_1+y_1=\tilde{S} (x_1, y_1)\]
            说明 \(\dim\coker S\leq\dim Y' < \infty\);
            (3) \(\img S\) 是闭的, 因为 \(\tilde{S}\) 是同构.\\
            从而 \(S\) 是一个 Fredholm 算子.
        \end{prf}

        \begin{thm}
            [index-fredholm-continuous]
            {Fredholm 指数是连续映射}
            \(\ind: \fred(X,Y)\to \Z\) 是连续映射. 即\\ 
            设 \(X,Y\) 是 Banach 空间. \(T\in\fred(X,Y)\). \(\epsilon>0\) 是小量, 则 \(\forall\,S\in\blo(X,Y)\land \norm{S}<\epsilon\), \(T+S\in\fred(X,Y)\), \(\ind(T+S)=\ind T\).
        \end{thm}

        \begin{thm}
            {紧算子与 Fredholm 算子的关系}
            \(X\) 是 Banach 空间. \(K\in\cpt(X,X)\). \(T=\Id+K\). 则 \(T\in\fred(X,X)\), \(\ind T=0\).
        \end{thm}

        \begin{prf}
            逐一验证 Frodholm 算子的条件:\\
            (1) \(\dim\ker\Id+K<\infty\): 设 \(\{x_n\}\subset\ker(\Id+K)\). \(\norm{x_n}\leq 1,\forall n\). 那么 \(Kx_n\) 有收敛子列, 但是 \((\Id+K)x_n=0\). 
            所以 \(x_n=\Id x_n\) 也存在收敛子列. 由于 \(\{x_n\}\) 是任取的, 所以 \(\ker\Id+K\) 的单位闭球是紧的. 所以是有限维的.\\
            (2) \(\img\Id+K\) 是闭的: 令 \(X=\ker T\oplus V\). 由于 \(V\cong\img T\), 所以 \(T\) 限制在 \(V\) 上是双射. 我们只需要考虑 \(T\) 限制在 \(V\) 上是不是有界的. \\
            反设如果不是双射, 则 \(\exists\{x_n\}\subset V\), 其中 \(\norm{x_n}\geq\epsilon\) 但是 \(Tx_n\to 0\). 
            利用线性性令 \(y_n=x_n/\norm{x_n}\), 则 \(\norm{y_n}=1\). 同样根据 \(K\) 是紧算子, \(Kx_n\) 有收敛到 \(z\) 的子列, 所以 \(y_n\to -z\), 其中 \(\norm{z}=1\) 显然 \(z\in V\). 但是按照假设 \((\Id+K)z=0\). 矛盾. 所以 \(\img T\) 是闭的.\\
            (3) \(\dim\coker T<+\infty\): 考虑其对偶算子 \(T^*=\Id+K^*\). 根据定理\ref{thm:schauder} \(K^*\) 是紧算子, 同样有 \(\dim\ker T^*<+\infty\). 而 \(\img T\) 是闭的, 有 
            \[\qty(\frac{X}{\img T})^*\cong\ker T^* \]
            所以 \(\dim\coker T=\dim\ker T^*<+\infty\).\\
            (4) \(\ind T=0\): 令 \(T_t:=\Id +tK\). \(\ind\Id=0\), 根据连续性, \(\ind T=0\).
        \end{prf}

        \begin{dfn}
            {算子的模 Frdholm 逆}
            [Parametrix of an Operator]
            设 \(X\) 和 \(Y\) 是 Banach 空间, \(T:X\to Y\) 是一个有界线性算子. 定义 \(S:Y\to X\) 是 \(T\) 的一个\textbf{模 Frdholm 逆}当且仅当 \(ST-\Id_X\) 和 \(TS-\Id_Y\) 都是紧算子.
        \end{dfn}

        \begin{thm}
            {Atkinson 定理}
            设 \(X\) 和 \(Y\) 是 Banach 空间, \(T:X\to Y\) 是一个有界线性算子. 则 \(T\) 是一个 Fredholm 算子当且仅当 \(T\) 存在模 Frdholm 逆.
        \end{thm}

        \begin{rmk}
            [Dedicatia]
            如果令
            \[\mathscr{C}(X,X)=\blo(X,X) / \cpt(X,X)\]
            那么 \(T\) 是 Fredholm 算子当且仅当 \(T\) 在 \(\mathscr{C}(X,X)\) 中是可逆的. 这给了我们一种更高级别的视角来观看 Fredholm 算子.
        \end{rmk}

    \subsection{【测度论】Riesz 表示}

        \begin{dfn}
            {拓扑群}
            [Topological Group]
            设 \(G\) 是一个群, \(G\) 上的一个拓扑使得 \(G\) 的乘法和取逆运算都是连续的, 则称 \(G\) 是一个\textbf{拓扑群}.
        \end{dfn}

        \begin{rmk}
            [Dedicatia]
            这个定义非常抽象. 我们将其详细解释: \\
            群上具有乘法运算 
            \[\mu:\qquad\begin{matrix}
                (G, G) & \to & G\\
                (x,y) & \mapsto & x\cdot y
            \end{matrix} \]
            连续指的是: 
            \paragraph{使用逆开性定义: 将开集拉回开集.} \(U\subseteq G\) 是一个开集, 那么它的原像必须是 \(G\times G\) 中的一个开集:
            \[U\in\tau_G \implies \mu^{-1}(U)\in\tau'_{G\times G}\]
            \paragraph{使用开集定义.} 对任意 \(x,y\in G\), 若 \(x\cdot y\in W\) 落在开邻域中, 那么存在 \(x\) 的开邻域 \(U\) 和 \(y\) 的开邻域 \(V\) 使得 \(U\cdot V \subseteq W\). 
            \paragraph{使用点定义.} 对任意 \(x,y\in G\), 若 \(x\cdot y\in W\) 落在开邻域中, 那么存在 \(x\) 的开邻域 \(U\) 和 \(y\) 的开邻域 \(V\) 使得 \(\forall u\in U, v\in V, u\cdot v \in W\). 
            类似地可以定义取逆运算的连续性.
        \end{rmk}

        \begin{ins}
            {\(\R\) 是一个拓扑群}
            \((\R,+)\) 是一个拓扑群. 因为 \(\R\) 上的加法 (\(+\)) 和取负运算都是连续的.
        \end{ins}

        \bysorry